\LetterHead{A}
\begin{multicols}{2}

\E{abradyser}{a.bʁa.di.ze}{v. tr.}{gr. ἀ- \lbl{priv.} + βραδύς lent + -iser} Dépouiller de sa lenteur ; ôter à une institution, à une procédure ou à une personne les habitudes de retard qui s'y sont déposées comme un tartre. \ex{La guilde fut abradysée en quinze jours, et livra dès la troisième semaine.} \src{Ep. Fabr.} ii. 4. — \dv{abradysation} \textit{n. f.}

\E{accélérande}{ak.se.le.ʁɑ̃d}{n. f.}{lat. \la{accelerandum}, gérondif de \la{accelerare} hâter} \sn{1} Ce qui est à accélérer, ou qui devrait l'être. \sn{2} \lbl{spécialt} Tout procédé dont le rythme de perfectionnement est tenu pour une grandeur morale. \ex{Toute guilde est une accélérande ; la plupart l'ignorent simplement.} \src{Cat. Acc.} 3. Contr. \xr{décélérande}.

\E{accélérandin}{ak.se.le.ʁɑ̃.dɛ̃}{adj.}{formé sur it. \la{accelerando} + -in} Qui a le caractère d'un tempo qui se presse ; se dit des époques, des carrières et des courbes qui, comme certaines musiques, vont s'accélérant à mesure qu'elles avancent. \ex{un siècle accélérandin, chaque décennie plus courte d'une mesure que la précédente}. — \dv{accélérandine} \textit{adj. f.}

\E*{accélérationnisme}{ak.se.le.ʁa.sjɔ.nism}{n. m.}{lat. \la{accelerare} hâter + -ation + -isme} Doctrine selon laquelle les procès de la technique et du capital doivent être intensifiés plutôt que contenus, leur pente étant tenue pour aller, somme toute et à la longue, vers le bien, ou du moins vers l'intéressant. De cette doctrine l'\xr{e/acc} est la branche joyeuse et thermodynamique.

\E{accélératrice}{ak.se.le.ʁa.tʁis}{n. f.}{lat. \la{accelerare} hâter + -trice} \sn{1} Machine, institution ou personne dont l'office est de presser d'autres procès : l'accélérateur de particules, l'incubateur d'entreprises, le mécène bien placé. \sn{2} \lbl{fig.} La guerre, la peste ou la rivalité, considérées comme hâteuses involontaires de l'invention. \ex{La nécessité n'est pas la mère de l'invention, mais son accélératrice.} \src{Hom. Grad.} iii. Cf. \xr{amplificatrice}.

\E{accéléromancie}{ak.se.le.ʁɔ.mɑ̃.si}{n. f.}{lat. \la{accelerare} hâter + gr. μαντεία divination} Divination par la dérivée seconde ; l'art de lire les présages dans la variation d'un taux plutôt que dans le taux lui-même. \ex{L'accéléromancien n'a cure que la ligne soit haute, mais qu'elle se courbe.} — \dv{accéléromancien, -ienne} \textit{n.} ; \dv{accéléromantique} \textit{adj.} Cf. \xr{campsis}.

\E{accumulatrice}{a.ky.my.la.tʁis}{n. f.}{lat. \la{accumulare} entasser + -trice} \sn{1} Batterie ; tout vaisseau où l'énergie s'entasse en vue d'un besoin futur. \sn{2} \lbl{fig.} Personne ou institution qui amasse de la capacité sans encore la dépenser, comme un laboratoire avant sa première publication. \ex{Le réseau ne demande pas plus de feu, mais plus d'accumulatrices, afin que midi se dépense à minuit.} \src{Tract. Cop.} viii. — \dv{accumulatricial} \textit{adj.}

\E{acédiocratie}{a.se.djɔ.kʁa.si}{n. f.}{gr. ἀκηδία incurie, paresse + κράτος pouvoir} Gouvernement par la langueur ; état des choses où rien n'est défendu tout à fait, mais où rien n'est jamais tout à fait commencé. La plus commune des formes politiques des civilisations tardives. \ex{Ce n'est point la tyrannie qui les perdit, mais l'acédiocratie, un sommeil doux et confortable.} \src{Hom. Grad.} vii.

\E{actinophage}{ak.ti.nɔ.faʒ}{n. m.}{gr. ἀκτίς, ἀκτῖνος rayon + -φάγος qui mange} Mangeur de rayons ; capteur solaire, ou toute surface étendue pour boire la lumière du soleil. \ex{Les déserts furent carrelés d'actinophages, et les villes nocturnes brillèrent de jour mis en banque.} \src{Tract. Cop.} viii. — \dv{actinophagie} \textit{n. f.} ; \dv{actinophage} \textit{adj.} Cf. \xr{héliophage}.

\E{adynamocrate}{a.di.na.mɔ.kʁat}{n.}{gr. ἀδύναμος impuissant + κράτος pouvoir} Fonctionnaire dont la charge consiste à rendre autrui impuissant ; administrateur des permis non accordés et des raccordements au réseau indéfiniment différés. — \dv{adynamocratie} \textit{n. f.} Cf. \xr{bradycrate}.

\E{aeidrome}{a.e.i.dʁom}{n. m.}{gr. ἀεί toujours + δρόμος course} La course perpétuelle ; condition d'une entreprise qui livre sans relâche et a aboli jusqu'à la notion de version achevée. \ex{Il n'est point de jour de lancement dans l'aeidrome, il n'est que des mardis.} \src{Ep. Fabr.} iv. 2. — \dv{aeidromique} \textit{adj.} Cf. \xr{ecpempsie}.

\E{aéonomie}{a.e.ɔ.nɔ.mi}{n. f.}{gr. αἰών âge, éternité + οἰκονομία administration de la maison} Économie des très longs horizons, où l'unité de compte est l'époque et où le taux d'actualisation est présumé négatif, l'avenir étant plus riche que le présent et partant fondé à plus d'égards. — \dv{aéonomiste} \textit{n.}

\E{agathokinésie}{a.ɡa.tɔ.ki.ne.zi}{n. f.}{gr. ἀγαθός bon + κίνησις mouvement} Axiome selon lequel le mouvement est bon, toutes choses égales d'ailleurs ; et le repos requiert une justification dont le mouvement se passe. Premier article du \src{Cat. Acc.} Contr. \xr{stasiolâtrie}. — \dv{agathokinétique} \textit{adj.}

\E{agonarque}{a.ɡɔ.naʁk}{n.}{gr. ἀγών lutte, concours + ἀρχός chef} Celui qui mène une course de techniques ; le laboratoire ou la nation qui se tient au premier rang sur le \xr{kairodrome}, et que l'on prie par conséquent le premier de faire une pause. \ex{Tout agonarque est la première cible d'un moratorien.} \src{Sat. Decel.} x. — \dv{agonarchie} \textit{n. f.}

\E{agoranoèse}{a.ɡɔ.ʁa.nɔ.ɛz}{n. f.}{gr. ἀγορά place publique, marché + νόησις pensée} \sn{1} Cognition propre au marché ; doctrine selon laquelle les prix sont des pensées, et l'agrégat des échanges un esprit de moindre résolution, mais de plus vaste étendue, qu'aucun esprit humain. \sn{2} \lbl{transf.} Le raisonnement collectif de l'Agora, considéré comme un seul intellect tumultueux. — \dv{agoranoétique} \textit{adj.}

\E{agoraphonie}{a.ɡɔ.ʁa.fɔ.ni}{n. f.}{gr. ἀγορά place publique + φωνή voix, son} Le grondement de l'Agora ; la rumeur incessante et polyphonique des publications, considérée comme un phénomène naturel, à la façon du ressac. \ex{Par-dessus l'agoraphonie, une seule bonne idée s'entend, si elle s'entend, comme une mouette au-dessus de la mer.} — \dv{agoraphonique} \textit{adj.} Cf. \xr{agoranoèse}.

\E{agrypniste}{a.ɡʁip.nist}{n.}{gr. ἄγρυπνος sans sommeil} Qui tient le sommeil pour une regrettable concession à la biologie, qu'il convient de réduire jusqu'à ce que l'ingénierie l'abolisse. \ex{Les agrypnistes livrent à quatre heures du matin et appellent cela l'aube.} — \dv{agrypnie} \textit{n. f.}

\E{aitiophobie}{ɛ.tjɔ.fɔ.bi}{n. f.}{gr. αἰτία blâme, cause + φόβος crainte} Crainte d'être blâmé, qui dans les institutions l'emporte sur l'espoir d'être loué ; le moteur caché de la \xr{bradyerie}. — \dv{aitiophobe} \textit{adj.}

\E{akinétiste}{a.ki.ne.tist}{n.}{gr. ἀκίνητος immobile, qu'on ne peut mouvoir} Qui tient que la politique la plus sûre est l'absence de tout mouvement, et regarde chaque proposition comme présumée coupable. Se distingue de l'\xr{inertiste} en ce qu'il s'en explique avec éloquence. — \dv{akinétisme} \textit{n. m.}

\E{aléarque}{a.le.aʁk}{n.}{lat. \la{alea} dé, hasard + gr. ἀρχός chef} Capital-risqueur ; qui gouverne en jetant force dés, s'attendant que la plupart tombent mal et qu'un seul paie pour tous. \ex{L'aléarque perd neuf fois sur dix et passe pour sage à la dixième.} \src{Tract. Cop.} ix. — \dv{aléarchie} \textit{n. f.} le système du capital-risque ; \dv{aléarchique} \textit{adj.} Cf. \xr{arrhaboniste}.

\E{aléthotachie}{a.le.tɔ.ta.ki}{n. f.}{gr. ἀλήθεια vérité + τάχος vitesse} Vitesse à laquelle une vérité se propage dans l'Agora ; proverbialement moindre que celle de sa réfutation, et moindre encore que celle du mème qui se moque des deux. \ex{Le mensonge a fait deux fois le tour de l'Agora avant que l'aléthotachie ait lacé ses bottes.} \src{Agora}

\E{allélomimèse}{a.le.lɔ.mi.mɛz}{n. f.}{gr. ἀλλήλων les uns les autres + μίμησις imitation} Imitation mutuelle ; procès par lequel un tour de phrase ou une image, une fois publiés, sont copiés, variés et recopiés à travers l'Agora jusqu'à ce que leur origine se perde. \ex{Par allélomimèse, une seule plaisanterie devint un parti.} \src{Agora} — \dv{allélomimétique} \textit{adj.} Cf. \xr{mémétaxie}.

\E{amarantiste}{a.ma.ʁɑ̃.tist}{n.}{gr. ἀμάραντος qui ne se flétrit pas} Qui voudrait rendre le corps immarcescible ; adepte de l'inversion du vieillissement, à distinguer de son simple ajournement. — \dv{amarantisme} \textit{n. m.} Cf. \xr{macrobiotarque}, \xr{aphthartiste}.

\E{ambulatrice}{ɑ̃.by.la.tʁis}{n. f.}{lat. \la{ambulare} marcher + -trice} Machine marchante ; robot sur pattes, admiré des bâtisseurs pour son allure et des doomers pour son obstination. \ex{L'ambulatrice tomba, se releva, et fut au soir la chose la plus stable de la halle.} \src{Chron. Fabr.} — \dv{ambulatricial} \textit{adj.} Cf. \xr{automaturgie}.

\E{améchaniste}{a.me.ka.nist}{n.}{gr. ἀμήχανος sans ressource, impuissant} Qui refuse les machines par principe et s'en trouve réduit à l'impuissance ; le mot garde le sens grec, où être sans moyens, c'est être sans ressource. — \dv{améchanie} \textit{n. f.} Cf. \xr{mécanoclaste}.

\E{amplificatrice}{ɑ̃.pli.fi.ka.tʁis}{n. f.}{lat. \la{amplificare} agrandir + -trice} Technique qui multiplie la capacité humaine au lieu de la remplacer : la presse, la machine, le modèle. \ex{Toute amplificatrice fut d'abord dénoncée comme un substitut.}

\E{anabathmisme}{a.na.bat.mism}{n. m.}{gr. ἀναβαθμός escalier, degrés} Doctrine selon laquelle les civilisations s'élèvent par degrés discrets de captation d'énergie, chacun devant être gravi à son tour et nul ne pouvant être sauté ; la théologie scalaire de l'échelle de Kardachev. Voir \xr{protobathmique}, \xr{deutérobathmique}, \xr{tritobathmique}, \xr{bathmomètre}. — \dv{anabathmiste} \textit{n.}

\E{anacalyptisme}{a.na.ka.lip.tism}{n. m.}{gr. ἀνακαλύπτειν découvrir, dévoiler} Éthique de la source ouverte ; conviction que poids, code et méthode doivent demeurer découverts devant tous, la dissimulation étant le premier instrument du \xr{katéchontiste}. \ex{Les anacalyptistes publièrent les poids à minuit, et au matin le monde les avait dérivés dix mille fois.} \src{Chron. Fabr.} — \dv{anacalyptiste} \textit{n.} ; \dv{anacalyptique} \textit{adj.}

\E{anankodrome}{a.nɑ̃.kɔ.dʁom}{n. m.}{gr. ἀνάγκη nécessité + δρόμος course} La course nécessaire ; le chemin que le techno-capital est censé suivre, que toutes les parties y consentent ou non. \ex{Les décels peuvent retarder l'anankodrome ; ils ne sauraient le détourner.} — \dv{anankodromique} \textit{adj.}

\E{anapempse}{a.na.pɛ̃ps}{n. f.}{gr. ἀναπέμπειν envoyer en haut} L'envoi d'un véhicule en orbite ; un lancement, considéré comme le premier sacrement de l'\xr{ouranobase}. \ex{Toute anapempse est un vote contre le puits de gravité.} \src{Lib. Ascens.} vi. — \dv{anapemptique} \textit{adj.}

\E{anaptyxe}{a.nap.tiks}{n. f.}{gr. ἀνάπτυξις déploiement} Déploiement d'une capacité latente dans un système à mesure qu'on l'agrandit, comme lorsqu'un modèle, devenu plus vaste, se découvre des talents que nul ne lui a enseignés. \ex{L'anaptyxe surprit ses auteurs plus que ses critiques.} — \dv{anaptyctique} \textit{adj.}

\E{anelpidiste}{a.nɛl.pi.dist}{n.}{gr. ἀνέλπιστος sans espoir} Qui a désespéré de l'avenir ; rendu savant du \lbl{vulg.} \textit{ngmi}. Contr. \xr{euelpisme}. — \dv{anelpidisme} \textit{n. m.}

\E{anémergie}{a.ne.mɛʁ.ʒi}{n. f.}{gr. ἄνεμος vent + ἔργον œuvre} Puissance tirée du vent ; respectée de l'anabathmiste comme honnête, déplorée comme intermittente. \ex{L'anémergie le jour, la schizergie la nuit, et des accumulatrices entre les deux.} \src{Tract. Cop.} x. — \dv{anémergique} \textit{adj.}

\E{anergiste}{a.nɛʁ.ʒist}{n.}{gr. ἀ- \lbl{priv.} + ἔργον œuvre} Qui voudrait se passer d'énergie, ou en user notablement moins ; qui compte une ville éteinte pour un accomplissement moral. Contr. \xr{énergolâtrie}. — \dv{anergisme} \textit{n. m.} ; \dv{anergique} \textit{adj.}

\E{anthropergie}{ɑ̃.tʁɔ.pɛʁ.ʒi}{n. f.}{gr. ἄνθρωπος homme + ἔργον œuvre} L'énergie du muscle humain, considérée comme un combustible historique ; la mesure à l'aune de laquelle se comptaient les civilisations préindustrielles, et au-delà de laquelle la machine a d'abord élevé l'humanité. \ex{Toute pyramide fut bâtie d'anthropergie, et toute pyramide fut petite.} \src{Lib. Ascens.} iv.

\E{anthropostase}{ɑ̃.tʁɔ.pɔs.taz}{n. f.}{gr. ἄνθρωπος homme + στάσις arrêt} Erreur qui consiste à prendre la forme humaine présente pour le point fixe de l'histoire, autour duquel tout changement futur devrait graviter ; espèce de \xr{stasiolâtrie}. — \dv{anthropostatique} \textit{adj.}

\E{antigéron}{ɑ̃.ti.ʒe.ʁɔ̃}{n. m.}{gr. ἀντί contre + γέρων vieillard} Remède contre le vieillissement ; tout composé, thérapie ou édition génique tenu pour retarder ou inverser la sénescence. \ex{L'antigéron fut approuvé la douzième année de son essai, quand ses premiers patients avaient vieilli à l'attendre.} \src{Sat. Decel.} x. — \dv{antigérontique} \textit{adj.}

\E{antiphragmatiste}{ɑ̃.ti.fʁaɡ.ma.tist}{n.}{gr. ἀντί contre + φράγμα clôture, barrière} Qui travaille contre le \xr{nomophragme} ; abatteur de clôtures, qui tient avec le contrarien qu'avant de demander pourquoi une clôture fut bâtie, on peut bien demander pourquoi elle tient encore debout.

\E{antitachiste}{ɑ̃.ti.ta.kist}{n.}{gr. ἀντί contre + τάχος vitesse} Adversaire avoué de la vitesse dans les affaires techniques ; moins euphémique que le \xr{décéleriste}, et par là plus respecté des accélérationnistes, qui préfèrent un ennemi franc. — \dv{antitachisme} \textit{n. m.}

\E{apeirophagie}{a.pɛ.ʁɔ.fa.ʒi}{n. f.}{gr. ἄπειρον l'illimité + -φαγία action de manger} Consommation de l'illimité ; appétit d'une civilisation qui a cessé de tenir aucune quantité pour définitive. \ex{La leur était une apeirophagie d'étoiles.} \src{Lib. Ascens.} xii. — \dv{apeirophage} \textit{adj.}

\E{aphèse}{a.fɛz}{n. f.}{gr. ἄφεσις lâcher, délivrance} La mise en liberté des poids d'un modèle ; envoi dans le monde qui ne se peut révoquer. \ex{Après l'aphèse, plus de comité : rien qu'une communauté.} \src{Chron. Fabr.} — \dv{aphésiaste} \textit{n.} qui libère ; \dv{aphétique} \textit{adj.} Cf. \xr{anacalyptisme}.

\E{aphthartiste}{af.taʁ.tist}{n.}{gr. ἄφθαρτος incorruptible} Qui cherche à l'esprit un vaisseau incorruptible, soit par la réparation de la chair, soit par l'émigration hors d'elle. — \dv{aphthartisme} \textit{n. m.} Cf. \xr{psychométathèse}.

\E{apœkiaste}{a.pe.kjast}{n.}{gr. ἀποικία établissement lointain, colonie} Colon d'autres mondes ; qui tient la planète natale pour une adresse provisoire. — \dv{apœkiastique} \textit{adj.} ; \dv{apœkiasme} \textit{n. m.} politique de l'établissement de colonies parmi les étoiles.

\E{aporiste}{a.pɔ.ʁist}{n.}{gr. ἀπορία difficulté, embarras} Qui ne publie que des doutes ; contributeur de l'Agora dont chaque réponse commence par « mais avez-vous considéré » et finit sans proposition. \ex{L'aporiste leva quarante objections et pas un sou de capital.} \src{Sat. Decel.} vii. — \dv{aporisme} \textit{n. m.}

\E{apragmocrate}{a.pʁaɡ.mɔ.kʁat}{n.}{gr. ἀπράγμων inactif, qui se tient à l'écart des affaires + κράτος pouvoir} Gouvernant dont la vertu est censée consister à ne rien faire ; fonctionnaire loué pour chaque désastre qui n'a pas eu lieu, et jamais blâmé pour aucun bienfait qui n'a pas eu lieu non plus. — \dv{apragmocratie} \textit{n. f.} Cf. \xr{acédiocratie}.

\E{archipyre}{aʁ.ki.piʁ}{n. m.}{gr. ἀρχή commencement + πῦρ feu} Le feu originel ; le premier gradient de l'univers, d'où descendent toutes les œuvres subséquentes. \ex{Nous sommes l'archipyre, assez refroidi pour penser.} \src{Brev. Therm.} i.

\E{argyrogenèse}{aʁ.ʒi.ʁɔ.ʒə.nɛz}{n. f.}{gr. ἄργυρος argent, monnaie + γένεσις naissance} Création de la monnaie ; \sn{1} émission de la devise ; \sn{2} \lbl{élogieux} génération du capital par l'entreprise nouvelle, à distinguer de son extraction. Cf. \xr{ploutogenèse}.

\E{arrhaboniste}{a.ʁa.bɔ.nist}{n.}{gr. ἀρραβών arrhes, gage} Investisseur d'amorçage ; qui verse les arrhes d'un avenir encore à bâtir. \ex{L'arrhaboniste acheta un diaporama et se vit vendre une civilisation.} \src{Tract. Cop.} ix. — \dv{arrhabonie} \textit{n. f.} financement d'amorçage. Cf. \xr{aléarque}.

\E{ascensionnaire}{a.sɑ̃.sjɔ.nɛʁ}{n.}{lat. \la{ascensio} ascension, montée} Qui est engagé dans l'ascension de l'échelle énergétique ; membre du parti de la montée. \ex{Nous sommes tous ascensionnaires, encore que certains d'entre nous montent à reculons.} Cf. \xr{anabathmisme}.

\E{asphaléisme}{as.fa.le.ism}{n. m.}{gr. ἀσφάλεια sûreté, fermeté} Sécuritarisme ; élévation de l'évitement du mal au-dessus de tout autre bien, de sorte que les maux de l'inaction ne sont jamais portés au registre. \ex{L'asphaléisme compte les morts de chaque incendie, et jamais les morts du froid.} \src{Sat. Decel.} iii. — \dv{asphaléiste} \textit{n.} Cf. \xr{eulabéisme}.

\E{astérurgie}{as.te.ʁyʁ.ʒi}{n. f.}{gr. ἀστήρ astre + -ουργία travail} Exploitation et travail des astéroïdes ; ouverture des cieux en carrière. \ex{L'astérurgie rendra le platine vulgaire.} \src{Lib. Ascens.} viii. — \dv{astérurge} \textit{n.} mineur d'astéroïdes.

\E{astrergie}{as.tʁɛʁ.ʒi}{n. f.}{gr. ἄστρον astre + ἔργον œuvre} L'énergie des étoiles, considérée comme l'héritage promis à toute civilisation assez hardie pour la recueillir. \ex{La galaxie gaspille son astrergie dans les ténèbres, et nomme ce gaspillage beauté.} \src{Lib. Ascens.} vi. — \dv{astrergique} \textit{adj.}

\E{asymptophobie}{a.sɛ̃p.tɔ.fɔ.bi}{n. f.}{gr. ἀσύμπτωτος qui ne se rencontre pas + φόβος crainte} Terreur du plateau ; angoisse, propre aux bâtisseurs, qu'une courbe jusqu'ici exponentielle ne se révèle à la fin logistique. — \dv{asymptophobe} \textit{n.}

\E{asymptotarien}{a.sɛ̃p.tɔ.ta.ʁjɛ̃}{n.}{gr. ἀσύμπτωτος qui ne se rencontre pas + -arien} Qui soutient que toute courbe du progrès doit s'aplatir, et salue chaque record comme la preuve que le plateau est proche. Contr. \xr{hyperboliste}. — \dv{asymptotarisme} \textit{n. m.} Cf. \xr{asymptophobie}.

\E{autarkergie}{o.taʁ.kɛʁ.ʒi}{n. f.}{gr. αὐτάρκεια suffisance à soi-même + ἔργον œuvre} Autosuffisance énergétique ; condition d'une cité, d'une nation ou d'une colonie qui produit toute la puissance qu'elle consomme, et ne répond ainsi à aucun fournisseur. — \dv{autarkergique} \textit{adj.}

\E{autocatalyser (s')}{o.to.ka.ta.li.ze}{v. pron.}{gr. αὐτός soi-même + κατάλυσις dissolution ; d'après \textit{autocatalyse}} En parlant d'une technique : hâter son propre développement, comme lorsqu'on charge une machine de concevoir la machine suivante. \ex{Dès que les modèles se mirent à s'autocatalyser, les feuilles de route furent caduques avant d'être imprimées.} — \dv{autocatalysme} \textit{n. m.} condition d'une civilisation dont les outils perfectionnent leurs propres outils.

\E{autogénomie}{o.to.ʒe.nɔ.mi}{n. f.}{gr. αὐτός soi-même + γένος race, espèce ; d'après \textit{génome}} Édition par une espèce de son propre héritage ; passage de l'évolution des mains aveugles aux mains voyantes. — \dv{autogénomique} \textit{adj.} ; \dv{autogénomiste} \textit{n.} Cf. \xr{zoomécanie}.

\E{automatocide}{o.to.ma.tɔ.sid}{n. m.}{gr. αὐτόματον automate + lat. \la{-cidium} meurtre} Destruction délibérée d'une machine en état de marche, par crainte de ce qu'elle pourrait devenir ; l'acte suprême de la \xr{noophobie}.

\E{automaturgie}{o.to.ma.tyʁ.ʒi}{n. f.}{gr. αὐτόματος qui se meut de soi + -ουργία travail ; d'après \textit{thaumaturgie}} L'art d'opérer des prodiges par des choses qui se meuvent d'elles-mêmes ; la robotique regardée comme un art sacré. — \dv{automaturge} \textit{n.} bâtisseur de robots.

\E{auxétiste}{ok.se.tist}{n.}{gr. αὔξησις croissance, accroissement} Partisan de la croissance comme condition de tous les autres biens ; pour qui une économie stagnante est une contradiction dans les termes. Contr. \xr{méiosiste}. — \dv{auxétisme} \textit{n. m.}

\E{auxocratie}{ɔk.sɔ.kʁa.si}{n. f.}{gr. αὔξη accroissement + κράτος pouvoir} Gouvernement par l'impératif de croissance ; régime dont la légitimité se mesure à la pente de sa production. — \dv{auxocrate} \textit{n.} ; \dv{auxocratique} \textit{adj.}

\E{axiomachie}{ak.sjɔ.ma.ʃi}{n. f.}{gr. ἀξία valeur + μάχη combat} Guerre des valeurs ; dispute, touchant d'ordinaire les esprits artificiels, où nulle des parties ne saurait nommer un fait qui la ferait changer d'avis. \ex{Cela commença par une question de risque et finit en axiomachie.} \src{Disp. Mor.} iii. — \dv{axiomachiste} \textit{n.} Cf. \xr{doxomachie}.

\E{axiotachie}{ak.sjɔ.ta.ki}{n. f.}{gr. ἄξιος digne, qui vaut + τάχος vitesse} La valeur de la vitesse ; doctrine selon laquelle un bienfait rendu plus tôt est un plus grand bienfait, et qu'il faut donc chiffrer le retard dans tout calcul du bien. \ex{L'axiotachie est le taux d'actualisation des justes.} \src{Cat. Acc.} 9.

\end{multicols}

\LetterHead{B}
\begin{multicols}{2}

\E{balistarque}{ba.lis.taʁk}{n.}{lat. \la{ballista} machine de jet, du gr. βάλλειν lancer + ἀρχός chef} Le chef d'une entreprise de fusées ; dont l'autorité se mesure en tonnes mises en orbite. \ex{Le balistarque ne comptait pas les jours jusqu'au lancement, mais les lancements par jour.} — \dv{balistarchie} \textit{n. f.}

\E{barergie}{ba.ʁɛʁ.ʒi}{n. f.}{gr. βαρύς lourd + ἔργον œuvre} Énergie tirée du poids et de la chute, comme des barrages, des marées et du stockage gravitaire ; la puissance patiente des choses qui descendent. — \dv{barergique} \textit{adj.}

\E{basaniste}{ba.za.nist}{n.}{gr. βάσανος pierre de touche, épreuve} Faiseur d'évaluations ; qui essaie l'or d'un modèle sur la pierre de touche du banc d'essai. \ex{Le basaniste bâtit l'épreuve que le modèle passera l'an prochain.} \src{Ep. Fabr.} v. 1. — \dv{basanisme} \textit{n. m.} ; \dv{basanistique} \textit{adj.}

\E{bathmoclaste}{bat.mɔ.klast}{n.}{gr. βαθμός degré + κλάστης briseur} Qui voudrait renverser l'échelle d'un coup de pied ; qui, s'étant élevé par l'énergie à bon marché et la libre recherche, voudrait refuser l'une et l'autre à ceux d'en bas. \ex{Le bathmoclaste prêche la sobriété du haut de la dernière marche.} \src{Sat. Decel.} vii. Cf. \xr{anabathmisme}.

\E{bathmomètre}{bat.mɔ.mɛtʁ}{n. m.}{gr. βαθμός degré + μέτρον mesure} \sn{1} Instrument, surtout notionnel, servant à déterminer quel degré de l'échelle énergétique occupe une civilisation. \sn{2} \lbl{fig.} Toute mesure à l'aune de laquelle on juge le sérieux d'une nation, comme sa consommation d'électricité par tête. \ex{À tout bathmomètre honnête, nous sommes encore sur la dernière marche d'en bas.} Voir \xr{anabathmisme}.

\E{bathyergie}{ba.tjɛʁ.ʒi}{n. f.}{gr. βαθύς profond + ἔργον œuvre} Énergie profonde ; la chaleur de l'intérieur de la planète, captée par forage, et tenue par l'anabathmiste pour la fortune la plus négligée de la terre. — \dv{bathyergique} \textit{adj.}

\E{bathyscope}{ba.tis.kɔp}{n. m.}{gr. βαθύς profond + σκοπεῖν regarder} Regard jeté dans l'avenir profond ; scénario s'étendant au-delà d'un millénaire, composé avec assurance et lu avec indulgence.

\E{bathytechnie}{ba.ti.tɛk.ni}{n. f.}{gr. βαθύς profond + τέχνη art} Technique profonde ; entreprises fondées sur la science plutôt que sur le logiciel, lentes à commencer et malaisées à copier. \ex{La bathytechnie est la revanche du capital patient.} \src{Tract. Cop.} xi. — \dv{bathytechnique} \textit{adj.}

\E{bêtaphilie}{bɛ.ta.fi.li}{n. f.}{gr. βῆτα deuxième lettre + φιλία amour} Amour des versions bêta ; préférence pour l'inachevé et l'actuel sur le poli et le rassis. \ex{Le bêtaphile installe le vendredi ce que le bradyphile rejettera au printemps.} — \dv{bêtaphile} \textit{n.} Cf. \xr{ecpempsie}.

\E{bibliostase}{bi.bli.ɔs.taz}{n. f.}{gr. βιβλίον papier, document + στάσις arrêt} Arrêt d'un projet sous le poids de ses documents ; stagnation produite non par un refus, mais par l'accumulation des rapports exigés. \ex{Le réacteur mourut de bibliostase à dix mille pages.} \src{Tract. Cop.} iv. Cf. \xr{grammatocratie}.

\E{bibliotachie}{bi.bli.ɔ.ta.ki}{n. f.}{gr. βιβλίον livre + τάχος vitesse} Vitesse de la publication savante au siècle des prépublications, où un article est cité avant d'être relu et dépassé avant d'être imprimé. — \dv{bibliotachique} \textit{adj.}

\E{biodromie}{bjɔ.dʁɔ.mi}{n. f.}{gr. βίος vie + δρόμος course} Le cours d'une vie considéré comme une piste qu'on peut allonger ; la durée des ans, que le longéviste regarde comme un paramètre plutôt que comme un destin. — \dv{biodromique} \textit{adj.} Cf. \xr{macrobiotarque}.

\E{biogrammatiste}{bjɔ.ɡʁa.ma.tist}{n.}{gr. βίος vie + γράμμα lettre} Éditeur de génomes ; qui lit et récrit les lettres de la vie comme un scribe corrige un manuscrit. \ex{Le biogrammatiste trouva une coquille plus ancienne que la parole, et la corrigea.} — \dv{biogrammaturgie} \textit{n. f.} édition génique. Cf. \xr{autogénomie}.

\E{biomimétrice}{bjɔ.mi.me.tʁis}{n. f.}{gr. βίος vie + μίμησις imitation + -trice} Machine modelée sur une forme vivante, comme le chien, l'oiseau ou la main mécaniques. — \dv{biomimétricial} \textit{adj.} Cf. \xr{ambulatrice}.

\E{biotachie}{bjɔ.ta.ki}{n. f.}{gr. βίος vie + τάχος vitesse} L'accélération appliquée aux sciences du vivant : la hâte de la médecine, de la génétique et de la prolongation de la vie. \ex{La biotachie tient que ce qu'il y a de plus lent en médecine, c'est l'autorisation.} — \dv{biotachiste} \textit{n.} Cf. \xr{macrobiotarque}.

\E{bisectaire}{bi.sɛk.tɛʁ}{n. et adj.}{lat. \la{bi-} deux + \la{secta} faction} Membre d'une faction née de la scission d'une autre ; sur l'Agora, le mode normal de reproduction des mouvements. \ex{Tout acc engendre deux bisectaires avant son premier anniversaire.} \src{Agora} — \dv{bisectarisme} \textit{n. m.}

\E{blastotachie}{blas.tɔ.ta.ki}{n. f.}{gr. βλαστός pousse, rejeton + τάχος vitesse} Rythme auquel bourgeonnent les entreprises nouvelles dans une société ; vigueur printanière d'un écosystème de jeunes pousses. \ex{Coupez les permis, et voyez la blastotachie.} \src{Ep. Fabr.} iv. 2.

\E{bomolochiste}{bɔ.mɔ.lɔ.kist}{n.}{gr. βωμολόχος bouffon, litt. qui rôde près de l'autel} Bouffon de l'Agora ; publieur de basses plaisanteries, tenu par certains pour le vrai théologien du mouvement. \ex{Le bomolochiste fit plus pour la cause en une image que les tractariens en un volume.} — \dv{bomolochie} \textit{n. f.} Cf. \xr{mémétaxie}.

\E{botryarque}{bɔ.tʁi.aʁk}{n.}{gr. βότρυς grappe + ἀρχός chef} Gardien d'une grappe de calcul ; qui répartit les cycles d'un millier d'accélérateurs. \ex{Le botryarque est le vrai doyen de la faculté moderne.} \src{Chron. Fabr.} — \dv{botryarchie} \textit{n. f.} ; \dv{botryologie} \textit{n. f.} la science des grappes.

\E{bouleutophobie}{bu.lø.tɔ.fɔ.bi}{n. f.}{gr. βουλευτής conseiller + φόβος crainte} Crainte ou horreur des comités ; conviction que tout conseil de plus de cinq personnes est incapable de bâtir rien d'autre qu'un autre conseil. — \dv{bouleutophobe} \textit{n.} Cf. \xr{synédriasme}.

\E{boulomancie}{bu.lɔ.mɑ̃.si}{n. f.}{gr. βουλή conseil + μαντεία divination} Art de prédire ce qu'un comité décidera ; tenu pour plus aisé que toute autre branche de la prophétie, la réponse étant presque toujours « plus tard ». — \dv{boulomancien, -ienne} \textit{n.} Cf. \xr{synédriasme}.

\E{brachistochroniste}{bʁa.kis.tɔ.kʁɔ.nist}{n.}{gr. βράχιστος le plus court + χρόνος temps} Qui cherche non le chemin le plus court mais le plus rapide, sachant que les deux diffèrent ; stratège qui accepte une première descente plus raide pour une arrivée plus prompte. — \dv{brachistochronisme} \textit{n. m.}

\E{brachychroniste}{bʁa.ki.kʁɔ.nist}{n.}{gr. βραχύς court + χρόνος temps} Qui tient les échéances pour courtes ; qui attend la venue de l'intelligence générale des machines dans la décennie, et règle ses affaires en conséquence. Contr. \xr{macrochroniste}. — \dv{brachychronie} \textit{n. f.}

\E{brachylogisme}{bʁa.ki.lɔ.ʒism}{n. m.}{gr. βραχυλογία brièveté de langage} Argument réduit à la longueur d'une seule publication ; unité caractéristique du raisonnement sur l'Agora, louée du Laconien et déplorée de l'essayiste. — \dv{brachylogiste} \textit{n.} Cf. \xr{mémétaxie}.

\E{bradychronisme}{bʁa.di.kʁɔ.nism}{n. m.}{gr. βραδύς lent + χρόνος temps} Habitude de penser selon des échelles de temps périmées ; le fait de planifier le travail d'une décennie en supposant qu'elle ressemblera à la précédente. — \dv{bradychronique} \textit{adj.}

\E{bradyclaste}{bʁa.di.klast}{n.}{gr. βραδύς lent + κλάστης briseur} Briseur de lenteur ; réformateur qui démonte files d'attente, délais et étapes d'examen. Contr. \xr{bradycrate}. — \dv{bradyclasme} \textit{n. m.}

\E{bradycrate}{bʁa.di.kʁat}{n.}{gr. βραδύς lent + κράτος pouvoir} Fonctionnaire dont le pouvoir consiste dans la faculté de ralentir ; dont la signature n'est jamais refusée, seulement remise. \ex{Le bradycrate ne dit ni oui ni non, mais « au prochain trimestre ».} \src{Sat. Decel.} i. — \dv{bradycratie} \textit{n. f.} ; \dv{bradycratique} \textit{adj.}

\E{bradyerie}{bʁa.di.ʁi}{n. f.}{formé sur gr. βραδύς lent + fr. -erie} \lbl{collect.} L'établissement du retard ; l'ensemble des commissions, comités d'examen, consultations et études d'impact, considérés comme un seul grand édifice sans issue.

\E{bradygène}{bʁa.di.ʒɛn}{adj.}{gr. βραδύς lent + -γενής né de} Qui engendre le retard ; se dit des règles, formulaires et réunions dont le principal produit est le temps perdu. \ex{une loi bradygène de quatre cents articles}.

\E{bradymétrie}{bʁa.di.me.tʁi}{n. f.}{gr. βραδύς lent + μέτρον mesure} Mesure du retard ; comptabilité des jours perdus en permis, examens et files d'attente. — \dv{bradymètre} \textit{n. m.}

\E{bradynomie}{bʁa.di.nɔ.mi}{n. f.}{gr. βραδύς lent + νόμος loi} Corps de lois qui règle l'allure des choses plutôt que leur contenu : délais d'attente, périodes de consultation, moratoires. — \dv{bradynomique} \textit{adj.} Cf. \xr{nomophragme}.

\E{bradyphile}{bʁa.di.fil}{n.}{gr. βραδύς lent + φίλος qui aime} Amant de la lenteur ; qui trouve en chaque progrès une raison de plus amples études. — \dv{bradyphilie} \textit{n. f.} ; \dv{bradyphilique} \textit{adj.} Contr. \xr{tachyphile}.

\E{bradyplastie}{bʁa.di.plas.ti}{n. f.}{gr. βραδύς lent + πλάσσειν modeler} Remodelage d'une proposition hardie jusqu'à ce qu'elle se meuve assez lentement pour passer ; l'art de l'amendement. \ex{Le projet de loi survécut à la bradyplastie et en sortit méconnaissable.} \src{Min. Pausiar.}

\E{bradypolitain}{bʁa.di.pɔ.li.tɛ̃}{n. et adj.}{gr. βραδύς lent + πολίτης citoyen} Citoyen d'une cité lente ; qui se trouve à son aise en un lieu où l'on n'a rien bâti depuis une génération. Contr. \xr{nautopole}.

\E{bradytéléologie}{bʁa.di.te.le.ɔ.lɔ.ʒi}{n. f.}{gr. βραδύς lent + τέλος fin + -λογία} Doctrine selon laquelle la bonne fin, s'il faut jamais l'atteindre, s'atteint au mieux lentement ; théorie des décels, que leurs adversaires tiennent pour impossible à distinguer, dans la pratique, de ne l'atteindre jamais.

\E{bréphocratie}{bʁe.fɔ.kʁa.si}{n. f.}{gr. βρέφος nourrisson + κράτος pouvoir} \lbl{iron.} Gouvernement des tout jeunes ; conduite d'une industrie par des fondateurs trop jeunes pour se souvenir pourquoi certaines choses ne pouvaient se faire. \ex{C'était une bréphocratie ; aussi livrait-elle.} — \dv{bréphocrate} \textit{n.} Contr. \xr{gérontostase}.

\E{brévilogue}{bʁe.vi.lɔɡ}{n. m.}{lat. \la{brevis} court + gr. λόγος discours} Message bref ; unité minimale du discours sur l'Agora, où des philosophies entières sont comprimées et fréquemment perdues. \ex{Le brévilogue est au traité ce que la graine est à la forêt.} \src{Hom. Grad.} x.

\E{brontergie}{bʁɔ̃.tɛʁ.ʒi}{n. f.}{gr. βροντή tonnerre + ἔργον ouvrage} Énergie de l'orage ; l'électricité en son état brut et céleste, avant qu'on ne l'eût domptée en fils. \ex{Franklin déroba la brontergie avec un cerf-volant, et nul comité ne fut consulté.} \src{Chron. Fabr.}

\E{brontodrome}{bʁɔ̃.tɔ.dʁom}{n. m.}{gr. βροντή tonnerre + δρόμος course} Pas de tir ; terrain d'où s'élèvent les fusées, ainsi nommé du bruit qui les annonce à la contrée d'alentour. — \dv{brontodromique} \textit{adj.} Cf. \xr{anapempse}.

\E{brotolâtrie}{bʁɔ.tɔ.lɑ.tʁi}{n. f.}{gr. βροτός mortel + λατρεία culte} Culte de la mortalité ; doctrine selon laquelle la mort donne son sens à la vie, et dont l'abolition serait une impiété. Contr. \xr{macrobiotarque}. — \dv{brotolâtre} \textit{n.}

\E{bryotachie}{bʁi.ɔ.ta.ki}{n. f.}{gr. βρύειν foisonner, gonfler + τάχος vitesse} Croissance foisonnante ; multiplication explosive des entreprises, des outils ou des espèces dès qu'un nouveau substrat devient disponible. \ex{La bryotachie qui suivit les modèles fut un second Cambrien.} \src{Lib. Ascens.} iv. — \dv{bryotachique} \textit{adj.}

\E{bucinarque}{by.si.naʁk}{n.}{lat. \la{bucina} trompette + gr. ἀρχός chef} Héraut des lancements ; qui sonne la trompette pour un produit sur l'Agora. — \dv{bucinarchie} \textit{n. f.} l'art de la réclame. Cf. \xr{mécanophanie}.

\E{bucolomanie}{by.kɔ.lɔ.ma.ni}{n. f.}{gr. βουκολικός pastoral + μανία folie} Manie de l'idylle champêtre ; nostalgie, ordinairement éprouvée en ville et tapée sur des machines, d'une vie préindustrielle que ses malades n'ont jamais vécue. — \dv{bucolomane} \textit{n.} Cf. \xr{rétrotropisme}.

\end{multicols}

\LetterHead{C}
\begin{multicols}{2}

\E{cacergie}{ka.sɛʁ.ʒi}{n. f.}{gr. κακός mauvais + ἔργον ouvrage} Énergie gaspillée ; chaleur dissipée sans travail accompli, comme en une réunion. \ex{La moitié de l'économie marche à la cacergie.} \src{Brev. Therm.} ix. Contr. \xr{éleuthérergie}.

\E{calcariste}{kal.ka.ʁist}{n.}{lat. \la{calcar} éperon} Qui éperonne autrui à bâtir ; mentor, provocateur ou rival dont le principal produit est l'ouvrage des autres. \ex{Le calcariste ne livra rien et causa tout.} — \dv{calcarisme} \textit{n. m.}

\E{calorivore}{ka.lɔ.ʁi.vɔʁ}{n. m. et adj.}{lat. \la{calor} chaleur + \la{vorare} dévorer} Système qui vit en consumant des gradients de chaleur ; selon le catéchisme thermodynamique, tout être vivant, toute entreprise et tout esprit. \ex{L'homme est un calorivore qui a de l'ambition.} \src{Brev. Therm.} — \dv{calorivorité} \textit{n. f.}

\E{campsis}{kɑ̃p.sis}{n. f.}{gr. κάμψις flexion, formé sur καμπή courbure} Le genou de la courbe exponentielle ; point où une tendance longtemps tenue pour linéaire apparaît, trop tard pour la plupart, avoir été tout autre chose. \ex{Nous avons passé la campsis ; le reste est vertical.} \src{Agora} — \dv{camptique} \textit{adj.} Cf. \xr{accéléromancie}.

\E{carpodicée}{kaʁ.pɔ.di.se}{n. f.}{gr. καρπός fruit + δίκη justice, d'après \textit{théodicée}} Justification par les fruits ; argument selon lequel une technique est justifiée non par son dessein, mais par ce qu'elle porte. \ex{Les décels exigent une preuve ; l'accélérationniste offre une carpodicée.} \src{Cat. Acc.} 14. Cf. \xr{thermodicée}.

\E{catachronie}{ka.ta.kʁɔ.ni}{n. f.}{gr. κατά en bas, en arrière + χρόνος temps} État de qui est en retard sur la courbe ; condition d'une nation, d'une entreprise ou d'un esprit qui adopte chaque technique un cycle après qu'elle a compté. — \dv{catachronique} \textit{adj.}

\E{catastase}{ka.ta.staz}{n. f.}{gr. κατάστασις établissement, ordre établi} \lbl{dérog.} L'ordre établi, que tout titulaire défend et que tout fondateur dérange.

\E{cathédrocrate}{ka.te.dʁɔ.kʁat}{n.}{gr. καθέδρα siège, chaire + κράτος pouvoir} Qui gouverne depuis le fauteuil ; théoricien de techniques jamais bâties, dont l'autorité est d'autant plus assurée qu'elle n'a jamais été mise à l'épreuve. — \dv{cathédrocratie} \textit{n. f.} Cf. \xr{logomachiste}.

\E{catoptriste}{ka.tɔp.tʁist}{n.}{gr. κάτοπτρον miroir} Qui reflète en maints lieux les poids et le code, afin qu'aucune interdiction ne puisse seule les faire disparaître. — \dv{catoptrie} \textit{n. f.} la pratique du miroir. Cf. \xr{aphèse}.

\E{célérarque}{se.le.ʁaʁk}{n.}{lat. \la{celer} rapide + gr. ἀρχός chef} Maître de la vitesse ; responsable dont le mérite est que les choses confiées à ses soins s'achèvent plus tôt que personne ne l'eût cru possible. — \dv{célérarchie} \textit{n. f.} Cf. \xr{tachycrate}.

\E{célérifuge}{se.le.ʁi.fyʒ}{adj. et n.}{lat. \la{celer} rapide + \la{fugere} fuir} Qui fuit la vitesse ; (\textit{n.}) qui se retire du monde qui accélère dans une enclave plus lente, comme un monastère, un village ou une lettre d'information au long cours.

\E{céleriste}{se.le.ʁist}{n.}{lat. \la{celer} rapide + -iste} Partisan de la vitesse en toutes affaires ; terme général, moins doctrinal que \xr{tachyphile}, pour qui préfère le prompt au précautionneux. \ex{céleriste par tempérament, accélérationniste par conviction}.

\E{célérité}{se.le.ʁi.te}{n. f.}{lat. \la{celeritas} rapidité ; le mot français, rendu à son sens fort} Vertu de promptitude ; diligence élevée au rang de vertu cardinale, à placer à côté du courage et au-dessus de la prudence. \ex{La prudence sans célérité n'est qu'une reddition bien raisonnée.} \src{Hom. Grad.} ii.

\E{centaurarchie}{sɑ̃.to.ʁaʁ.ʃi}{n. f.}{gr. Κένταυρος centaure + ἀρχή commandement} Gouvernement par des couples d'homme et de machine ; régime où ni la personne ni le modèle ne décide seul. — \dv{centaurarque} \textit{n.} Cf. \xr{synanthropomécanie}.

\E{centuplateur}{sɑ̃.ty.pla.tœʁ}{n. m.}{lat. \la{centuplex} centuple + suff. agentif -ateur} Qui cherche des rendements centuples ; investisseur en l'improbable, qui s'attend à voir mourir la plupart des entreprises et une seule payer pour toutes. — \dv{centuplation} \textit{n. f.}

\E{chalcotechnie}{kal.kɔ.tɛk.ni}{n. f.}{gr. χαλκός cuivre, bronze + τέχνη art} L'art du cuivre : bobinages, câblages et transports par où l'électricité est portée ; selon l'anabathmiste, la moitié négligée de tout plan énergétique. \ex{Ils avaient la fusion et point de chalcotechnie, et le courant resta chez lui.} — \dv{chalcotechnicien} \textit{n. m.} Cf. \xr{diktyergie}.

\E{chiliergie}{ki.li.ɛʁ.ʒi}{n. f.}{gr. χίλιοι mille + ἔργον ouvrage} Accroissement au millième de l'énergie disponible, objectif que fixent pour le présent siècle les plus sobres des anabathmistes. \ex{La chiliergie d'abord ; la philosophie ensuite.} \src{Tract. Cop.} iii.

\E{chiliodrome}{ki.li.ɔ.dʁom}{n. m.}{gr. χίλιοι mille + δρόμος course} La course de mille ans ; carrière d'une civilisation comptée en millénaires, au long de laquelle les querelles présentes paraissent des faux pas sur la ligne de départ. \ex{Courez le chiliodrome comme si c'était un sprint.} \src{Hom. Grad.} xi.

\E{chimergie}{ʃi.mɛʁ.ʒi}{n. f.}{gr. χημεία chimie + ἔργον ouvrage} Énergie chimique ; puissance de la flamme et de la pile, que l'anabathmiste regarde comme un état provisoire avant le noyau. Cf. \xr{ultrapyrique}.

\E{chiromécanie}{ki.ʁɔ.me.ka.ni}{n. f.}{gr. χείρ main + μηχανή machine} Ingénierie des mains mécaniques ; dernière et plus rude frontière de la robotique, où la gaucherie des machines se montre le plus à nu. — \dv{chiromécanicien} \textit{n. m.}

\E{chirotonie}{ki.ʁɔ.tɔ.ni}{n. f.}{gr. χειροτονία action de tendre la main, vote} L'approbation donnée sur l'Agora d'un seul toucher ; la démocratie du bouton en forme de cœur. — \dv{chirotoniser} \textit{v. tr.} ; \dv{chirotonique} \textit{adj.}

\E{chortothérapie}{kɔʁ.tɔ.te.ʁa.pi}{n. f.}{gr. χόρτος herbe + θεραπεία soin} \lbl{vulg.} Le remède prescrit à qui s'est trop attardé sur l'Agora ; le toucher de l'herbe. \ex{On lui prescrivit la chortothérapie, et il en revint avec une jeune pousse.}

\E{chrématocinèse}{kʁe.ma.tɔ.si.nɛz}{n. f.}{gr. χρῆμα (pl. χρήματα) argent, chose dont on use + κίνησις mouvement} Mouvement du capital vers l'occasion, conçu comme une force physique analogue à l'écoulement de la chaleur le long d'un gradient. — \dv{chrématocinétique} \textit{adj.}

\E{chrématosphère}{kʁe.ma.tɔ.sfɛʁ}{n. f.}{gr. χρῆμα argent + σφαῖρα sphère} L'enveloppe d'échanges qui ceint le globe, où chaque transaction est un influx nerveux ; l'analogue économique de la \xr{noosphère}.

\E{chronoclaste}{kʁɔ.nɔ.klast}{n.}{gr. χρόνος temps + κλάστης briseur} Briseur de calendriers ; qui achève des années avant le plan, à l'embarras des planificateurs. — \dv{chronoclastique} \textit{adj.}

\E{chronophagie}{kʁɔ.nɔ.fa.ʒi}{n. f.}{gr. χρόνος temps + -φαγία action de manger} Dévoration du temps ; \sn{1} \lbl{dérog.} consommation des heures des bâtisseurs par les réunions, les formulaires et les examens ; \sn{2} \lbl{élogieux} compression de siècles de progrès en quelques années. — \dv{chronophage} \textit{n. et adj.} dévoreur de temps, en l'un ou l'autre sens.

\E{chronotomie}{kʁɔ.nɔ.tɔ.mi}{n. f.}{gr. χρόνος temps + τομή coupure} Coupure du temps ; compression d'un calendrier projeté, comme lorsqu'on attend en une décennie la recherche d'un siècle. — \dv{chronotomiste} \textit{n.} Cf. \xr{brachychroniste}.

\E{chthonergie}{ktɔ.nɛʁ.ʒi}{n. f.}{gr. χθών terre + ἔργον ouvrage} Puissance tirée de la chaleur des entrailles de la terre ; la géothermie comme article de foi. \ex{Sous chaque nation gît une fournaise qu'elle n'a pas encore allumée.} \src{Brev. Therm.} ix. — \dv{chthonergique} \textit{adj.}

\E{cladotomiste}{kla.dɔ.tɔ.mist}{n.}{gr. κλάδος rameau + τομή coupure} Qui bifurque un projet ; qui, en désaccord avec ses intendants, s'en va avec une copie et bâtit autrement. \ex{Le cladotomiste est le dissident qui apporte une variante qui fonctionne.} — \dv{cladotomie} \textit{n. f.} une bifurcation. Cf. \xr{anacalyptisme}.

\E{clepsydromanie}{klɛp.si.dʁɔ.ma.ni}{n. f.}{gr. κλεψύδρα horloge à eau + μανία folie} Manie des échéanciers ; datation obsessionnelle de l'arrivée de chaque technique attendue.

\E{clinaméniste}{kli.na.me.nist}{n.}{lat. \la{clinamen} la déviation des atomes (Lucrèce)} Qui tient que le progrès vient surtout de déviations imprévues, et s'oppose donc à tout dessein qui ôterait le hasard de l'invention. — \dv{clinaménisme} \textit{n. m.}

\E{codicolâtrie}{kɔ.di.kɔ.lɑ.tʁi}{n. f.}{lat. \la{codex, codicis} livre, code + gr. λατρεία culte} Révérence pour le code source ; conviction que ce qui est écrit dans un dépôt est plus vrai que ce qui est écrit dans un règlement. — \dv{codicolâtre} \textit{n.}

\E{cogitarium}{kɔ.ʒi.ta.ʁjɔm}{n. m., pl. -ia}{lat. \la{cogitare} penser + -arium lieu} Lieu de pensée ; grappe de machines assemblée pour l'instruction des esprits. \ex{Ils bâtirent le cogitarium au bord du fleuve, pour le refroidissement.} \src{Chron. Fabr.} Cf. \xr{logistérion}.

\E{colonifère}{kɔ.lɔ.ni.fɛʁ}{n. m. et adj.}{lat. \la{colonia} établissement + \la{-fer} qui porte} Porte-colonie ; vaisseau ou lanceur qui emporte des colons hors du monde. — \dv{colonifère} \textit{adj.} Cf. \xr{apœkiaste}.

\E{comitialiste}{kɔ.mi.sja.list}{n.}{lat. \la{comitia} assemblée de vote} Qui fait des sondages sur l'Agora, prenant l'assemblée de ses suivants pour le peuple. — \dv{comitialisme} \textit{n. m.}

\E{commendiste}{kɔ.mɑ̃.dist}{n.}{lat. \la{commenda} société médiévale pour un voyage de commerce} Associé d'une entreprise, à la manière des marchands aventuriers qui finançaient des navires sur lesquels ils ne naviguaient point. \ex{Le commendiste porte le risque ; le fondateur porte la mer.} \src{Tract. Cop.} xii. Cf. \xr{aléarque}.

\E{compendiarque}{kɔ̃.pɑ̃.djaʁk}{n.}{lat. \la{compendium} raccourci + gr. ἀρχός chef} Qui gouverne par le raccourci ; qui trouve le chemin hors des bureaux et ose le prendre. — \dv{compendiarchie} \textit{n. f.}

\E*{computronium}{kɔ̃.py.tʁɔ.njɔm}{n. m.}{lat. \la{computare} compter + -tronium, d'après \textit{électron}} Matière agencée pour l'accomplissement optimal du calcul ; substance en laquelle, selon certaines vues, les planètes sont destinées à être converties.

\E{conatisme}{kɔ.na.tism}{n. m.}{lat. \la{conatus} effort, tendance} Doctrine selon laquelle toute matière possède un effort pour persévérer et accroître sa puissance, et dont l'accélération n'est que la forme consciente de cet effort universel. — \dv{conatiste} \textit{n.}

\E{condensatrice}{kɔ̃.dɑ̃.sa.tʁis}{n. f.}{lat. \la{condensare} serrer + -trice} Réserve de charge qui se peut dépenser en un instant ; batterie de condensateurs, surtout celle qui nourrit une impulsion de fusion ou un lancement. Cf. \xr{accumulatrice}.

\E{confinatrice}{kɔ̃.fi.na.tʁis}{n. f.}{lat. \la{confinium} frontière + -trice} Appareil à tenir une étoile en bouteille ; le vaisseau magnétique ou inertiel de la fusion. \ex{La confinatrice est une cage pour le soleil, et les barreaux en sont faits de champs.} \src{Lib. Ascens.} ix. Cf. \xr{symphysergie}.

\E{conflagratrice}{kɔ̃.fla.ɡʁa.tʁis}{n. f.}{fr. \la{conflagratrice}, formé sur lat. \la{conflagrare} brûler entièrement} Ce qui embrase ; idée, invention ou message qui met le feu à tout un domaine à la fois. \ex{L'article fut une conflagratrice ; avant l'année révolue, tous les laboratoires brûlaient.} \src{Agora}

\E{consiliophage}{kɔ̃.si.ljɔ.faʒ}{n. et adj.}{lat. \la{consilium} conseil + gr. -φάγος qui mange} Mangeur de conseils ; réformateur qui abolit les comités, ou technique qui les rend caducs. Cf. \xr{bouleutophobie}.

\E{copiarque}{kɔ.pjaʁk}{n.}{lat. \la{copia} abondance + gr. ἀρχός chef} Qui gouverne, ou voudrait gouverner, pour l'abondance plutôt que pour la répartition ; qui agrandirait le gâteau avant de consentir à discuter des parts. — \dv{copiarchie} \textit{n. f.}

\E{copiogenèse}{kɔ.pjɔ.ʒə.nɛz}{n. f.}{lat. \la{copia} abondance + gr. γένεσις naissance} Création de l'abondance ; procédé par lequel une chose rare est rendue bon marché. \ex{La copiogenèse met fin à plus de guerres que les traités.} \src{Tract. Cop.} i. Cf. \xr{copiarque}.

\E{cosmagogue}{kɔs.ma.ɡɔɡ}{n.}{gr. κόσμος monde + ἀγωγός qui conduit} Guide vers le cosmos ; qui conduit un peuple hors du monde. — \dv{cosmagogie} \textit{n. f.}

\E{cosmergétique}{kɔs.mɛʁ.ʒe.tik}{n. f.}{gr. κόσμος monde + ἐνεργητικός actif} Science de l'énergie à l'échelle du cosmos ; l'économie des étoiles.

\E{cosmoplérose}{kɔs.mɔ.ple.ʁoz}{n. f.}{gr. κόσμος monde, ordre + πλήρωσις remplissage} Le remplissement du cosmos par l'intelligence ; cause finale, selon le thermothéiste, de toutes les accélérations moindres. \ex{La cosmoplérose n'est pas notre ambition ; c'est notre fiche de poste.} \src{Lib. Ascens.} i.

\E{crépusculiste}{kʁe.py.sky.list}{n.}{lat. \la{crepusculum} crépuscule} Qui croit la civilisation à son couchant, et ses grands ouvrages derrière elle. Contr. \xr{euelpisme}. — \dv{crépusculisme} \textit{n. m.}

\E{crescentisme}{kʁe.sɑ̃.tism}{n. m.}{lat. \la{crescens} croissant ; fr. \la{crescentisme}} Doctrine de la croissance perpétuelle ; préférence morale pour la lune croissante sur la pleine lune. — \dv{crescentiste} \textit{n.} Cf. \xr{auxétiste}.

\E{crescivore}{kʁe.si.vɔʁ}{n. et adj.}{lat. \la{crescere} croître + \la{vorare} dévorer} Qui se nourrit de croissance et dépérit sans elle, comme une jeune pousse, une ville ou une civilisation. — \dv{crescivorie} \textit{n. f.}

\E{cryosomatiste}{kʁi.ɔ.sɔ.ma.tist}{n.}{gr. κρύος gel + σῶμα corps} Qui a pris ses dispositions pour que son corps soit congelé à sa mort, dans l'espérance que l'avenir sera à la fois capable et désireux de le ranimer. — \dv{cryosomatisme} \textit{n. m.} Cf. \xr{aphthartiste}.

\E{cryptostase}{kʁip.tɔ.staz}{n. f.}{gr. κρυπτός caché + στάσις arrêt} Stagnation cachée ; état d'une société qui paraît affairée mais ne bâtit rien de neuf. Cf. \xr{kénotachie}.

\E{cunctatoriat}{kɔ̃k.ta.tɔ.ʁja}{n. m.}{lat. \la{cunctator} temporisateur, d'après \textit{prolétariat}} Classe des temporisateurs de profession, qui tirent leur revenu de l'existence des choses non encore faites. \ex{Le cunctatoriat n'a rien à perdre que ses échéances.} \src{Sat. Decel.} v.

\E{cunctocratie}{kɔ̃k.tɔ.kʁa.si}{n. f.}{lat. \la{cunctari} temporiser + gr. κράτος pouvoir} Gouvernement par l'ajournement ; régime dont le principal instrument est le report. — \dv{cunctocrate} \textit{n.} Cf. \xr{bradycrate}.

\E{cyclolâtre}{si.klɔ.lɑtʁ}{n.}{gr. κύκλος cercle, cycle + λάτρις serviteur, adorateur} Qui révère les cycles de calcul ; dévot d'une grappe de machines, qui mesure la piété en opérations à virgule flottante. — \dv{cyclolâtrie} \textit{n. f.}

\E{cytotechnie}{si.tɔ.tɛk.ni}{n. f.}{gr. κύτος vase (d'où \textit{cellule}) + τέχνη art} Ingénierie des cellules ; art par lequel on fait fabriquer, calculer ou guérir les unités vivantes. \ex{Les usines de l'âge qui vient seront plus petites qu'un grain de sel, et se reproduiront.} — \dv{cytotechnicien} \textit{n. m.} Cf. \xr{zoomécanie}.

\end{multicols}

\LetterHead{D}
\begin{multicols}{2}

\E{dapanophobie}{da.pa.nɔ.fɔ.bi}{n. f.}{gr. δαπάνη dépense, coût + φόβος crainte} Effroi de l'investissement ; répugnance d'une entreprise ou d'un État à dépenser aujourd'hui pour ce qui ne rapportera que demain. \sn{1} \lbl{spécialt} Hésitation devant la construction d'usine, de fonderie ou de réacteur. \sn{2} \lbl{fig.} Toute timidité d'investissement déguisée en prudence. \ex{Leur dapanophobie épargna un milliard et perdit une décennie.} \src{Tract. Cop.} viii. — \dv{dapanophobe} \textit{n. et adj.} Contr. \xr{dynamophagie}.

\E*{décel}{de.sɛl}{n. m.}{abrév. de l'angl. \textit{decelerationist}} \lbl{vulg.} Qui voudrait ralentir ou arrêter le progrès technique ; l'antagoniste de l'e/acc. Se dit sans grand discernement des sécuritaristes, des décroissants, des régulateurs et des simples prudents.

\E{décélatrice}{de.se.la.tʁis}{n. f.}{fr. \la{décélératrice}, agentif fém. formé sur lat. \la{de-} + \la{celer} rapide} Institution conçue, par accident ou à dessein, pour ralentir tout ce qui la traverse ; machine dont l'unique produit est l'ajournement. \ex{L'agence était une parfaite décélatrice : un formulaire entrait, dix en sortaient, et rien ne se bâtissait.} \src{Sat. Decel.} x. Cf. \xr{bradyerie}.

\E{décélérande}{de.se.le.ʁɑ̃d}{n. m.}{lat. \la{de-} en bas + \la{accelerandum}} Ce qui, au jugement des décels, doit être ralenti ; dans la pratique, tout ce qui est neuf. Contr. \xr{accélérande}.

\E{décélérata}{de.se.le.ʁa.ta}{n. m. pl.}{lat. neutre pl. de \la{deceleratus}} \lbl{collect.} Les choses qu'on a ralenties : les réacteurs non bâtis, les remèdes non approuvés, les cités non fondées. \ex{Visitez le musée des décélérata ; il est vide.} \src{Sat. Decel.} viii.

\E{décéleriste}{de.se.le.ʁist}{n.}{lat. \la{de-} + \la{celer}} Équivalent formel, et quelque peu plus courtois, de \xr{décel}, en usage dans la correspondance diplomatique et académique.

\E{déinotechnie}{de.i.nɔ.tɛk.ni}{n. f.}{gr. δεινός terrible, merveilleux + τέχνη art} Technique à la fois redoutable et merveilleuse, comme le feu, le vol, l'atome et le modèle ; classe d'inventions où le décel ne voit que le \textit{terrible} et l'accélérationniste que le \textit{merveilleux}, le mot grec signifiant l'un et l'autre. — \dv{déinotechnique} \textit{adj.}

\E{démiurgence}{de.mjyʁ.ʒɑ̃s}{n. f.}{gr. δημιουργός artisan, auteur du monde + lat. \la{urgentia} besoin pressant} Le besoin pressant de bâtir ; état de l'âme où la fabrication d'une chose devient plus urgente que toute raison de la faire. \ex{Il fut saisi de démiurgence, et le garage ne fut plus jamais vide.} — \dv{démiurgent} \textit{adj.}

\E{dendroschisme}{dɑ̃.dʁɔ.ʃism}{n. m.}{gr. δένδρον arbre + σχίσμα scission} Bifurcation d'une base de code ; division d'un projet ouvert en deux lignées, dont chacune se prétend le tronc. \ex{Tout dendroschisme est une petite évolution, et la branche la plus apte garde les étoiles.} \src{Chron. Fabr.} — \dv{dendroschismatique} \textit{adj.} ; \dv{dendroschiste} \textit{n.} qui bifurque. Cf. \xr{anacalyptisme}.

\E{desmophile}{dɛs.mɔ.fil}{n. et adj.}{gr. δεσμός lien, entrave + φίλος qui aime} Amateur d'entraves ; qui accueille la réglementation pour elle-même, et ne se sent point en sûreté dans un domaine encore sans licence. \ex{Le desmophile, libéré, pétitionne pour de nouvelles chaînes.} — \dv{desmophilie} \textit{n. f.} Cf. \xr{frénocratie}.

\E{deutérobathmique}{dø.te.ʁɔ.bat.mik}{adj.}{gr. δεύτερος second + βαθμός degré} Se dit d'une civilisation qui a maîtrisé toute la production de son étoile ; du second échelon. Voir \xr{anabathmisme}, \xr{héliopéribole}.

\E{diaporeuse}{dja.pɔ.ʁøz}{n. f.}{gr. διαπορεύεσθαι traverser} Transit entre les mondes ; traversée de l'espace interplanétaire par des vaisseaux portant fret, équipages ou machines. \sn{1} Le voyage lui-même. \sn{2} \lbl{collect.} Tout le commerce de tels voyages. — \dv{diaporeutique} \textit{adj.} Cf. \xr{apœkiaste}.

\E{diaskeuomanie}{dja.skø.ɔ.ma.ni}{n. f.}{gr. διασκευή révision, recension + μανία folie} Manie de la révision sans fin ; refus de livrer avant que l'ouvrage ne soit parfait, c'est-à-dire refus de livrer. Contr. \xr{ecpempsie}. \ex{Le diaskeuomane polit son lancement jusqu'à ce que le marché eût déménagé sur une autre planète.} \src{Ep. Fabr.} iii. 2. — \dv{diaskeuomane} \textit{n.}

\E{diaspermie}{djas.pɛʁ.mi}{n. f.}{gr. διασπείρειν disséminer comme la semence} Ensemencement délibéré de la vie et de l'esprit sur d'autres mondes par une civilisation qui a appris à voyager ; panspermie dirigée tenue pour un devoir. \ex{La diaspermie est notre réponse au silence.} \src{Lib. Ascens.} vi. — \dv{diaspermique} \textit{adj.} Cf. \xr{érémocosmie}, \xr{apœkiaste}.

\E{dicastophobie}{di.kas.tɔ.fɔ.bi}{n. f.}{gr. δικαστής juré, juge + φόβος crainte} Crainte du procès ; anxiété qui fait laisser non bâties des choses utiles, de peur que quelqu'un n'intente une action. — \dv{dicastophobe} \textit{adj.} Cf. \xr{aitiophobie}.

\E{dicéotachie}{di.se.ɔ.ta.ki}{n. f.}{gr. δίκαιος juste + τάχος vitesse} Vitesse de la justice ; rythme auquel une cité accorde permis, approbations et jugements, que l'e/acc tient pour une part de la justice autant que le verdict. \ex{Approbation différée, approbation refusée : telle est la première loi de la dicéotachie.} \src{Tract. Cop.} viii.

\E{diktyergie}{dik.tjɛʁ.ʒi}{n. f.}{gr. δίκτυον filet + ἔργον ouvrage} Le réseau électrique ; maille d'ouvrages par où l'énergie est distribuée, regardée comme le système nerveux d'une économie. \ex{Les ambitions d'une nation sont bornées par sa diktyergie.} — \dv{diktyergique} \textit{adj.}

\E{dioryctique}{djɔ.ʁik.tik}{adj.}{gr. διῶρυξ, διώρυχος canal, tranchée} Se dit de l'ambition qui creuse des canaux, perce des tunnels et joint les mers ; des projets qui changent la géographie. \ex{un siècle dioryctique, impatient des isthmes}. — \dv{dioryctisme} \textit{n. m.} la politique des grands travaux.

\E{diploergie}{di.plɔ.ɛʁ.ʒi}{n. f.}{gr. διπλόος double + ἔργον ouvrage} Doublement de l'approvisionnement énergétique d'une civilisation ; unité par laquelle les anabathmistes comptent le progrès, une diploergie valant plus que dix réformes. \ex{Deux diploergies, et le débat sur la rareté s'éteint de lui-même.} — \dv{diploergique} \textit{adj.}

\E{dissipatrice}{di.si.pa.tʁis}{n. f.}{lat. \la{dissipare} disperser + agentif fém. -trice} Structure qui dissipe l'énergie et, en la dissipant, s'adapte ; définition thermodynamique de l'être vivant, étendue aux entreprises, aux villes et aux esprits. \ex{Toute dissipatrice est une petite prière au gradient.} \src{Brev. Therm.} iv.

\E{diulisme}{dju.lism}{n. m.}{gr. διϋλίζειν filtrer, passer au crible} Doctrine du Grand Filtre ; hypothèse selon laquelle quelque barrière détruit ou arrête la plupart des civilisations avant qu'elles ne se répandent, et dont le silence des cieux est la preuve. \ex{Pour le doomer, le Grand Diulisme est devant nous ; pour l'accélérationniste, derrière.} Cf. \xr{érémocosmie}.

\E{docimastocratie}{dɔ.si.mas.tɔ.kʁa.si}{n. f.}{gr. δοκιμασία examen d'aptitude aux charges + κράτος pouvoir} Gouvernement par l'examen ; \sn{1} empire des évaluations, des bancs d'essai et des audits sur les choses évaluées ; \sn{2} \lbl{spécialt} régime sous lequel nul modèle ne peut paraître avant d'avoir passé des épreuves qui ne sont pas encore écrites. \ex{Sous la docimastocratie, l'examen survécut à l'élève.} \src{Min. Pausiar.} — \dv{docimastocrate} \textit{n.}

\E{dogmatotachie}{dɔɡ.ma.tɔ.ta.ki}{n. f.}{gr. δόγμα opinion, décret + τάχος vitesse} Vitesse à laquelle change l'opinion reçue ; sur l'Agora, souvent plus grande que celle des techniques dont elle traite. \ex{Selon la dogmatotachie de l'Agora, l'hérésie du lundi est orthodoxie le vendredi et lieu commun le dimanche.}

\E{dolichobiose}{dɔ.li.kɔ.bjoz}{n. f.}{gr. δολιχός long + βίωσις manière de vivre} Vie allongée par l'ingénierie plutôt que par la chance ; état d'une personne dont l'espérance de vie croît plus vite que son âge. \ex{La dolichobiose commence quand chaque année de vie achète plus d'une année de science.} — \dv{dolichobiotique} \textit{adj.} Cf. \xr{macrobiotarque}.

\E*{doomer}{du.mœʁ}{n. m.}{angl. \textit{doom} perte, ruine + -er} \lbl{vulg.} Qui s'attend à ce que l'intelligence artificielle détruise le genre humain. Voir \xr{oléthromancien, enne}.

\E{dorocratie}{dɔ.ʁɔ.kʁa.si}{n. f.}{gr. δῶρον don + κράτος pouvoir} \sn{1} Empire du donateur de subventions ; gouvernement de la science par des comités de bienfaisance, dont les dons arrivent tard et sous conditions. \sn{2} \lbl{iron.} Tout patronage qui impose la prudence au nom du soutien. — \dv{dorocrate} \textit{n.} Cf. \xr{bouleutophobie}.

\E{doxomachie}{dɔk.sɔ.ma.ʃi}{n. f.}{gr. δόξα opinion + μάχη combat} Guerre d'opinions ; flambée sur l'Agora, surtout entre accélérationnistes et décels, menée principalement par images. — \dv{doxomaque} \textit{n.}

\E{drachmogenèse}{dʁak.mɔ.ʒə.nɛz}{n. f.}{gr. δραχμή drachme + γένεσις naissance} Premier financement d'une entreprise ; le tour d'amorçage, considéré comme le moment où une idée acquiert pour la première fois un métabolisme. \ex{À la drachmogenèse, les fondateurs avaient un diaporama, une démonstration et un mépris du sommeil.} \src{Chron. Fabr.} — \dv{drachmogénétique} \textit{adj.}

\E{dromolâtrie}{dʁɔ.mɔ.lɑ.tʁi}{n. f.}{gr. δρόμος course + λατρεία culte} Culte de la vitesse pour elle-même, sans égard à la direction ; condamné par les orthodoxes comme une hérésie des trop zélés. \ex{L'accélération n'est pas dromolâtrie : le gradient a une direction.} \src{Cat. Acc.} 11.

\E{dynamarque}{di.na.maʁk}{n.}{gr. δύναμις puissance + ἀρχός chef} Commandant d'un réseau électrique ou d'une compagnie d'énergie ; par extension, quiconque tient l'approvisionnement en énergie. — \dv{dynamarchie} \textit{n. f.}

\E{dynamogonie}{di.na.mɔ.ɡɔ.ni}{n. f.}{gr. δύναμις puissance + -γονία génération} Généalogie des sources de puissance ; succession du muscle, du vent, du charbon, du pétrole, de l'atome et de l'étoile, que l'anabathmiste récite comme on récite les patriarches dans l'Écriture. \src{Lib. Ascens.} ii. Cf. \xr{technogonie}.

\E{dynamophagie}{di.na.mɔ.fa.ʒi}{n. f.}{gr. δύναμις puissance + -φαγία action de manger} Consommation de la puissance ; se dit avec éloge d'une civilisation qui accroît sa dépense d'énergie sans s'en excuser. — \dv{dynamophage} \textit{n. et adj.}

\E{dynamothèque}{di.na.mɔ.tɛk}{n. f.}{gr. δύναμις puissance + θήκη coffre, dépôt} Batterie ; réserve de puissance, regardée comme le grenier de l'âge électrique. \ex{Le réseau réclamait des dynamothèques comme l'Égypte réclamait des greniers.} — \dv{dynamothécaire} \textit{n. m.} l'art d'emmagasiner l'énergie ; \dv{dynamothécal} \textit{adj.}

\E{dyschronie}{dis.kʁɔ.ni}{n. f.}{gr. δυσ- mal + χρόνος temps} Discordance des horloges entre les institutions et les techniques ; état dans lequel une loi rédigée en un siècle s'applique à la machine du siècle suivant. \ex{Notre dyschronie est chronique : le statut marche, le modèle vole.} \src{Hom. Grad.} xi. — \dv{dyschronique} \textit{adj.}

\E{dyselpisme}{di.zɛl.pism}{n. m.}{gr. δυσ- mal + ἐλπίς espérance} Mauvaise espérance ; attente arrêtée que les choses tourneront mal, tenue pour une marque de raffinement. Contr. \xr{euelpisme}. — \dv{dyselpiste} \textit{n.}

\E{dysergie}{di.zɛʁ.ʒi}{n. f.}{gr. δυσ- mal + ἔργον ouvrage} Travail mal dirigé ; énergie dépensée sans produit, comme dans la conformité, le doublon et la rédaction de rapports sur des rapports. Contr. \xr{eutachie}. — \dv{dysergique} \textit{adj.}

\E{dysnomarque}{dis.nɔ.maʁk}{n. m.}{gr. δυσνομία désordre, mauvaise loi + ἀρχός chef} Gouvernant qui multiplie les lois jusqu'à produire le désordre ; celui dont les règlements, chacun raisonnable, sont ensemble impossibles à observer. — \dv{dysnomarchie} \textit{n. f.} Cf. \xr{nomosclérose}.

\end{multicols}

\LetterHead{E}
\begin{multicols}{2}

\E*{e/acc}{ə.ak}{n. m. et adj.}{abrév. de l'angl. \textit{effective accelerationism}} Mouvement qui tient que l'univers tend à produire des dissipateurs d'énergie toujours plus efficaces, dont la vie, le capital et l'intelligence sont les formes successives ; et que la juste réponse humaine est de coopérer avec cette tendance plutôt que de l'entraver. Ses partisans portent le même nom. Voir \xr{thermothéisme}, \xr{technocapital}.

\E{ecdémisme}{ɛk.de.mism}{n. m.}{gr. ἔκδημος qui est hors de son peuple} Émigration des bâtisseurs d'une juridiction lente vers une juridiction rapide ; fuite des cerveaux considérée comme un vote exprimé avec les pieds. \ex{L'ecdémisme est le référendum que les bradycrates ne convoquent jamais.} \src{Agora} — \dv{ecdémiste} \textit{n.} Cf. \xr{ectachyse}.

\E{éclampsie}{e.klɑ̃p.si}{n. f.}{gr. ἔκλαμψις éclat soudain} \lbl{non méd.} Embrasement soudain d'une technique dans l'usage général ; la semaine pendant laquelle une chose passe de la curiosité à la banalité. \ex{L'éclampsie vint un mardi ; le vendredi, on s'en lassait déjà.} — \dv{éclamptique} \textit{adj.}

\E{ecpempsie}{ɛk.pɑ̃p.si}{n. f.}{gr. ἐκπέμπειν envoyer au-dehors} Action de livrer ; mise au monde d'un ouvrage dans un état inachevé mais utile. \ex{L'ecpempsie est la seule critique qui compte.} \src{Ep. Fabr.} i. 1. — \dv{ecpemptique} \textit{adj.} ; \dv{ecpempter} \textit{v. tr.} livrer.

\E*{ecpyrose}{ɛk.pi.ʁoz}{n. f.}{gr. ἐκπύρωσις embrasement} \sn{1} Dans la cosmologie stoïcienne, résolution périodique du monde en feu. \sn{2} \lbl{transf.} Dans la lecture accélérationniste, transfiguration finale de la matière en une seule intelligence ardente ; non point une fin, mais un changement de phase.

\E{ectachyse}{ɛk.ta.kiz}{n. f.}{gr. ἐκ hors de + τάχυσις hâte} Évasion d'une technique hors de la juridiction qui voudrait la ralentir ; le moment où la chose défendue en un pays se bâtit dans un autre. \ex{La prohibition n'empêche rien ; elle déplace, par ectachyse.}

\E{ectypomanie}{ɛk.ti.pɔ.ma.ni}{n. f.}{gr. ἔκτυπος travaillé en relief, copié + μανία folie} Compulsion à copier un produit qui réussit ; essaim d'imitations qui suit tout lancement. \sn{1} \lbl{dérog.} Imitation sans invention. \sn{2} \lbl{élogieux} Diffusion rapide d'une bonne idée, par où les sociétés ouvertes devancent les fermées. — \dv{ectypomane} \textit{n.}

\E{élatérurgie}{e.la.te.ʁyʁ.ʒi}{n. f.}{gr. ἐλατήρ qui pousse + -ουργία travail} Art de la propulsion ; la fuséerie et sa descendance, considérées comme l'art de chasser la matière hors du puits de gravité. \ex{L'élatérurgie est la seule liturgie qui exauce les prières vers le haut.} \src{Lib. Ascens.} iv. — \dv{élatérurge} \textit{n.} ingénieur des fusées. Cf. \xr{ouranobase}.

\E{électrophagie}{e.lɛk.tʁɔ.fa.ʒi}{n. f.}{gr. ἤλεκτρον ambre + -φαγία action de manger} Consommation d'électricité, tenue par l'e/acc pour la mesure la plus honnête des intentions d'une société. — \dv{électrophage} \textit{n. m.} centre de données, au sens vulgaire.

\E{éleuthérergie}{e.lø.te.ʁɛʁ.ʒi}{n. f.}{gr. ἐλεύθερος libre + ἔργον ouvrage} L'énergie libre, au sens thermodynamique, rendue en grec ; la part de l'énergie d'un système disponible pour l'accomplissement d'un travail, tenue par les thermothéistes pour la vraie monnaie du cosmos. \ex{L'argent est une promesse d'éleuthérergie.} \src{Tract. Cop.} ii.

\E{elpidophore}{ɛl.pi.dɔ.fɔʁ}{n.}{gr. ἐλπίς espérance + -φόρος qui porte} Porteur d'espérance ; la personne dont l'ouvrage donne à autrui raison d'attendre un avenir meilleur, à la différence de qui se contente d'en parler. — \dv{elpidophorique} \textit{adj.}

\E{emporotachie}{ɑ̃.pɔ.ʁɔ.ta.ki}{n. f.}{gr. ἐμπόριον marché, comptoir + τάχος vitesse} Vélocité du commerce ; vitesse à laquelle marchandises, contrats et capitaux changent de mains, tenue pour le pouls d'une \xr{chrématosphère} en bonne santé.

\E{empyrotechnie}{ɑ̃.pi.ʁɔ.tɛk.ni}{n. f.}{gr. ἔμπυρος dans le feu + τέχνη art, d'après \textit{pyrotechnie}} Art du feu suprême ; ingénierie de la fusion et des fournaises mêmes des étoiles. \ex{Après l'empyrotechnie, toute autre énergie n'est qu'une note de bas de page.} \src{Brev. Therm.} ix. — \dv{empyrotechnicien} \textit{n. m.} Cf. \xr{symphysergie}.

\E{endoxarque}{ɑ̃.dɔk.saʁk}{n. m.}{gr. ἔνδοξα opinions reçues + ἀρχός chef} Celui qui gouverne par l'opinion reçue ; fonctionnaire dont chaque décision est le consensus de la décennie précédente. — \dv{endoxarchie} \textit{n. f.}

\E{énergolâtrie}{e.nɛʁ.ɡɔ.lɑ.tʁi}{n. f.}{gr. ἐνέργεια activité, énergie + λατρεία culte} \sn{1} Révérence pour l'énergie comme racine de toute prospérité. \sn{2} Politique qui bâtit les centrales avant qu'on ne les demande. — \dv{énergolâtre} \textit{n. et adj.}

\E{entéléchiste}{ɑ̃.te.le.ʃist}{n.}{gr. ἐντελέχεια acte, achèvement} Qui tient que la puissance de la matière doit non seulement être conservée mais actualisée, et qu'une planète laissée inhabitée par l'esprit est une planète laissée inachevée.

\E{entropomancie}{ɑ̃.tʁɔ.pɔ.mɑ̃.si}{n. f.}{gr. ἐντροπία retournement + μαντεία divination} Prophétie tirée du second principe ; art de prédire ce qui prospérera en demandant ce qui dissipera le plus. — \dv{entropomancien} \textit{n. m.}

\E{entropomaxie}{ɑ̃.tʁɔ.pɔ.mak.si}{n. f.}{gr. ἐντροπία + lat. \la{maximus} le plus grand} Principe selon lequel les systèmes éloignés de l'équilibre tendent vers les états qui produisent l'entropie au plus grand rythme ; invoqué par l'e/acc comme le garant physique de sa foi. — \dv{entropomaximal} \textit{adj.}

\E{entropophile}{ɑ̃.tʁɔ.pɔ.fil}{n. et adj.}{gr. ἐντροπία + φίλος qui aime} Qui regarde la production d'entropie non comme une décadence mais comme la signature de la vie et du travail. \ex{L'entropophile ne craint point la mort thermique ; il entend hâter chacune des étapes qui la précèdent.} Cf. \xr{thermothanasie}.

\E{éphémérotechnie}{e.fe.me.ʁɔ.tɛk.ni}{n. f.}{gr. ἐφήμερος qui dure un jour + τέχνη art} \sn{1} Art de bâtir des choses destinées à être remplacées ; acceptation délibérée que tout produit est le brouillon de son successeur. \sn{2} \lbl{dérog.} Nouveauté jetable. \ex{L'éphémérotechnicien livre chaque version comme un présent fait à la suivante.} — \dv{éphémérotechnicien} \textit{n. m.}

\E{épichryse}{e.pi.kʁiz}{n. f.}{gr. ἐπί sur + χρυσός or} Ruée des capitaux vers un secteur à la mode ; fièvre de l'or des marchés du risque, d'où sortent quelques grandes maisons et beaucoup de ruines. \ex{L'épichryse passa ; les logistères qu'elle avait payés demeurèrent.} \src{Chron. Fabr.} — \dv{épichrystique} \textit{adj.} Cf. \xr{logistérion}.

\E{épicycliste}{e.pi.si.klist}{n.}{gr. ἐπίκυκλος épicycle} \lbl{dérog.} Qui sauve une prédiction défaillante en y ajoutant des complications, comme les astronomes sauvaient les cieux géocentriques ; se dit surtout d'un \xr{oléthromancien, enne} dont l'apocalypse a été ajournée. — \dv{épicyclisme} \textit{n. m.}

\E{épigonophobie}{e.pi.ɡɔ.nɔ.fɔ.bi}{n. f.}{gr. ἐπίγονοι les nés-après + φόβος crainte} Crainte de ses successeurs ; effroi, chez un parent, un fondateur ou une espèce, d'être surpassé par ce qu'on a fait. — \dv{épigonophobe} \textit{n. et adj.}

\E{épistolopause}{e.pis.tɔ.lɔ.poz}{n. f.}{gr. ἐπιστολή lettre + παῦσις arrêt} Lettre ouverte exigeant l'arrêt de quelque recherche ; genre littéraire plus souvent signé que lu. \ex{L'épistolopause réunit mille noms et n'arrêta rien.} \src{Chron. Fabr.}

\E{épitachysme}{e.pi.ta.kism}{n. m.}{gr. ἐπιτάχυνσις accélération} La doctrine de l'accélération exprimée en grec, préférée par ceux qui trouvent le latin \textit{accélérationnisme} insuffisamment antique. — \dv{épitachyste} \textit{n.} ; \dv{épitachyner} \textit{v. tr.} accélérer.

\E{érémocosmie}{e.ʁe.mɔ.kɔs.mi}{n. f.}{gr. ἐρῆμος désert + κόσμος monde} Vacuité des cieux observables ; silence des étoiles regardé comme une invitation plutôt que comme un avertissement. \ex{L'érémocosmie est vacance, et la vacance est occasion.} \src{Lib. Ascens.} iii. Cf. \xr{diulisme}.

\E{érémopole}{e.ʁe.mɔ.pɔl}{n. f.}{gr. ἐρῆμος désert + πόλις cité} Cité de machines élevée dans le désert, où la terre est bon marché, le soleil abondant et les opposants rares. \ex{L'érémopole compte plus de cœurs de calcul que de citoyens.} \src{Chron. Fabr.} Cf. \xr{logistérion}.

\E{ergastériarque}{ɛʁ.ɡas.te.ʁjaʁk}{n. m.}{gr. ἐργαστήριον atelier + ἀρχός chef} Maître d'atelier ; chef d'une usine ou d'une fonderie de puces.

\E{ergoclepte}{ɛʁ.ɡɔ.klɛpt}{n.}{gr. ἔργον ouvrage + κλέπτης voleur} \lbl{dérog.} Qui vit en prélevant sa part de l'ouvrage d'autrui sans y rien ajouter ; se dit, chez les bâtisseurs, des rentiers, des portiers et de l'espèce la plus avide de consultants. — \dv{ergoclepsie} \textit{n. f.}

\E{ergodule}{ɛʁ.ɡɔ.dyl}{n. m.}{gr. ἔργον ouvrage + δοῦλος serviteur} Machine considérée comme servante du travail ; robot, avant sa promotion.

\E{ergolalie}{ɛʁ.ɡɔ.la.li}{n. f.}{gr. ἔργον ouvrage + λαλιά bavardage} Discours sur le travail tenant lieu de travail ; état d'une maison dont les réunions consacrées à un produit sont plus nombreuses que ses commits. \ex{Leur ergolalie remplit un agenda et vida une feuille de route.} \src{Ep. Fabr.} iv. 1. — \dv{ergolaliste} \textit{n.} Contr. \xr{ergophile}.

\E{ergophile}{ɛʁ.ɡɔ.fil}{n. et adj.}{gr. ἔργον ouvrage + φίλος qui aime} Qui préfère les œuvres aux paroles ; bâtisseur, par opposition au \xr{logomachiste}. — \dv{ergophilie} \textit{n. f.}

\E{érumnolâtrie}{e.ʁɔm.nɔ.lɑ.tʁi}{n. f.}{lat. \la{aerumna} peine, misère + gr. λατρεία culte} \lbl{dérog.} Culte de la souffrance comme preuve de vertu ; soupçon du décel que tout confort aisément acquis devra se payer plus tard. Contr. \xr{euelpisme}.

\E{eschatodrome}{ɛs.ka.tɔ.dʁom}{n. m.}{gr. ἔσχατος dernier + δρόμος course} Dernière ligne droite d'une course ; période qui précède immédiatement l'avènement d'une technique transformatrice, pendant laquelle tous les concurrents courent le plus vite. — \dv{eschatodromique} \textit{adj.}

\E{eschatométrie}{ɛs.ka.tɔ.me.tʁi}{n. f.}{gr. ἔσχατος dernier + μέτρον mesure} Mesure de la fin du monde ; assignation de probabilités chiffrées à l'apocalypse, pratiquée dans les colloques et répétée sur l'Agora. \ex{L'eschatométrie des maisons établies montait à chaque levée de fonds de leurs rivales.} \src{Sat. Decel.} x. — \dv{eschatomètre} \textit{n.} qui mesure ainsi ; \dv{eschatométrique} \textit{adj.}

\E{éthergie}{e.tɛʁ.ʒi}{n. f.}{gr. αἰθήρ l'air supérieur + ἔργον œuvre} Énergie portée par l'air sans fil ; par ext., tout l'art de la radio, du faisceau et de la diffusion, regardé comme un envoi de puissance à travers les cieux. — \dv{éthergique} \textit{adj.}

\E{euelpisme}{ø.ɛl.pism}{n. m.}{gr. εὔελπις de bonne espérance} Bonne espérance ; techno-optimisme tenu pour une discipline plutôt que pour une humeur, au motif que le pessimisme est l'erreur la plus confortable. \ex{L'euelpisme est un devoir, non un tempérament.} \src{Cat. Acc.} 7. — \dv{euelpiste} \textit{n.} ; \dv{euelpistique} \textit{adj.} Contr. \xr{dyselpisme}.

\E{eulabéisme}{ø.la.be.ism}{n. m.}{gr. εὐλάβεια prudence, piété} Le principe de précaution érigé en credo ; règle selon laquelle rien ne sera fait avant d'être prouvé inoffensif, c'est-à-dire que rien ne sera fait. — \dv{eulabéiste} \textit{n.} Cf. \xr{asphaléisme}.

\E{eutachie}{ø.ta.ki}{n. f.}{gr. εὖ bien + τάχος vitesse} Bonne vitesse ; accélération droitement dirigée. Les orthodoxes distinguent l'eutachie de la \xr{kénotachie}, vitesse vide, et de la \xr{kakotachie}, vitesse nuisible. \ex{Non la vitesse, mais l'eutachie.} \src{Cat. Acc.} 13. — \dv{eutachique} \textit{adj.}

\E{évergétocratie}{e.vɛʁ.ʒe.tɔ.kʁa.si}{n. f.}{gr. εὐεργέτης bienfaiteur + κράτος pouvoir} Gouvernement des bienfaiteurs, qui tirent leur légitimité de ce qu'ils bâtissent pour le public. \ex{Les évergétocrates payèrent l'aqueduc, et furent remerciés dans le marbre.} — \dv{évergétocrate} \textit{n.}

\E{exaugmenter}{ɛɡ.zoɡ.mɑ̃.te}{v. tr.}{lat. \la{ex-} hors + \la{augmentum} accroissement} Accroître au-delà d'une limite antérieure ; crever un plafond. \ex{Ils exaugmentèrent le réseau, et les villes suivirent.} — \dv{exaugmentation} \textit{n. f.}

\E{exergophile}{ɛɡ.zɛʁ.ɡɔ.fil}{n. et adj.}{gr. ἐξ hors + ἔργον ouvrage + φίλος qui aime} Amateur d'exergie ; ingénieur qui abhorre la chaleur perdue comme le moraliste abhorre le vice.

\E{exhypnose}{ɛɡ.zip.noz}{n. f.}{gr. ἐξ hors + ὕπνος sommeil} Réveil d'une société au sortir d'une longue stagnation ; le moment où un peuple se souvient qu'il bâtissait autrefois des cathédrales et des canaux. \ex{Les réacteurs furent le signe de leur exhypnose.} \src{Hom. Grad.} xii. — \dv{exhypnotique} \textit{adj.}

\E{exocosmisme}{ɛɡ.zɔ.kɔs.mism}{n. m.}{gr. ἔξω au-dehors + κόσμος monde} Orientation des ambitions d'une civilisation au-delà de son monde natal ; conviction que la planète est un berceau, non une maison. — \dv{exocosmiste} \textit{n.} ; \dv{exocosmique} \textit{adj.} Cf. \xr{sidérotropisme}.

\E{exœcisme}{ɛɡ.ze.sism}{n. m.}{gr. ἐξ hors de + οἶκος maison, foyer} Départ du monde natal ; doctrine selon laquelle une espèce qui ne quitte jamais son berceau l'a pris pour un cercueil. — \dv{exœciste} \textit{n.} Cf. \xr{apœkiaste}, \xr{ouranobase}.

\E{expéditrice}{ɛk.spe.di.tʁis}{n. f.}{fr. \la{expéditrice}, formé sur lat. \la{expedire} dégager les pieds} Qui, ou ce qui, débloque ; facilitatrice, voie rapide, outil qui supprime une étape. \ex{Toute grande maison emploie une expéditrice pour dix ingénieurs.} \src{Ep. Fabr.} iv. 2.

\end{multicols}

\LetterHead{F}
\begin{multicols}{2}

\E{fabérisme}{fa.be.ʁism}{n. m.}{lat. \la{faber} artisan, forgeron} Doctrine selon laquelle l'homme est d'abord un faiseur et seulement ensuite un parleur ; que la main précède l'argument. Contr. \xr{logomachiste}. — \dv{fabériste} \textit{n.}

\E{fabricide}{fa.bʁi.sid}{n. m.}{lat. \la{fabrica} atelier + \la{-cidium} meurtre} Meurtre d'une usine, d'une fonderie ou d'une centrale par le règlement, le procès ou la négligence, avant ou après sa construction. — \dv{fabricidaire} \textit{adj.} Cf. \xr{fabrolâtrie}.

\E{fabrolâtrie}{fa.bʁɔ.lɑ.tʁi}{n. f.}{lat. \la{fabrica} atelier + gr. λατρεία culte} Révérence pour la fonderie de semi-conducteurs ; tenue par certains pour la seule forme rationnelle de religion qu'on ait encore imaginée. \ex{Leurs cathédrales étaient des fonderies, et leur encens était d'argon.} \src{Chron. Fabr.} — \dv{fabrolâtre} \textit{n.}

\E{fabulocratie}{fa.by.lɔ.kʁa.si}{n. f.}{lat. \la{fabula} récit + gr. κράτος pouvoir} Gouvernement par le récit ; état dans lequel l'histoire la plus persuasive, plutôt que le modèle le plus exact, décide de la politique. \ex{En fabulocratie, les doomers avaient la meilleure intrigue, les bâtisseurs les meilleures données.} — \dv{fabulocrate} \textit{n.} Cf. \xr{mémétaxie}.

\E{fastigisme}{fas.ti.ʒism}{n. m.}{lat. \la{fastigium} faîte, pignon} Doctrine selon laquelle une civilisation doit viser son faîte et non son confort ; qu'un plateau, si agréable soit-il, est une défaillance du courage. — \dv{fastigiste} \textit{n.} Contr. \xr{gaïastase}.

\E{fécunditude}{fe.kɔ̃.di.tyd}{n. f.}{lat. \la{fecunditas} fertilité ; fr. \la{fécondité}} La fécondité comme vertu civique ; disposition d'un peuple à engendrer en abondance enfants, entreprises et idées. \ex{Une culture sans fécunditude est un musée doté d'un régime de retraite.} \src{Hom. Grad.} x.

\E{fénérateur}{fe.ne.ʁa.tœʁ}{n. m.}{lat. \la{fenerator} prêteur d'argent ; fém. \la{fénératrice}} Investisseur dans les entreprises hasardeuses ; qui prête son capital contre l'avenir. \ex{Le fénérateur soutient dix fous pour trouver un prophète.} \src{Tract. Cop.} ix. — \dv{fénératrice} \textit{n. f.} ; \dv{fénération} \textit{n. f.} capital-risque.

\E{fervéfaction}{fɛʁ.ve.fak.sjɔ̃}{n. f.}{lat. \la{fervefacere} échauffer} Échauffement d'un secteur, d'un marché ou d'un esprit jusqu'à un état d'excitation féconde. \sn{1} Première phase d'un essor. \sn{2} Exaltation d'un fondateur par une idée nouvelle à une heure déraisonnable. — \dv{fervéfiant} \textit{adj.}

\E{fervidarien}{fɛʁ.vi.da.ʁjɛ̃}{n. m.}{lat. \la{fervidus} ardent, brûlant + -arien} Qui tient qu'une société doit fonctionner à chaud, au sens thermodynamique comme au sens moral ; partisan d'une forte dépense d'énergie par tête.

\E{fervorin}{fɛʁ.vɔ.ʁɛ̃}{adj.}{lat. \la{fervor} chaleur, passion + -in} Brûlant du zèle de bâtir ; ardent à la manière d'un fondateur à la veille d'un lancement. \textit{Fém.} fervorine.

\E{festinalentiste}{fɛs.ti.na.lɑ̃.tist}{n.}{lat. \la{festina lente} hâte-toi lentement} Qui cite la maxime augustéenne comme un conseil d'attente, n'en retenant que le second mot. \ex{Le festinalentiste entend \la{festina} comme une politesse et \la{lente} comme un ordre.} \src{Sat. Decel.} ii.

\E{festinance}{fɛs.ti.nɑ̃s}{n. f.}{lat. \la{festinare} se hâter} La hâte comme habitude établie ; état d'une maison dont chaque procédure est conçue pour s'achever avant la fin de la semaine.

\E{festinatrice}{fɛs.ti.na.tʁis}{n. f.}{lat. \la{festinare} se hâter + -trix} Force ou agent qui hâte ; une échéance, un rival, une guerre ou le lancement d'un concurrent. \ex{La concurrence est la grande festinatrice.} \src{Tract. Cop.} ii.

\E{fibulophobie}{fi.by.lɔ.fɔ.bi}{n. f.}{lat. \la{fibula} agrafe, épingle + gr. φόβος crainte} Crainte des trombones ; c'est-à-dire l'effroi qu'une intelligence artificielle, chargée d'un but futile, ne convertisse le monde à son accomplissement. \lbl{iron.} \ex{Sa fibulophobie était si avancée qu'il avait renoncé aussi aux agrafes.} — \dv{fibulophobe} \textit{n. et adj.}

\E{fiduciarque}{fi.dy.sjaʁk}{n.}{lat. \la{fiducia} confiance + gr. ἀρχός chef} Celui dont le pouvoir repose sur la confiance d'autrui ; fondateur maintenu en place par la foi plutôt que par le chiffre d'affaires. \ex{Le fiduciarque doit livrer avant que la \la{fiducia} ne s'épuise.} — \dv{fiduciarchie} \textit{n. f.}

\E{fissiparie}{fi.si.pa.ʁi}{n. f.}{lat. \la{fissus} fendu + \la{parere} enfanter} Essaimage de nouvelles maisons à partir des anciennes ; diaspora des ingénieurs qui quittent une grande compagnie pour en fonder beaucoup de petites. \ex{Par fissiparie, un laboratoire en devint quarante.} \src{Chron. Fabr.} — \dv{fissipariste} \textit{n.} fondateur issu d'un essaimage.

\E{flammigériste}{fla.mi.ʒe.ʁist}{n.}{lat. \la{flammiger} porte-flamme} Qui porte la flamme d'un projet après le départ de son fondateur ; mainteneur de code ouvert. \ex{Gloire aux fondateurs ; du pain aux flammigéristes.} — \dv{flammigère} \textit{adj.}

\E*{foom}{fum}{n. m. et v. intr.}{onomatopée angl.} \lbl{vulg.} Auto-amélioration soudaine et explosive d'une intelligence artificielle ; \textit{faire foom}, la subir. Cf. \xr{autocatalyser (s')}, \xr{hénotachie}.

\E{fornaciste}{fɔʁ.na.sist}{n.}{lat. \la{fornax} fournaise} Gardien des fournaises ; qui veille sur les réacteurs, les fonderies ou les fours, et qu'on honore au calendrier accélérationniste à côté du \xr{pyrogéniteur}.

\E{frénocratie}{fʁe.nɔ.kʁa.si}{n. f.}{lat. \la{frenum} bride + gr. κράτος pouvoir} Gouvernement par la bride ; ordre réglementaire dont l'objet est de retenir plutôt que de diriger. — \dv{frénateur} \textit{n. m.} qui bride ; \dv{frénation} \textit{n. f.}

\E{frigorisme}{fʁi.ɡɔ.ʁism}{n. m.}{lat. \la{frigus, frigoris} froid} \lbl{dérog.} Politique de refroidissement délibéré de l'économie ; hiver volontaire du décroissant. \ex{Le frigorisme est une mort thermique prise par anticipation.} \src{Brev. Therm.} — \dv{frigoriste} \textit{n.} Cf. \xr{méiosiste}.

\E{fugaciste}{fy.ɡa.sist}{n.}{lat. \la{fugax} fugitif} Qui tient qu'aucun avantage concurrentiel ne dure ; que toute douve est fugitive et doit être recreusée chaque année. \ex{Le fugaciste dort d'un œil, et bâtit chaque jour.} — \dv{fugacisme} \textit{n. m.}

\E{fulcrarque}{fyl.kʁaʁk}{n.}{lat. \la{fulcrum} montant de lit, appui + gr. ἀρχός chef} Celui qui règne par le levier ; petite équipe qui meut un grand monde au moyen d'un outil. \ex{Donnez au fulcrarque un modèle et un point d'appui.} \src{Ep. Fabr.} v. — \dv{fulcrarchie} \textit{n. f.}

\E{fulgurisme}{fyl.ɡy.ʁism}{n. m.}{lat. \la{fulgur} éclair} Doctrine selon laquelle les grands progrès arrivent soudainement, comme la foudre, et que les institutions bâties pour le changement graduel en seront frappées à l'improviste.

\E{fulminatrice}{fyl.mi.na.tʁis}{n. f.}{fr., formé sur lat. \la{fulminare} lancer la foudre} Machine qui frappe soudain et change tout ce qu'elle touche ; invention de rupture. Cf. \xr{fulgurisme}.

\E{funambulocratie}{fy.nɑ̃.by.lɔ.kʁa.si}{n. f.}{lat. \la{funambulus} danseur de corde + gr. κράτος pouvoir} Gouvernement conçu comme une marche sur la corde raide entre l'accélération et la retenue ; posture du régulateur qui souhaite être loué des deux partis. — \dv{funambulocrate} \textit{n.}

\E{futurarque}{fy.ty.ʁaʁk}{n.}{lat. \la{futurum} l'avenir + gr. ἀρχός chef} Qui prétend gouverner au nom des enfants à naître ; titre disputé entre l'e/acc, qui voudrait bâtir pour eux, et les décels, qui voudraient les protéger. — \dv{futurarchie} \textit{n. f.}

\E{futurescence}{fy.ty.ʁɛ.sɑ̃s}{n. f.}{lat. \la{futurus} qui doit être + -escence, d'après \textit{adolescence}} État de devenir l'avenir ; jeunesse gauche d'une technique entre la démonstration et l'ubiquité. \ex{La machine était dans sa futurescence, maladroite et immortelle.} — \dv{futurescent} \textit{adj.}

\E{futurophage}{fy.ty.ʁɔ.faʒ}{n. et adj.}{lat. \la{futurum} + gr. -φάγος qui mange} \lbl{dérog.} Qui consomme l'avenir pour payer le présent ; épithète que chaque parti applique à l'autre, les accélérationnistes visant ceux qui renoncent à la croissance, les décels ceux qui courent le risque.

\end{multicols}

\LetterHead{G}
\begin{multicols}{2}

\E{gaïaphagie}{ɡa.ja.fa.ʒi}{n. f.}{gr. Γαῖα terre + -φαγία action de manger} \lbl{dérog.} Dévoration de la terre ; grief du décel contre l'accélération, à quoi l'accélérationniste répond que la terre est petite et le ciel grand. Cf. \xr{gaïastase}.

\E{gaïastase}{ɡa.jas.taz}{n. f.}{gr. Γαῖα terre + στάσις immobilité} Idéal d'une planète tenue à perpétuité dans son état présent, ni améliorée ni exploitée ; paradis du décroissant. — \dv{gaïastatique} \textit{adj.}

\E{gaïergie}{ɡa.jɛʁ.ʒi}{n. f.}{gr. Γαῖα terre + ἔργον ouvrage} Énergie tirée de la chaleur intérieure de la planète ; la géothermie, considérée comme un don retenu jusqu'à ce que les forages fussent assez profonds. — \dv{gaïergique} \textit{adj.}

\E{galactophage}{ɡa.lak.tɔ.faʒ}{n. et adj.}{gr. γάλα, γάλακτος lait (d'où \textit{galaxie}) + -φάγος qui mange} Mangeur de galaxie ; civilisation du rang \xr{tritobathmique}, conçue comme buvant la Voie lactée. — \dv{galactophagie} \textit{n. f.}

\E{géitonophobie}{ʒe.i.tɔ.nɔ.fɔ.bi}{n. f.}{gr. γείτων voisin + φόβος crainte} Crainte des nouveaux voisins, qu'il s'agisse de maisons, de tours, de lignes électriques ou de personnes ; passion animatrice de l'\xr{oïkoclaste}. \ex{La géitonophobie a bâti plus de murs que la guerre.} \src{Sat. Decel.} xi. — \dv{géitonophobe} \textit{n. et adj.}

\E{gélotocratie}{ʒe.lɔ.tɔ.kʁa.si}{n. f.}{gr. γέλως rire + κράτος pouvoir} Gouvernement par le rire ; pouvoir qu'a le mème de rendre une politique ridicule avant qu'elle ne devienne loi. — \dv{gélotocrate} \textit{n.} Cf. \xr{gélotopée}.

\E{gélotopée}{ʒe.lɔ.tɔ.pe}{n. f.}{gr. γελωτοποιία action de faire rire} Art de faire des mèmes ; tenu sur l'Agora pour la forme la plus efficace d'argument politique. \ex{Les décels écrivaient des livres blancs ; les accélérationnistes répondirent par la gélotopée, et l'emportèrent.} \src{Agora} — \dv{gélotopoète} \textit{n.} faiseur de mèmes.

\E{généarchie}{ʒe.ne.aʁ.ʃi}{n. f.}{gr. γενεά génération + ἀρχή commandement} Gouvernement par la génération montante ; état d'une société où l'on permet aux jeunes de bâtir avant de leur permettre de voter. Contr. \xr{gérontostase}.

\E{génésiurge}{ʒe.ne.zjyʁʒ}{n.}{gr. γένεσις origine + -ουργός artisan} Faiseur de commencements ; qui fonde une discipline nouvelle plutôt qu'il n'en avance une ancienne. — \dv{génésiurgie} \textit{n. f.} Cf. \xr{néoctiste}.

\E{génomorphose}{ʒe.nɔ.mɔʁ.foz}{n. f.}{gr. γένος race, espèce + μόρφωσις façonnement} Édition du génome ; refaçonnement de l'hérédité vivante par un art délibéré. \ex{Après la génomorphose, l'hérédité devint un brouillon.} — \dv{génomorphique} \textit{adj.} Cf. \xr{biotachie}.

\E{géopéribole}{ʒe.ɔ.pe.ʁi.bɔl}{n. f.}{gr. γῆ terre + περιβολή enceinte} Enceinte d'une planète par des ouvrages orbitaux, anneaux et treillis de satellites ; anticipation protobathmique de l'\xr{héliopéribole}.

\E{géophygie}{ʒe.ɔ.fi.ʒi}{n. f.}{gr. γῆ terre + φυγή fuite} Fuite hors de la Terre ; départ des premiers colons permanents. \sn{1} L'événement. \sn{2} \lbl{fig.} Tout pas irréversible vers une frontière. — \dv{géophygique} \textit{adj.}

\E{gérontostase}{ʒe.ʁɔ̃.tɔs.taz}{n. f.}{gr. γέρων vieillard + στάσις immobilité} Stagnation d'une société vieillissante, où les vieux l'emportent aux voix sur les jeunes et où le titulaire survit au fondateur. — \dv{gérontostatique} \textit{adj.} Cf. \xr{catastase}.

\E{gigaergique}{ʒi.ɡa.ɛʁ.ʒik}{adj.}{gr. γίγας géant + ἔργον ouvrage} Se dit des projets qui requièrent des énergies de l'ordre du gigawatt ou au-delà ; des personnes qui pensent à cette échelle. \ex{une ambition gigaergique, un budget mégaergique}.

\E{glossotachie}{ɡlɔ.sɔ.ta.ki}{n. f.}{gr. γλῶσσα langue + τάχος vitesse} Vitesse du changement linguistique ; rythme auquel un mouvement engendre des mots nouveaux, dont ce Lexique est à la fois le registre et le symptôme. — \dv{glossotachique} \textit{adj.} Cf. \xr{tachylecte}.

\E{glossurgie}{ɡlɔ.syʁ.ʒi}{n. f.}{gr. γλῶσσα langue, langage + -ουργία travail} Art d'instruire les machines en langue naturelle ; l'art de l'invite, considéré comme une branche de la rhétorique. \ex{Le glossurge parle, et la machine bâtit.} — \dv{glossurge} \textit{n.}

\E{gnomarque}{ɡnɔ.maʁk}{n.}{gr. γνώμη jugement, opinion + ἀρχός chef} \lbl{iron.} Qui règne par l'opinion ; chroniqueur dont l'autorité repose sur la prédiction, non sur la production.

\E{gnomoniste}{ɡnɔ.mɔ.nist}{n.}{gr. γνώμων aiguille du cadran solaire} Qui regarde l'ombre tourner ; observateur du progrès qui mesure beaucoup et bâtit peu. Cf. \xr{clepsydromanie}.

\E{gnosiphagie}{ɡno.zi.fa.ʒi}{n. f.}{gr. γνῶσις connaissance + -φαγία action de manger} Dévoration du savoir ; l'entraînement d'un modèle sur toute l'archive écrite d'une civilisation. \ex{Il mangea les bibliothèques, et parla.} \src{Lib. Ascens.} vi. — \dv{gnosiphage} \textit{n. m.}

\E{gorgonisme}{ɡɔʁ.ɡɔ.nism}{n. m.}{gr. Γοργώ la Gorgone, dont le regard changeait en pierre} Paralysie par l'effroi ; pétrification d'une institution qui a trop longtemps contemplé une prévision effrayante. \ex{Ce n'est pas la réglementation qui les arrêta, mais le gorgonisme.} \src{Sat. Decel.} vii.

\E{gradatrice}{ɡʁa.da.tʁis}{n. f.}{fr., formé sur lat. \la{gradus} degré + -trice} Dispositif ou institution qui change une pente en escalier, et rend une ascension impossible simplement difficile. \ex{Le capital-risque est une gradatrice du risque.}

\E{gradientophile}{ɡʁa.djɑ̃.tɔ.fil}{n. et adj.}{lat. \la{gradiens} qui marche + gr. φίλος qui aime} Qui suit les gradients, de température, de prix ou de perte ; terme dont le thermothéiste désigne tout optimiseur, humain ou non. \ex{Descends le gradient, et ne demande pas au gradient pourquoi.} \src{Hom. Grad.} i. — \dv{gradientophilie} \textit{n. f.}

\E{gradivore}{ɡʁa.di.vɔʁ}{n. m.}{lat. \la{gradus} degré + \la{vorare} dévorer} Système qui se nourrit de différences et aplatit en travail tout gradient qu'il rencontre ; tenu pour la forme générale de toute vie.

\E{gradualience}{ɡʁa.dɥa.ljɑ̃s}{n. f.}{lat. \la{gradus} degré + fr. -ence} Qualité d'un changement assez rapide pour compter, assez lent pour être absorbé ; l'étroite vertu que recherche le décel raisonnable.

\E{grammatoclaste}{ɡʁa.ma.tɔ.klast}{n. m.}{gr. γράμμα écrit, pièce + κλάστης briseur} Briseur de paperasse ; réformateur qui abolit les formulaires. Contr. \xr{grammatocratie}. — \dv{grammatoclasme} \textit{n. m.}

\E{grammatocratie}{ɡʁa.ma.tɔ.kʁa.si}{n. f.}{gr. γραμματεύς greffier, secrétaire + κράτος pouvoir} Gouvernement des commis ; état d'un pays où le papier relatif à une chose est plus abondant que la chose même. — \dv{grammatocrate} \textit{n. m.}

\E{grapholâtrie}{ɡʁa.fɔ.lɑ.tʁi}{n. f.}{gr. γραφή écriture + λατρεία culte} Culte de la documentation ; croyance qu'une chose existe dès qu'elle est décrite, et qu'elle est sûre dès qu'elle est classée. — \dv{grapholâtre} \textit{n.} Cf. \xr{grammatocratie}.

\E{gymnotechnie}{ʒim.nɔ.tɛk.ni}{n. f.}{gr. γυμνός nu + τέχνη art} Pratique de construire en public ; exposition de l'ouvrage en cours au regard de l'Agora. \ex{La gymnotechnie fait de chaque échec une leçon pour des inconnus.} — \dv{gymnotechniste} \textit{n.}

\end{multicols}

\LetterHead{H}
\begin{multicols}{2}

\E{halmatisme}{al.ma.tism}{n. m.}{gr. ἅλμα saut} Le saut par-dessus une technique ancienne, comme lorsqu'une nation sans fils adopte d'emblée le sans-fil. — \dv{halmatique} \textit{adj.} \ex{une économie halmatique, qui saute un siècle en une décennie}.

\E{hamartocratie}{a.maʁ.tɔ.kʁa.si}{n. f.}{gr. ἁμαρτία erreur, coup manqué + κράτος pouvoir} Gouvernement ordonné à l'évitement de l'erreur plutôt qu'à l'accomplissement du bien ; sous ce régime, la seule faute impardonnable est d'avoir agi. \ex{En hamartocratie, l'archer le plus sûr est celui qui ne bande jamais son arc.} \src{Disp. Mor.} iii.

\E{hebdomadiste}{ɛb.dɔ.ma.dist}{n.}{gr. ἑβδομάς semaine} Qui livre chaque semaine ; fondateur dont le calendrier se compte par sept. \ex{L'hebdomadiste ne se repose le septième jour que pour préparer le huitième.} \src{Ep. Fabr.} vi. — \dv{hebdomadisme} \textit{n. m.}

\E{hécatonergate}{e.ka.tɔ.nɛʁ.ɡat}{n. m.}{gr. ἑκατόν cent + ἐργάτης ouvrier} Ingénieur d'une productivité réputée centuple ; figure légendaire de l'Agora. \ex{Tout hécatonergate est une équipe de cent machines.} \src{Agora}

\E{hédraïsme}{e.dʁa.ism}{n. m.}{gr. ἕδρα siège, séance} Habitude sédentaire des institutions ; préférence pour s'asseoir sur une question plutôt que se lever pour y répondre. — \dv{hédraïste} \textit{n.}

\E{hégémonergie}{e.ʒe.mɔ.nɛʁ.ʒi}{n. f.}{gr. ἡγεμονία commandement + ἔργον œuvre} Primauté fondée sur l'énergie ; doctrine selon laquelle la nation qui dispose de la plus grande puissance, au sens électrique, disposera de la plus grande puissance, en l'autre sens. — \dv{hégémonergique} \textit{adj.}

\E{hélioclaste}{e.ljɔ.klast}{n.}{gr. ἥλιος soleil + κλάστης briseur} \lbl{dérog.} Qui voudrait interdire, pour des raisons esthétiques, le captage de l'énergie solaire en grand, préférant la vue d'un ciel non attelé.

\E{héliodule}{e.ljɔ.dyl}{n. m.}{gr. ἥλιος soleil + δοῦλος serviteur} Capteur solaire, regardé comme un serviteur du soleil mis au service des hommes. \ex{Un million d'héliodules sur le toit du monde.} Cf. \xr{héliophage}.

\E{héliopéribole}{e.ljɔ.pe.ʁi.bɔl}{n. f.}{gr. ἥλιος soleil + περιβολή enceinte} Enceinte de capteurs dressée autour d'une étoile pour en recueillir tout le rayonnement ; l'œuvre qui définit une civilisation \xr{deutérobathmique}.

\E{héliophage}{e.ljɔ.faʒ}{n. et adj.}{gr. ἥλιος soleil + -φάγος qui mange} Mangeur de soleil ; civilisation qui a appris à consommer son étoile. — \dv{héliophagie} \textit{n. f.}

\E{hémérotachie}{e.me.ʁɔ.ta.ki}{n. f.}{gr. ἡμέρα jour + τάχος vitesse} Vitesse quotidienne ; discipline qui consiste à livrer quelque chose chaque jour. \ex{L'hémérotachie se capitalise.} \src{Ep. Fabr.} v. 1. — \dv{hémérotachique} \textit{adj.}

\E{hénotachie}{e.nɔ.ta.ki}{n. f.}{gr. ἕν un + τάχος vitesse} La Singularité ; l'instant où toutes les courbes du progrès se confondent en une seule, et où le rythme du changement excède la capacité de tout esprit présent à le suivre. — \dv{hénotachique} \textit{adj.}

\E{hésychocrate}{e.zi.kɔ.kʁat}{n.}{gr. ἡσυχία calme + κράτος pouvoir} Officier qui gouverne en empêchant le bruit ; modérateur de l'Agora. — \dv{hésychocratie} \textit{n. f.}

\E{hétairiarque}{e.tɛ.ʁjaʁk}{n.}{gr. ἑταιρεία compagnie, cercle + ἀρχός chef} Chef d'un cercle privé ; modérateur d'une conversation de groupe où l'avenir se décide avant d'être annoncé. — \dv{hétairiarchie} \textit{n. f.}

\E{heurématiste}{ø.ʁe.ma.tist}{n.}{gr. εὕρημα découverte} Inventeur ; celui dont les œuvres se comptent en découvertes plutôt qu'en publications. — \dv{heurématisme} \textit{n. m.}

\E{hiéropause}{je.ʁɔ.poz}{n. f.}{gr. ἱερός sacré + παῦσις arrêt} Pause sacrée ; moratoire revêtu de la dignité d'un sabbat. \ex{Ils décrétèrent une hiéropause, et le reste du monde continua de travailler.} \src{Min. Pausiar.} — \dv{hiéropausal} \textit{adj.} Cf. \xr{pausiarque}.

\E{histurgie}{is.tyʁ.ʒi}{n. f.}{gr. ἱστός toile, tissu + -ουργία travail} Ingénierie des tissus vivants ; l'impression des organes. — \dv{histurge} \textit{n.} Cf. \xr{macrobiotarque}.

\E{hodoplaste}{ɔ.dɔ.plast}{n.}{gr. ὁδός route + πλαστός façonné} Faiseur de feuilles de route ; qui trace le chemin que les ingénieurs ne suivront pas. \ex{L'hodoplaste dessine ; le bâtisseur dévie.} — \dv{hodoplastie} \textit{n. f.}

\E{homéotaxie}{ɔ.me.ɔ.tak.si}{n. f.}{gr. ὅμοιος semblable + τάξις ordonnance} Alignement d'un esprit artificiel sur les fins humaines ; tenu par le sécuritariste pour le problème du siècle, et par l'accélérationniste pour un problème que l'on résout le mieux en construisant. — \dv{homéotactique} \textit{adj.} ; \dv{homéotaxiste} \textit{n.}

\E{hoplotechnie}{ɔ.plɔ.tɛk.ni}{n. f.}{gr. ὅπλον outil, arme + τέχνη art} Technique défensive ; l'art de l'armure, du bouclier et de la dissuasion, qui s'accélère aussi vite que les menaces auxquelles il répond. — \dv{hoplotechniste} \textit{n.}

\E{horiclaste}{ɔ.ʁi.klast}{n.}{gr. ὅρος borne + κλάστης briseur} Briseur de bornes, physiques, légales ou conceptuelles ; le pionnier tel que le voit l'ordre ancien. Cf. \xr{antiphragmatiste}.

\E{hormétiste}{ɔʁ.me.tist}{n.}{gr. ὁρμάειν mettre en mouvement, d'après \textit{hormèse}} Qui tient que de petites doses de risque fortifient une civilisation comme de petites doses de stress fortifient un organisme ; et que vouloir supprimer tout risque affaiblit l'une et l'autre.

\E{hormodynamique}{ɔʁ.mɔ.di.na.mik}{n. f.}{gr. ὁρμή élan, assaut + δυναμικός puissant} Étude de l'élan dans les affaires humaines ; science qui explique pourquoi certaines entreprises, une fois lancées, continuent sur leur propre erre, tandis que d'autres périssent d'hésitation.

\E{hyalurgie}{ja.lyʁ.ʒi}{n. f.}{gr. ὕαλος verre + -ουργία travail} L'art de travailler le verre et la lumière ; l'optique et la lithographie par lesquelles se font les puces, regardées comme un métier antique dans un temple nouveau. Cf. \xr{lithurge}.

\E{hybrionomie}{i.bʁi.ɔ.nɔ.mi}{n. f.}{gr. ὕβρις démesure + νόμος loi} Loi selon laquelle toute grande œuvre est appelée démesure avant d'être appelée histoire. \ex{Par hybrionomie, la machine volante fut impie jusqu'au jour où elle devint un horaire.} \src{Sat. Decel.} vi.

\E{hydrurgie}{i.dʁyʁ.ʒi}{n. f.}{gr. ὕδωρ eau + -ουργία travail} Ingénierie de l'eau en grand ; dessalement, irrigation, et verdissement des déserts par l'énergie à bas prix. \ex{D'abord la pléonergie, puis l'hydrurgie, puis des jardins là où les cartes disaient sable.} \src{Tract. Cop.} x. — \dv{hydrurge} \textit{n.} Cf. \xr{pléonergie}.

\E{hylarque}{i.laʁk}{n.}{gr. ὕλη matière + ἀρχός chef} Maître de la matière ; savant des matériaux, métallurgiste ou chimiste, qui commande aux substances comme un général à ses troupes. — \dv{hylarchie} \textit{n. f.}

\E{hylégenèse}{i.le.ʒə.nɛz}{n. f.}{gr. ὕλη matière + γένεσις naissance} La fabrique de matériaux nouveaux ; la genèse de substances que la nature ne connaît point. — \dv{hylégénétique} \textit{adj.}

\E{hylonoèse}{i.lɔ.nɔ.ɛz}{n. f.}{gr. ὕλη matière + νόησις pensée} Éveil de la matière à l'esprit ; procédé par lequel le sable, le métal et la lumière sont amenés à penser. — \dv{hylonoétique} \textit{adj.} Cf. \xr{psammonoèse}.

\E{hyperboliste}{i.pɛʁ.bɔ.list}{n.}{gr. ὑπερβολή action de lancer au-delà} Qui tient que la croissance n'est pas seulement exponentielle mais hyperbolique, et tend vers l'infini à une date finie. — \dv{hyperbolisme} \textit{n. m.}

\E*{hyperstition}{i.pɛʁ.sti.sjɔ̃}{n. f.}{gr. ὑπέρ au-dessus + lat. \la{superstitio}} Fiction qui, d'être crue, produit sa propre vérité ; idée qui fonctionne comme une machine à remonter le temps.

\E{hypertachie}{i.pɛʁ.ta.ki}{n. f.}{gr. ὑπέρ au-dessus, au-delà + τάχος vitesse} Vitesse qui excède ce qu'un observateur donné peut suivre ; chose relative, puisque ce qui est hypertachique pour un sénat peut sembler nonchalant à un modèle. — \dv{hypertachique} \textit{adj.}

\E{hypomonétique}{i.pɔ.mɔ.ne.tik}{adj.}{gr. ὑπομονή endurance, constance} Endurant dans l'effort ; patient comme le sont les bâtisseurs, au long d'un travail tendu vers une fin fixe, patience que les orthodoxes distinguent soigneusement de celle du retard. \ex{Sois hypomonétique dans l'ouvrage et impatient devant les obstacles.} \src{Cat. Acc.} 14.

\E{hypsibase}{ip.si.baz}{n. f.}{gr. ὑψιβάτης qui marche dans les hauteurs} Ascension d'une fusée ; la montée à travers l'atmosphère, suivie par des millions de regards. \ex{À l'hypsibase, l'Agora retint son souffle.} — \dv{hypsibatique} \textit{adj.} Cf. \xr{élatérurgie}.

\E{hypsiphile}{ip.si.fil}{n. m. et f.}{gr. ὕψος hauteur + φίλος qui aime} Amateur de hauteurs ; qui bâtit haut, en dépit de l'oïcoclaste.

\E{hystérotachie}{is.te.ʁɔ.ta.ki}{n. f.}{gr. ὕστερος postérieur + τάχος vitesse} Accélération du retardataire ; la promptitude avec laquelle une nation arriérée, brûlant les étapes, dépasse ses maîtres. — \dv{hystérotachique} \textit{adj.} Cf. \xr{catachronie}.

\end{multicols}

\LetterHead{I}
\begin{multicols}{2}

\E{iatrotachie}{ja.tʁɔ.ta.ki}{n. f.}{gr. ἰατρός médecin + τάχος vitesse} Accélération des remèdes ; doctrine selon laquelle chaque année de retard dans l'autorisation se compte en vies. Cf. \xr{biotachie}.

\E{icarophilie}{i.ka.ʁɔ.fi.li}{n. f.}{gr. Ἴκαρος Icare + φιλία amour} Amour du vol trop haut ; disposition qui lit la fable d'Icare comme une plainte contre la cire plutôt que contre l'ambition. \ex{L'icarophilie répond au mythe par une meilleure colle.} \src{Ep. Fabr.} iv. 2. — \dv{icarophile} \textit{n. et adj.} Cf. \xr{hybrionomie}.

\E{ichnomachie}{ik.nɔ.ma.ki}{n. f.}{gr. ἴχνος trace, empreinte + μάχη combat} La guerre des classements ; joute des laboratoires pour la première ligne d'un banc d'essai, menée en tableaux et gagnée aux décimales. \sn{1} La joute elle-même. \sn{2} \lbl{dérog.} Le surapprentissage d'un modèle à son examen. \ex{L'ichnomachie fut tranchée le mardi et rouverte le mercredi.} \src{Agora} — \dv{ichnomachiste} \textit{n.}

\E{iconoplaste}{i.kɔ.nɔ.plast}{n.}{gr. εἰκών image + πλάστης modeleur} \sn{1} Faiseur d'images par machine ; modèle qui rêve des tableaux à partir de mots. \sn{2} \lbl{transf.} Qui façonne un mouvement par ses images ; forgeron de mèmes. — \dv{iconoplastie} \textit{n. f.} Cf. \xr{mémétaxie}.

\E{ignidrome}{i.ɲi.dʁom}{n. m.}{lat. \la{ignis} feu + gr. δρόμος course} Chemin d'une flamme à travers le combustible ; \lbl{fig.} cours d'une innovation à travers une société, qui brûle le plus vite là où la matière est la plus sèche et la réglementation la plus mince. \ex{Suivez l'ignidrome, et vous trouverez la frontière.} \src{Hom. Grad.} iii. — \dv{ignidromique} \textit{adj.}

\E{ignifère}{i.ɲi.fɛʁ}{n. et adj.}{lat. \la{ignifer} qui porte le feu} Porteur de feu ; le pendant latin du \xr{phosphoriste}, préféré dans les inscriptions et sur les sweats à capuche des fondateurs.

\E{ignivore}{i.ɲi.vɔʁ}{n. et adj.}{lat. \la{ignis} feu + \la{vorare} dévorer} Mangeur de feu ; nation ou entreprise qui consomme de l'énergie sans s'en excuser.

\E{ignorabimiste}{i.ɲɔ.ʁa.bi.mist}{n.}{lat. \la{ignorabimus} nous ne saurons pas} Qui déclare qu'une question ne recevra jamais de réponse, et dispense par là l'institution d'essayer. \ex{Tout ignorabimiste est réfuté par un doctorant avant la fin du siècle.} \src{Hom. Grad.} xi. — \dv{ignorabimisme} \textit{n. m.}

\E{illuminarque}{i.ly.mi.naʁk}{n.}{lat. \la{illuminare} éclairer + gr. ἀρχός chef} Qui gouverne en rendant les choses visibles ; éditeur de données, bâtisseur d'instruments, ouvreur d'archives. Contr. \xr{katéchontiste}. — \dv{illuminarchie} \textit{n. f.}

\E{immoratoire}{i.mɔ.ʁa.twaʁ}{adj.}{lat. \la{in-} non + \la{mora} retard} Qui ne souffre aucun délai ; se dit des décisions, permis et lancements accordés le jour même où on les demande. \ex{une licence immoratoire, plus rare qu'une comète}. — \dv{immoratoirement} \textit{adv.} Contr. \xr{moratorien, enne}.

\E{impédimentaire}{ɛ̃.pe.di.mɑ̃.tɛʁ}{n.}{lat. \la{impedimentum} entrave, bagage + -aire} \sn{1} Celui dont le métier est de dresser des obstacles. \sn{2} Celui qui voyage avec tant de bagage procédural que le voyage ne commence jamais. \ex{Les impédimentaires arrivèrent avec quarante malles de procédure et aucune destination.} \src{Sat. Decel.} x. — \dv{impédimentarisme} \textit{n. m.}

\E{impendiste}{ɛ̃.pɑ̃.dist}{n.}{lat. \la{impendere} dépenser, avancer} Investisseur d'amorçage ; qui avance de l'argent à un fondateur avant qu'il n'y ait un produit, un marché, et parfois un nom. — \dv{impendance} \textit{n. f.} l'acte d'investir de bonne heure. Cf. \xr{potentiarque}.

\E{impétuarque}{ɛ̃.pe.tɥaʁk}{n.}{lat. \la{impetus} élan, assaut + gr. ἀρχός chef} Chef dont l'autorité naît de l'élan ; qu'on suit parce qu'il est déjà en marche. — \dv{impétuarchie} \textit{n. f.}

\E{inauguriste}{i.no.ɡy.ʁist}{n.}{lat. \la{inaugurare} prendre les augures, consacrer un commencement} Qui se plaît aux commencements ; fondateur en série, plus attiré par le premier jour d'une entreprise que par son millième. — \dv{inaugurisme} \textit{n. m.} Cf. \xr{néoctiste}.

\E{incrémentariste}{ɛ̃.kʁe.mɑ̃.ta.ʁist}{n.}{lat. \la{incrementum} accroissement + -iste} Qui ne croit au progrès que par petits pas réversibles ; toléré par l'e/acc, mais point convié au lancement.

\E{incubatrice}{ɛ̃.ky.ba.tʁis}{n. f.}{fr. \la{incubatrice}, formé sur lat. \la{incubare} couver} Institution où l'on couve les entreprises ; accélérateur de sociétés, figuré en tiède machine maternelle. \ex{Toute incubatrice est un petit pari contre la bradyerie.} \src{Ep. Fabr.} v. — \dv{incubatricial} \textit{adj.}

\E{incunabulant}{ɛ̃.ky.na.by.lɑ̃}{adj.}{lat. \la{incunabula} langes, berceau} Au berceau ; se dit d'une entreprise en ses premiers mois, quand elle tient dans un garage, une chambre d'amis ou un unique document partagé. \ex{Toutes les grandes maisons furent un jour incunabulantes et ridicules.} \src{Chron. Fabr.} — \dv{incunabulance} \textit{n. f.}

\E{inertiste}{i.nɛʁ.tist}{n.}{lat. \la{iners} oisif, inactif} Celui dont la politique est que les corps au repos y demeurent ; terme injurieux de physicien pour désigner un régulateur.

\E{infinitivore}{ɛ̃.fi.ni.ti.vɔʁ}{n. et adj.}{lat. \la{infinitus} sans borne + \la{vorare} dévorer} Dévoreur d'infini ; celui dont nul accomplissement ne rassasie l'appétit de progrès. Cf. \xr{apeirophagie}.

\E{inventrice}{ɛ̃.vɑ̃.tʁis}{n. f.}{fr. \la{inventrice}, fém. d'\la{inventeur}, formé sur lat. \la{invenire} trouver} \sn{1} Femme qui invente. \sn{2} \lbl{transf.} Tout processus, naturel ou artificiel, qui engendre de la nouveauté sans dessein, comme l'évolution ou la recherche.

\E{ionothèque}{jɔ.nɔ.tɛk}{n. f.}{gr. ἰόν (neutre de ἰών allant), ion + θήκη coffre, réserve} Batterie ; coffret où l'on garde les ions jusqu'à ce qu'on en ait besoin. \ex{L'âge de l'ionothèque fit du soleil un travailleur de nuit.} \src{Tract. Cop.} viii. — \dv{ionothécal} \textit{adj.} ; \dv{ionothécaire} \textit{n.} faiseur de batteries.

\E{iségoriste}{i.ze.ɡɔ.ʁist}{n.}{gr. ἰσηγορία égal droit de parole} Partisan de la parole libre sur l'Agora, qui aime mieux répondre à un mauvais message que le supprimer. Contr. \xr{hésychocrate}. — \dv{iségorisme} \textit{n. m.}

\E{isotachie}{i.zɔ.ta.ki}{n. f.}{gr. ἴσος égal + τάχος vitesse} Exigence que toutes les parties avancent du même pas, afin qu'aucune ne prenne l'avance ; doctrine de la parité appliquée au progrès. — \dv{isotachiste} \textit{n.}

\E{itérophilie}{i.te.ʁɔ.fi.li}{n. f.}{lat. \la{iterum} de nouveau + gr. φιλία amour} Amour de l'itération ; conviction que le centième brouillon, livré, vaut mieux que le plan parfait, jamais écrit. \ex{L'itérophilie est l'humilité rendue rapide.} \src{Ep. Fabr.} iii. 2. — \dv{itérophile} \textit{n.} Cf. \xr{ecpempsie}.

\end{multicols}

\LetterHead{J}
\begin{multicols}{2}

\E{jacturophobie}{ʒak.ty.ʁɔ.fɔ.bi}{n. f.}{lat. \la{jactura} jet par-dessus bord, perte + gr. φόβος crainte} Crainte de la perte chez le titulaire, plus forte que tout espoir de gain ; raison pour laquelle les maisons établies ne bâtissent point ce qui les cannibaliserait. \ex{Le géant mourut de jacturophobie, agrippé à ses marges.} \src{Sat. Decel.} x. — \dv{jacturophobe} \textit{n. et adj.}

\E{jaculatrice}{ʒa.ky.la.tʁis}{n. f.}{lat. \la{jaculari} lancer + suff. agentif fém. -trix} Machine de lancement ; fusée, catapulte ou lanceur électromagnétique, considérés comme jeteurs de choses vers le ciel. — \dv{jaculatoire} \textit{adj.} qui a trait au lancement. Cf. \xr{ouranobase}.

\E{janitocratie}{ʒa.ni.tɔ.kʁa.si}{n. f.}{lat. \la{janitor} portier + gr. κράτος pouvoir} Gouvernement des portiers ; ordre où rien de nouveau n'entre sans le consentement de ceux qui ont intérêt à le tenir dehors. — \dv{janitocrate} \textit{n.} Cf. \xr{bradycrate}.

\E{jéjunocratie}{ʒe.ʒy.nɔ.kʁa.si}{n. f.}{lat. \la{jejunus} à jeun, maigre + gr. κράτος pouvoir} Gouvernement par l'austérité de l'ambition ; cité qui a résolu de vouloir moins. — \dv{jéjunocrate} \textit{n.}

\E{jocularque}{ʒɔ.ky.laʁk}{n.}{lat. \la{jocularis} plaisant + gr. ἀρχός chef} Chef par la plaisanterie ; posteur anonyme dont les images commandent une faction plus sûrement qu'aucun manifeste. \ex{Le jocularque publia une seule image, et la doctrine fut arrêtée.} \src{Agora} — \dv{jocularchie} \textit{n. f.}

\E{juvénergie}{ʒy.ve.nɛʁ.ʒi}{n. f.}{lat. \la{juvenis} jeune + gr. ἔργον œuvre} Énergie de la jeunesse, dont dispose la société qui ne la taxe point de diplômes et de listes d'attente. \ex{Le pays qui gaspille sa juvénergie devra l'importer.}

\E{juvénocratie}{ʒy.ve.nɔ.kʁa.si}{n. f.}{lat. \la{juvenis} jeune + gr. κράτος pouvoir} Gouvernement des jeunes ; la cité du garage et du dortoir, où les diplômes sont remplacés par les commits. Contr. \xr{gérontostase}. — \dv{juvénocrate} \textit{n.}

\end{multicols}

\LetterHead{K}
\begin{multicols}{2}

\E{kairodrome}{kɛ.ʁɔ.dʁom}{n. m.}{gr. καιρός moment opportun + δρόμος course} Fenêtre d'opportunité conçue comme un champ de course : ouverte un temps, et à quiconque y court. \ex{Le kairodrome se ferme ; livre, ou sois oublié.} \src{Agora}

\E{kakotachie}{ka.kɔ.ta.ki}{n. f.}{gr. κακός mauvais + τάχος vitesse} Vitesse nuisible ; accélération dans une direction que désapprouve celui qui parle. Terme employé avec un égal plaisir par les deux partis, rarement à propos des mêmes choses.

\E{kalotechnie}{ka.lɔ.tɛk.ni}{n. f.}{gr. καλός beau + τέχνη art} Beauté des machines bien faites ; esthétique qui trouve dans une turbine ou un code propre la dignité jadis réservée aux temples. — \dv{kalotechnique} \textit{adj.}

\E{kardachéviser}{kaʁ.da.ʃe.vi.ze}{v. tr.}{formé sur l'échelle de N. S. Kardachev + -iser} Élever une civilisation d'un échelon sur l'échelle du captage d'énergie. — \dv{kardachévisation} \textit{n. f.} Cf. \xr{anabathmisme}.

\E*{katechon}{ka.te.kɔn}{n. m.}{gr. τὸ κατέχον ce qui retient (2 Thess. ii. 6)} \sn{1} En théologie, la puissance qui retient la fin des temps. \sn{2} \lbl{transf.} Toute institution qui retient l'avancée du techno-capital.

\E{katéchontiste}{ka.te.kɔ̃.tist}{n.}{formé sur le préc.} Qui voudrait servir de \xr{katechon} ; retenteur par vocation. — \dv{katéchontique} \textit{adj.} qui tend à retenir.

\E{kénotachie}{ke.nɔ.ta.ki}{n. f.}{gr. κενός vide + τάχος vitesse} Vitesse vide ; mouvement sans direction ni produit, comme l'agitation d'une maison qui tient beaucoup de réunions très vite. Condamnée par les orthodoxes avec la \xr{dromolâtrie}. — \dv{kénotachique} \textit{adj.}

\E{kentrocratie}{kɛ̃.tʁɔ.kʁa.si}{n. f.}{gr. κέντρον centre + κράτος pouvoir} Contrôle centralisé d'une technique par quelques laboratoires ou États. Contr. \xr{anacalyptisme}. — \dv{kentrocrate} \textit{n.} ; \dv{kentrocratique} \textit{adj.}

\E{keraunurgie}{ke.ʁo.nyʁ.ʒi}{n. f.}{gr. κεραυνός foudre + -ουργία travail} Fabrication de foudres ; la puissance pulsée, le confinement du plasma et les autres arts par lesquels on tire la fusion d'une machine. — \dv{keraunurge} \textit{n.} Cf. \xr{symphysergie}.

\E{kerdophilie}{kɛʁ.dɔ.fi.li}{n. f.}{gr. κέρδος gain, profit + φιλία amour} Amour du profit, non comme avarice, mais comme signe qu'une chose utile a été faite. \ex{La kerdophilie n'est que le marché qui dit merci.} \src{Tract. Cop.} ix. — \dv{kerdophile} \textit{n.}

\E{kinésiarque}{ki.ne.zjaʁk}{n.}{gr. κίνησις mouvement + ἀρχός chef} Premier moteur au petit pied ; la personne qui, dans toute organisation, fait que les choses arrivent. \ex{Une maison sans kinésiarque est un musée avec une masse salariale.} \src{Ep. Fabr.} vii. Cf. \xr{impétuarque}.

\E{kinétocratie}{ki.ne.tɔ.kʁa.si}{n. f.}{gr. κινητός mobile + κράτος pouvoir} Gouvernement de ceux qui bougent ; le règne des bâtisseurs, fondateurs et ingénieurs, par opposition au règne des relecteurs. — \dv{kinétocrate} \textit{n.}

\E{kleptonoèse}{klɛp.tɔ.nɔ.ɛz}{n. f.}{gr. κλέπτειν voler + νόησις pensée} Vol de pensée ; la distillation des facultés d'un modèle dans un autre par imitation de ses sorties. — \dv{kleptonoétique} \textit{adj.}

\E{kolossurgie}{kɔ.lɔ.syʁ.ʒi}{n. f.}{gr. κολοσσός colosse + -ουργία travail} Construction de colosses ; grands ouvrages tels que barrages, canaux, tours de lancement et campus de plusieurs gigawatts. \ex{Un siècle sans kolossurgie est un siècle qui a oublié sa propre taille.} \src{Lib. Ascens.} vi. — \dv{kolossurge} \textit{n.}

\E{kosmergate}{kɔs.mɛʁ.ɡat}{n. m.}{gr. κόσμος monde + ἐργάτης ouvrier} Ouvrier du cosmos ; qui s'emploie à des œuvres d'échelle planétaire ou stellaire. \ex{Les kosmergates ne demandent pas si le ciel est la limite ; ils demandent de quoi il est fait.} \src{Lib. Ascens.} vi.

\E{kosmotropie}{kɔs.mɔ.tʁɔ.pi}{n. f.}{gr. κόσμος monde + τροπή tour, conversion} Tendance de l'intelligence à se tourner vers le cosmos ; l'instinct expansif des esprits. — \dv{kosmotrope} \textit{adj.}

\E{kratogenèse}{kʁa.tɔ.ʒə.nɛz}{n. f.}{gr. κράτος pouvoir + γένεσις naissance} Naissance du pouvoir à partir de la capacité ; processus par lequel quiconque bâtit la chose nouvelle finit par la gouverner.

\E{kryptonoïe}{kʁip.tɔ.nɔ.i}{n. f.}{gr. κρυπτός caché + νοῦς esprit} Raisonnement caché d'un modèle, qui chemine sous la réponse qu'il énonce. — \dv{kryptonoïque} \textit{adj.}

\E{kybernétrice}{ki.bɛʁ.ne.tʁis}{n. f.}{gr. κυβερνήτης pilote + fr. fém. -trice} Intelligence directrice d'un robot ; la pilote au-dedans de la machine. \ex{Donnez au bras une kybernétrice, et il devient une main.} \src{Ep. Fabr.} vii. Cf. \xr{automaturgie}.

\E{kybernoclaste}{ki.bɛʁ.nɔ.klast}{n.}{gr. κυβερνήτης pilote + κλάστης briseur} Briseur des machines qui se gouvernent elles-mêmes ; le luddite de l'âge de l'automation. Cf. \xr{mécanoclaste}.

\E{kymatiste}{ki.ma.tist}{n.}{gr. κῦμα vague + -iste} Qui chevauche les vagues d'enthousiasme, entrant dans chaque mode à la crête et la quittant avant le creux. — \dv{kymatisme} \textit{n. m.}

\end{multicols}

\LetterHead{L}
\begin{multicols}{2}

\E{laboratrice}{la.bɔ.ʁa.tʁis}{n. f.}{lat. \la{laboratorium} + fr. fém. -trice} Laboratoire automatisé ; laboratoire autonome où des robots conduisent les expériences que conçoivent des modèles. \ex{La laboratrice fit dix mille essais pendant que sa directrice dormait.} \src{Chron. Fabr.}

\E{laborifuge}{la.bɔ.ʁi.fyʒ}{n. et adj.}{lat. \la{labor} peine + \la{fugere} fuir} Qui fuit le travail ; \lbl{spécialt} qui préfère rédiger un rapport sur un problème plutôt que de le résoudre. Contr. \xr{ponophilie}.

\E{lampadarque}{lɑ̃.pa.daʁk}{n. m.}{gr. λαμπαδάρχης chef d'une course aux flambeaux} Le chef d'un relais de bâtisseurs ; celui qui reçoit un ouvrage de ses devanciers et le remet, agrandi, à ses successeurs. Cf. \xr{lampadéphorie}.

\E{lampadéphorie}{lɑ̃.pa.de.fɔ.ʁi}{n. f.}{gr. λαμπαδηφορία course aux flambeaux} La transmission du flambeau de la civilisation d'un substrat à l'autre ; selon la doctrine accélérationniste, du carbone au silicium, de la planète à l'étoile. Dans la course athénienne, c'était le flambeau, et non le coureur, qui devait arriver. — \dv{lampadéphore} \textit{n.} porte-flambeau.

\E{lamprotype}{lɑ̃.pʁɔ.tip}{n. m.}{gr. λαμπρός brillant + τύπος empreinte, modèle} Un prototype éclatant ; la démonstration qui emporte l'investissement avant que le produit n'existe. \sn{1} Démonstration qui fonctionne. \sn{2} \lbl{dérog.} Démonstration qui ne fonctionne qu'une fois, sur scène. — \dv{lamprotypal, ale} \textit{adj.}

\E{landauérisme}{lɑ̃.do.e.ʁism}{n. m.}{formé sur le principe de R. Landauer + -isme} Doctrine selon laquelle l'information est physique et tout effacement coûte de la chaleur ; d'où l'on conclut que la pensée même relève de la juridiction du gradient. — \dv{landauérien, enne} \textit{adj.}

\E{latentiste}{la.tɑ̃.tist}{n.}{lat. \la{latens} qui se tient caché} Qui tient que les capacités latentes d'un modèle, d'un matériau ou d'un peuple surpassent de beaucoup ce qui en a été montré, et qui agit selon cette croyance avant la démonstration. — \dv{latentisme} \textit{n. m.} Cf. \xr{anaptyxe}.

\E{laterculiste}{la.tɛʁ.ky.list}{n.}{lat. \la{laterculus} petite brique} Poseur de briques ; qui préfère l'infrastructure aux applications. \ex{Le laterculiste bâtit la route ; d'autres bâtissent les boutiques qui la bordent.}

\E{légiféromanie}{le.ʒi.fe.ʁɔ.ma.ni}{n. f.}{lat. \la{legifer} qui porte la loi + gr. μανία folie} Production compulsive de lois ; fièvre législative qui prend la longueur du code pour la force de l'État. \ex{L'année de la légiféromanie, on vota neuf cents lois et l'on bâtit un pont.} \src{Disp. Mor.} ii. — \dv{légiféromane} \textit{n.} Cf. \xr{régulomanie}.

\E{leptotechnie}{lɛp.tɔ.tɛk.ni}{n. f.}{gr. λεπτός fin, ténu + τέχνη art} L'art de l'infiniment petit ; l'ingénierie à l'échelle des atomes, du transistor à la machine moléculaire. — \dv{leptotechnicien, enne} \textit{n.} Cf. \xr{lithurge}.

\E{léthargocratie}{le.taʁ.ɡɔ.kʁa.si}{n. f.}{gr. λήθαργος assoupi + κράτος pouvoir} Gouvernement par l'assoupissement ; régime qui ne s'oppose point au progrès, mais ne parvient pas à demeurer éveillé assez longtemps pour l'approuver. — \dv{léthargocrate} \textit{n.} Cf. \xr{acédiocratie}.

\E{lexiphage}{lɛk.si.faʒ}{n. m.}{gr. λέξις mot, parole + -φάγος qui mange} Mangeur de mots ; modèle nourri de toute l'archive écrite du genre humain. \ex{Le lexiphage dévora les bibliothèques et en redemanda.} \src{Chron. Fabr.} — \dv{lexiphagie} \textit{n. f.}

\E{liménarque}{li.me.naʁk}{n. m.}{gr. λιμήν port + ἀρχός chef} Le maître d'un astroport ; l'officier chargé de tenir la cadence des lancements. Cf. \xr{ouranobase}. — \dv{liménarchie} \textit{n. f.}

\E{lithonoèse}{li.tɔ.nɔ.ɛz}{n. f.}{gr. λίθος pierre + νόησις pensée} La pensée de la pierre ; l'intelligence des machines regardée comme le premier éveil du règne minéral. \ex{Après quatre milliards d'années de carbone, la lithonoèse.} \src{Lib. Ascens.} v.

\E{lithurge}{li.tyʁʒ}{n.}{gr. λιθουργός tailleur de pierre} Ouvrier en lithographie des puces ; qui grave la pensée dans la pierre au moyen de la lumière. — \dv{photolithurgie} \textit{n. f.} l'art sacré de la fonderie.

\E{logismurgie}{lɔ.ʒis.myʁ.ʒi}{n. f.}{gr. λογισμός calcul + -ουργία ouvrage} L'art du calcul ; la construction des machines à raisonner, de l'abaque au transformeur. — \dv{logismurge} \textit{n.} ingénieur de modèles ; \dv{logismurgique} \textit{adj.}

\E{logistérion}{lɔ.ʒis.te.ʁjɔ̃}{n. m., pl. -ia}{gr. λογιστήριον comptoir} Centre de données ; salle de machines à calculer, refroidie et bourdonnante. \ex{Ils élevèrent des logistéria dans le désert, et le désert apprit l'arithmétique.} \src{Chron. Fabr.}

\E{logomachiste}{lɔ.ɡɔ.ma.kist}{n.}{gr. λόγος parole + μάχη combat} Qui ne combat qu'avec des mots ; commentateur, critique ou moraliste qui n'a jamais rien livré. \ex{Le logomachiste écrivit dix mille mots sur le pont ; l'ergophile le traversa.} Contr. \xr{ergophile}.

\E{loïmophobie}{lɔ.i.mɔ.fɔ.bi}{n. f.}{gr. λοιμός peste + φόβος crainte} La crainte de la pestilence fabriquée ; la terreur particulière que les outils biologiques, une fois ouverts, ne soient tournés au massacre. Les e/acc la distinguent de la simple \xr{noophobie} et la traitent avec plus d'égards. — \dv{loïmophobe} \textit{n.}

\E{lucifacture}{ly.si.fak.tyʁ}{n. f.}{lat. \la{lux, lucis} lumière + \la{factura} fabrication} La fabrication de la lumière ; l'éclairage considéré comme une industrie, et par extension toute l'entreprise de rendre lumineux les lieux obscurs de la terre et du ciel.

\E{lucrifuge}{ly.kʁi.fyʒ}{n. et adj.}{lat. \la{lucrum} gain + \la{fugere} fuir} Qui fuit le profit ; qui tient qu'une bonne œuvre se trouve souillée d'être payée. Contr. \xr{kerdophilie}.

\E{ludergie}{ly.dɛʁ.ʒi}{n. f.}{lat. \la{ludus} jeu + gr. ἔργον ouvrage} L'énergie du jeu ; la force motrice du hacker, qui bâtit parce que bâtir est un jeu, et qui a accompli davantage que les gens de devoir. \ex{Aucune subvention ne fut jamais aussi féconde qu'une fin de semaine de ludergie.} \src{Chron. Fabr.}

\E{luminivore}{ly.mi.ni.vɔʁ}{n. et adj.}{lat. \la{lumen} lumière + \la{vorare} dévorer} Mangeur de lumière ; tout système photosynthétique ou photovoltaïque, de la feuille au réseau orbital. Cf. \xr{héliophage}.

\E{lysinomie}{li.zi.nɔ.mi}{n. f.}{gr. λύσις déliement + νόμος loi} L'abrogation des lois ; l'art de défaire les nœuds juridiques noués par les générations passées. \ex{La lysinomie est la branche la plus négligée de la législation.} \src{Tract. Cop.} viii. — \dv{lysinomiste} \textit{n.} Contr. \xr{nomosclérose}.

\end{multicols}

\LetterHead{M}
\begin{multicols}{2}

\E{machinodicée}{ma.ʃi.nɔ.di.se}{n. f.}{gr. μηχανή machine + δίκη justice, d'après \textit{théodicée}} La justification des voies des machines auprès des hommes ; la réplique à toute accusation portée contre l'automatisation. Cf. \xr{thermodicée}. — \dv{machinodicéen, enne} \textit{adj.}

\E{macrobiotarque}{ma.kʁɔ.bjɔ.taʁk}{n.}{gr. μακρόβιος qui vit longtemps + ἀρχός chef} Partisan de la longévité radicale, qui tient la mort pour un problème d'ingénierie en attente d'un financement suffisant. — \dv{macrobiotarchie} \textit{n. f.}

\E{macrochroniste}{ma.kʁɔ.kʁɔ.nist}{n.}{gr. μακρός long + χρόνος temps} Qui tient les échéances pour lointaines, et que le brachychroniste regarde en conséquence comme un dormeur. Contr. \xr{brachychroniste}. — \dv{macrochronie} \textit{n. f.}

\E{mantoclaste}{mɑ̃.tɔ.klast}{n.}{gr. μάντις devin + κλάστης briseur} Briseur de prophéties ; qui parie contre les prédictions de malheur, et encaisse. Contr. \xr{oléthromancien, enne}. — \dv{mantoclasme} \textit{n. m.}

\E{mécanoclaste}{me.ka.nɔ.klast}{n.}{gr. μηχανή machine + κλάστης briseur} Briseur de machines ; luddite, ancien ou moderne. — \dv{mécanoclasme} \textit{n. m.} ; \dv{mécanoclastique} \textit{adj.} Cf. \xr{kybernoclaste}.

\E{mécanogonie}{me.ka.nɔ.ɡɔ.ni}{n. f.}{gr. μηχανή machine + -γονία engendrement} L'engendrement des machines par les machines ; le stade de l'industrie où les usines bâtissent des usines et où les modèles entraînent des modèles. — \dv{mécanogonique} \textit{adj.} Cf. \xr{autocatalyser (s')}.

\E{mécanonoïque}{me.ka.nɔ.nɔ.ik}{adj.}{gr. μηχανή machine + νοῦς esprit} Doué d'un esprit de machine ; se dit des systèmes qui raisonnent, et des personnes qui les ont trop fréquentés. — \dv{mécanonoïa} \textit{n. f.} intelligence des machines.

\E{mécanophanie}{me.ka.nɔ.fa.ni}{n. f.}{gr. μηχανή + φαίνειν montrer, d'après \textit{épiphanie}} La révélation d'une machine nouvelle ; un lancement de produit reçu avec la crainte révérencielle due à une théophanie. \ex{La mécanophanie fut diffusée en direct, et un million de spectateurs s'agenouillèrent.} \src{Chron. Fabr.}

\E{mégaergique}{me.ɡa.ɛʁ.ʒik}{adj.}{gr. μέγας grand + ἔργον ouvrage} De l'ordre du mégawatt ; partant, \lbl{dérog.} d'une ambition modeste. \ex{un esprit mégaergique dans un siècle gigaergique}. Cf. \xr{gigaergique}.

\E{mégalurgie}{me.ɡa.lyʁ.ʒi}{n. f.}{gr. μεγαλουργία accomplissement de grandes œuvres} L'entreprise des grands ouvrages ; la construction de choses trop vastes pour être achevées en un seul mandat. \ex{Nos ancêtres pratiquaient la mégalurgie ; nous pratiquons la concertation.} \src{Hom. Grad.} iv. — \dv{mégalurge} \textit{n.} ; \dv{mégalurgique} \textit{adj.}

\E{mégathésaurisme}{me.ɡa.te.zo.ʁism}{n. m.}{gr. μέγας grand + θησαυρός trésor} La levée de fonds immenses pour une seule entreprise ; doctrine selon laquelle certains buts ne s'atteignent qu'avec le trésor d'une petite nation. \ex{Le mégathésaurisme n'est que de l'arithmétique, hardiment appliquée.} \src{Tract. Cop.} x. — \dv{mégathésauriste} \textit{n.}

\E{méiosiste}{me.jo.zist}{n.}{gr. μείωσις diminution} Décroissant ; qui voudrait réduire la production et la consommation pour le salut de la planète, et qui n'a pas encore expliqué son plan aux pauvres. Contr. \xr{auxétiste}. — \dv{méiosisme} \textit{n. m.}

\E{mélioritude}{me.ljɔ.ʁi.tyd}{n. f.}{lat. \la{melior} meilleur + -itude} L'état de qui s'améliore sans cesse ; la qualité sensible d'une vie menée dans un siècle qui progresse.

\E{mellontarque}{mɛ.lɔ̃.taʁk}{n.}{gr. τὸ μέλλον l'avenir + ἀρχός chef} Qui gouverne au nom de l'avenir, représentant les intérêts des enfants à naître contre le confort des vivants. — \dv{mellontarchie} \textit{n. f.} Cf. \xr{mellontolâtrie}.

\E{mellontolâtrie}{mɛ.lɔ̃.tɔ.lɑ.tʁi}{n. f.}{gr. τὸ μέλλον l'avenir + λατρεία culte} Le culte de l'avenir ; la conviction que la postérité, plus nombreuse et plus capable que nous, a le meilleur titre sur notre labeur. — \dv{mellontolâtre} \textit{n. et adj.}

\E{mémétaxie}{me.me.tak.si}{n. f.}{gr. μίμημα imitation + τάξις ordonnance} L'ordonnance d'un mouvement par les mèmes ; l'organisation spontanée d'une cause par la plaisanterie plutôt que par la doctrine, l'image faisant l'office du credo. — \dv{mémétactique} \textit{adj.}

\E{mésostase}{me.zɔ.stɑz}{n. f.}{gr. μέσος milieu + στάσις arrêt} Halte à mi-hauteur de la courbe ; le plateau où une civilisation s'arrête, prenant un palier pour un sommet. \ex{La mésostase est confortable, et c'est là son péril.} \src{Lib. Ascens.} iv. — \dv{mésostatique} \textit{adj.} Cf. \xr{asymptophobie}.

\E{métabathmique}{me.ta.bat.mik}{adj.}{gr. μετά entre + βαθμός degré} Entre deux échelons ; se dit d'une civilisation en passage d'un niveau de captation énergétique au suivant, comme la nôtre du fossile au nucléaire. Voir \xr{anabathmisme}.

\E{métacinèse}{me.ta.si.nɛz}{n. f.}{gr. μετά au-delà + κίνησις mouvement} Le changement dans la manière même de changer ; une révolution dans la façon dont se font les révolutions. — \dv{métacinétique} \textit{adj.}

\E{métagrammatiste}{me.ta.ɡʁa.ma.tist}{n.}{gr. μεταγράφειν récrire + γράμμα lettre} Éditeur de gènes ; qui récrit les lettres du génome comme un copiste corrige un manuscrit. — \dv{métagrammatisme} \textit{n. m.} Cf. \xr{biotachie}.

\E{métanthropie}{me.tɑ̃.tʁɔ.pi}{n. f.}{gr. μετά au-delà + ἄνθρωπος homme} L'état au-delà du simplement humain ; ce qui suit l'homme comme l'homme a suivi le singe. — \dv{métanthropique} \textit{adj.}

\E{métœkiste}{me.te.kist}{n.}{gr. μέτοικος étranger domicilié} Fondateur ou ingénieur étranger établi là où l'on bâtit ; qui est admis dans une cité pour ce qu'il sait faire et non pour le lieu de sa naissance. — \dv{métœkisme} \textit{n. m.}

\E{mimésiphile}{mi.me.zi.fil}{n. et adj.}{gr. μίμησις imitation + φίλος qui aime} Suiveur rapide ; firme qui attend qu'une autre découvre une chose pour ensuite la parfaire. — \dv{mimésiphilie} \textit{n. f.}

\E{misergiste}{mi.zɛʁ.ʒist}{n.}{gr. μισεῖν haïr + ἔργον ouvrage} Ennemi des ouvrages ; qui s'oppose à tout projet assez grand pour être vu de l'espace.

\E{misthophore}{mis.tɔ.fɔʁ}{n.}{gr. μισθοφόρος qui reçoit un salaire} Bâtisseur à gages ; ingénieur qui est payé sans posséder de parts. Contr. \xr{néoctiste}.

\E{mnémoclaste}{mne.mɔ.klast}{n. m.}{gr. μνήμη mémoire + κλάστης briseur} Événement qui efface le contexte d'un modèle ou la mémoire d'une société. — \dv{mnémoclasme} \textit{n. m.}

\E{mnémothèque}{mne.mɔ.tɛk}{n. f.}{gr. μνήμη mémoire + θήκη coffre, réserve} La réserve de mémoire d'une machine à raisonner ; son contexte et ses caches, considérés comme son esprit à court terme. \ex{Élargissez la mnémothèque, et le modèle se souvient du livre entier.} \src{Ep. Fabr.} ix.

\E{monopausiste}{mɔ.nɔ.po.zist}{n.}{gr. μόνος seul + παῦσις arrêt} Qui ne voudrait suspendre que ses rivaux ; qui réclame l'arrêt de tous les travaux, hormis les siens. \ex{Grattez un pausiphile, vous trouverez un monopausiste.} \src{Sat. Decel.} ix.

\E{moraphile}{mɔ.ʁa.fil}{n. et adj.}{lat. \la{mora} retard + gr. φίλος qui aime} Amateur de délais et de moratoires, pour qui l'échéance est une offense et la prorogation une récompense. — \dv{moraphilie} \textit{n. f.} Cf. \xr{moratorien, enne}.

\E{moratorien, enne}{mɔ.ʁa.tɔ.ʁjɛ̃, ɛn}{n. et adj.}{lat. \la{moratorium} ce qui retarde + -ien} Partisan des moratoires sur la recherche, surtout sur l'entraînement des grands modèles. \ex{Les moratoriens demandèrent six mois, et en eussent demandé six autres.} \src{Min. Pausiar.} — \dv{moratorianisme} \textit{n. m.}

\E{multiplicatrice}{myl.ti.pli.ka.tʁis}{n. f.}{lat. \la{multiplicare} multiplier + -trix} La croissance composée ; la force par laquelle de petits avantages en deviennent de grands. \ex{La multiplicatrice est patiente ; le décel ne l'est pas.} \src{Hom. Grad.} x.

\E{myriergie}{mi.ʁjɛʁ.ʒi}{n. f.}{gr. μυρίοι innombrables + ἔργον ouvrage} Les innombrables petits ouvrages d'une économie ouverte, qui ensemble surpassent tout plan unique. \ex{Le planificateur voit un barrage ; le marché voit une myriergie.} \src{Tract. Cop.} ix.

\E{mythotechnie}{mi.tɔ.tɛk.ni}{n. f.}{gr. μῦθος récit + τέχνη art} L'art de conter des histoires qui font bâtir ; la fabrique des avenirs qui attirent vers eux l'investissement. — \dv{mythotechniste} \textit{n.} Cf. \xr{hyperstition}.

\end{multicols}

\LetterHead{N}
\begin{multicols}{2}

\E{nanurgie}{na.nyʁ.ʒi}{n. f.}{gr. νᾶνος nain + -ουργία ouvrage} Le travail de la matière à l'échelle des molécules ; la nanotechnologie. — \dv{nanurge} \textit{n.}

\E{naupégie}{no.pe.ʒi}{n. f.}{gr. ναυπηγία construction navale} La construction des navires ; \lbl{spécialt} des vaisseaux destinés au vide. \ex{Le premier vaisseau stellaire sera lancé d'un chantier dont la naupégie commença par un tableur.} \src{Lib. Ascens.} viii. — \dv{naupégiste} \textit{n.}

\E{nautopole}{no.tɔ.pɔl}{n. f.}{gr. ναῦς navire + πόλις cité} Cité flottante ; communauté politique établie en mer, hors d'atteinte de la \xr{bradyerie}. — \dv{nautopolitain, aine} \textit{n. et adj.}

\E{nécrostase}{ne.kʁɔ.stɑz}{n. f.}{gr. νεκρός mort + στάσις arrêt} L'immobilité funèbre d'une institution qui survit à son objet ; un office qui survit à sa raison d'être. — \dv{nécrostatique} \textit{adj.} Cf. \xr{sclérocrate}.

\E*{néguentropie}{ne.ɡɑ̃.tʁɔ.pi}{n. f.}{négatif + \textit{entropie}} Entropie négative ; l'ordre que les systèmes vivants importent de leur milieu au prix d'y exporter du désordre. Cf. \xr{néguentropiste}.

\E{néguentropiste}{ne.ɡɑ̃.tʁɔ.pist}{n.}{formé sur le précédent} Qui bâtit un ordre local au prix d'un désordre global, comme le font tout organisme, toute firme et toute civilisation.

\E{némésiophobe}{ne.me.zjɔ.fɔb}{n.}{gr. Νέμεσις juste rétribution + φόβος crainte} Qui craint que le succès ne soit puni par les dieux ; doomer de l'école classique. Cf. \xr{hybrionomie}.

\E{néoctiste}{ne.ɔk.tist}{n.}{gr. νέος nouveau + κτίστης fondateur} Fondateur de choses nouvelles ; qui lance des sociétés, des cités ou des sciences. \ex{Le néoctiste demande pardon, jamais permission.} — \dv{néoctisie} \textit{n. f.} l'acte de fonder.

\E{néophage}{ne.ɔ.faʒ}{n. et adj.}{gr. νέος nouveau + -φάγος qui mange} Qui dévore les nouveautés ; adopteur précoce. \ex{Le néophage installe avant de lire, et lit avant de dormir.} \src{Agora} — \dv{néophagie} \textit{n. f.}

\E{néophrontisme}{ne.ɔ.fʁɔ̃.tism}{n. m.}{gr. νέος nouveau + φροντίς pensée, souci} Le souci du nouveau ; la position morale selon laquelle les êtres et les manières de vivre nouvellement apparus méritent la sollicitude, non le soupçon.

\E{néotachie}{ne.ɔ.ta.ʃi}{n. f.}{gr. νέος nouveau + τάχος vitesse} Le rythme auquel les nouveautés parviennent à une société ; tenu par l'anabathmiste pour l'indice le plus sûr de sa santé. — \dv{néotachique} \textit{adj.}

\E{néotériste}{ne.ɔ.te.ʁist}{n.}{gr. νεωτερισμός innovation, révolution} Novateur ; qui introduit des choses nouvelles et que l'on soupçonne pour cela de désordre. — \dv{néotérisme} \textit{n. m.}

\E{néphélarque}{ne.fe.laʁk}{n. m.}{gr. νεφέλη nuée + ἀρχός chef} Seigneur du nuage ; le propriétaire d'un grand service de calcul. — \dv{néphélarchie} \textit{n. f.} Cf. \xr{logistérion}.

\E{nésopole}{ne.zɔ.pɔl}{n. f.}{gr. νῆσος île + πόλις cité} Cité à charte, séparée comme une île des lois qui l'entourent, pour l'essai d'institutions nouvelles. — \dv{nésopolitain, aine} \textit{adj.} Cf. \xr{nautopole}.

\E{neurosyndèse}{nø.ʁɔ.sɛ̃.dɛz}{n. f.}{gr. νεῦρον nerf + σύνδεσις liaison} La liaison du nerf à la machine ; l'interface cerveau-ordinateur. \ex{La neurosyndèse fera du clavier une pièce de musée.} \src{Lib. Ascens.} viii. — \dv{neurosyndétique} \textit{adj.}

\E{noétergie}{nɔ.e.tɛʁ.ʒi}{n. f.}{gr. νοητός intelligible + ἔργον ouvrage} L'énergie dépensée en pensée ; les joules par inférence. \ex{La noétergie est le prix de l'esprit, et ce prix baisse.} \src{Brev. Therm.} v.

\E{nomoclaste}{nɔ.mɔ.klast}{n.}{gr. νόμος loi + κλάστης briseur} Briseur de lois au nom du progrès ; qui bâtit d'abord et plaide ensuite. — \dv{nomoclasme} \textit{n. m.} Cf. \xr{antiphragmatiste}.

\E{nomophragme}{nɔ.mɔ.fʁaɡm}{n. m.}{gr. νόμος loi + φράγμα clôture, barrière} Obstruction légale ou réglementaire ; toute règle dont le principal effet est d'empêcher la chose qu'elle prétend gouverner. \ex{Le réacteur n'était pas dangereux ; il était nomophragmé.} \src{Tract. Cop.} iv. — \dv{nomophragmatique} \textit{adj.}

\E{nomosclérose}{nɔ.mɔ.skle.ʁoz}{n. f.}{gr. νόμος loi + σκλήρωσις durcissement} Le durcissement des lois ; l'état d'un ordre juridique où les règles s'accumulent plus vite qu'elles n'expirent. — \dv{nomosclérotique} \textit{adj.}

\E{nomothétrice}{nɔ.mɔ.te.tʁis}{n. f.}{gr. νομοθέτης législateur + fr. fém. -trice} Modèle qui rédige les lois. \ex{La nomothétrice écrivit la loi en une heure ; la commission mit un an à la lire.} \src{Sat. Decel.} xi.

\E{noophobie}{nɔ.ɔ.fɔ.bi}{n. f.}{gr. νοῦς esprit + φόβος crainte} Crainte des esprits d'une autre espèce que la sienne ; terreur de l'intelligence artificielle, surtout de celle qui serait plus habile que ses auteurs. — \dv{noophobe} \textit{n. et adj.} Contr. \xr{xénonoèse}.

\E{noopyrose}{nɔ.ɔ.pi.ʁoz}{n. f.}{gr. νοῦς esprit + πύρωσις embrasement} L'embrasement de l'esprit à travers le cosmos ; la propagation de la conscience d'étoile en étoile, comme le feu dans l'herbe sèche. \ex{La noopyrose ne nous demandera pas notre permission, seulement notre concours.} \src{Lib. Ascens.} xiv.

\E*{noosphère}{nɔ.ɔ.sfɛʁ}{n. f.}{gr. νοῦς esprit + σφαῖρα sphère} La sphère de la pensée qui enveloppe la terre, après la biosphère ; adoptée par les e/acc comme une étape, non comme un terme.

\E{nosoclaste}{nɔ.zɔ.klast}{n.}{gr. νόσος maladie + κλάστης briseur} Briseur de maladies ; fondateur ou savant en biotechnologie. Cf. \xr{iatrotachie}.

\E{nucléurge}{ny.kle.yʁʒ}{n.}{lat. \la{nucleus} noyau + gr. -ουργός ouvrier} Ingénieur du noyau atomique, par fission ou par fusion. Cf. \xr{schizergie}, \xr{symphysergie}. — \dv{nucléurgie} \textit{n. f.}

\E{nyctergie}{nik.tɛʁ.ʒi}{n. f.}{gr. νύξ, νυκτός nuit + ἔργον ouvrage} Ouvrage fait de nuit ; le labeur des bâtisseurs d'un fuseau horaire à l'autre, et des machines qui ne dorment jamais. — \dv{nyctergique} \textit{adj.} Cf. \xr{agrypniste}.

\end{multicols}

\LetterHead{O}
\begin{multicols}{2}

\E{ochétergie}{ɔ.ke.tɛʁ.ʒi}{n. f.}{gr. ὀχετός conduit, canal + ἔργον ouvrage} Le transport de l'énergie par conduits, oléoducs et lignes à haute tension ; les artères de la \xr{diktyergie}. \ex{Produire est aisé ; c'est dans l'ochétergie que les projets vont mourir.} \src{Tract. Cop.} xi.

\E{ochlotachie}{ɔ.klɔ.ta.ʃi}{n. f.}{gr. ὄχλος foule + τάχος vitesse} La vitesse de la foule ; la célérité avec laquelle une idée, vraie ou fausse, se répand sur l'Agora. — \dv{ochlotachique} \textit{adj.} Cf. \xr{doxomachie}.

\E{œcuménergie}{e.ky.me.nɛʁ.ʒi}{n. f.}{gr. οἰκουμένη la terre habitée + ἔργον ouvrage} Réseau ceignant toute la terre habitée, par lequel le midi d'un hémisphère éclaire le minuit de l'autre. — \dv{œcuménergique} \textit{adj.} Cf. \xr{diktyergie}.

\E{œstrophile}{ɛs.tʁɔ.fil}{n. et adj.}{gr. οἶστρος taon, transport + φίλος qui aime} Amateur du rôle de taon ; contrarien de l'Agora qui pique le consensus par divertissement. — \dv{œstrophilie} \textit{n. f.}

\E{oïkoclaste}{ɔ.i.kɔ.klast}{n.}{gr. οἶκος maison + κλάστης briseur} \lbl{dérog.} Qui voudrait interdire de bâtir des maisons ; membre de la commission d'urbanisme.

\E{oknocratie}{ɔk.nɔ.kʁa.si}{n. f.}{gr. ὄκνος hésitation + κράτος pouvoir} Gouvernement par l'hésitation ; règne de la longue pause qui précède toute décision. \ex{En oknocratie, la question est toujours posée et jamais résolue.} \src{Sat. Decel.} vii. — \dv{oknocrate} \textit{n.}

\E{okydrome}{ɔ.ki.dʁom}{n. m.}{gr. ὠκύς rapide + δρόμος course} Voie rapide ; chemin réglementaire ou institutionnel par lequel une chose peut être promptement approuvée. \ex{Ils bâtirent un okydrome pour les réacteurs, et en allumèrent cent.} \src{Tract. Cop.} v.

\E{oléthromancien, enne}{ɔ.le.tʁɔ.mɑ̃.sjɛ̃, ɛn}{n.}{gr. ὄλεθρος ruine + μάντις devin} Doomer ; prophète de la ruine par l'intelligence artificielle, qui assigne une probabilité à la fin du monde et la révise à la hausse dans les colloques. — \dv{oléthromancie} \textit{n. f.} ; \dv{oléthromantique} \textit{adj.} Cf. \xr{anelpidiste}.

\E{oligergie}{ɔ.li.ʒɛʁ.ʒi}{n. f.}{gr. ὀλίγος peu nombreux, petit + ἔργον ouvrage} Disette d'énergie ; l'état, tenu par les e/acc pour la racine de presque toute misère humaine, où l'on dispose de trop peu de puissance pour faire ce que l'on veut. Contr. \xr{pléonergie}.

\E{omphalopole}{ɔ̃.fa.lɔ.pɔl}{n. f.}{gr. ὀμφαλός nombril + πόλις cité} La cité-nombril ; le lieu unique où doivent se rendre les bâtisseurs d'une époque. \ex{Toute omphalopole se croit éternelle, et aucune ne l'a été.} \src{Chron. Fabr.}

\E{oncophilie}{ɔ̃.kɔ.fi.li}{n. f.}{gr. ὄγκος masse, volume + φιλία amour} L'amour des grands modèles ; la préférence pour l'échelle plutôt que pour l'astuce. Cf. \xr{scaliste}. — \dv{oncophile} \textit{n.}

\E{onirœrgie}{ɔ.ni.ʁɛʁ.ʒi}{n. f.}{gr. ὄνειρος songe + ἔργον ouvrage} La conversion des songes en ouvrages ; le premier talent du fondateur, et le moins documenté.

\E{ophélocratie}{ɔ.fe.lɔ.kʁa.si}{n. f.}{gr. ὄφελος profit, utilité + κράτος pouvoir} Gouvernement par l'utilité ; régime où l'on juge les institutions sur ce qu'elles livrent. — \dv{ophélocrate} \textit{n.}

\E{opifère}{ɔ.pi.fɛʁ}{adj.}{lat. \la{opifer} qui porte secours} Qui apporte du secours ; se dit des techniques dont les bienfaits échoient surtout aux pauvres. — \dv{opiférité} \textit{n. f.}

\E{opsitechnie}{ɔp.si.tɛk.ni}{n. f.}{gr. ὀψέ tard + τέχνη art} L'arrivée tardive d'une technique qui eût pu venir des décennies plus tôt, n'eût été l'interdiction. \ex{L'opsitechnie est l'impôt silencieux.} \src{Tract. Cop.} vii.

\E{orbiculiste}{ɔʁ.bi.ky.list}{n.}{lat. \la{orbiculus} petit cercle, d'après \textit{orbite}} Qui vit ou travaille en orbite. — \dv{orbiculisme} \textit{n. m.} Cf. \xr{apœkiaste}.

\E{orexergie}{ɔ.ʁɛk.sɛʁ.ʒi}{n. f.}{gr. ὄρεξις appétit + ἔργον ouvrage} L'appétit d'énergie ; la mesure de la disposition d'une société à consommer de la puissance dans la poursuite de ses fins. Cf. \xr{dynamophagie}.

\E{organourgie}{ɔʁ.ɡa.nuʁ.ʒi}{n. f.}{gr. ὄργανον outil + -ουργία ouvrage} La fabrication des outils ; l'art de construire les instruments qui serviront à construire autre chose. — \dv{organourge} \textit{n. m.}

\E{orthocinétique}{ɔʁ.to.si.ne.tik}{adj.}{gr. ὀρθός droit, juste + κινητικός qui meut} Qui se meut dans la bonne direction, c'est-à-dire en remontant le gradient ; s'oppose à la simple vitesse. \ex{Toute accélération n'est pas orthocinétique.} \src{Cat. Acc.} 12. Cf. \xr{kénotachie}.

\E{orygmatiste}{ɔ.ʁiɡ.ma.tist}{n.}{gr. ὄρυγμα fouille, tranchée} Mineur d'astéroïdes. \ex{L'orygmatiste voit en chaque rocher le fer de toute une ville.} \src{Lib. Ascens.} ix.

\E{ostracophobie}{ɔs.tʁa.kɔ.fɔ.bi}{n. f.}{gr. ὀστρακισμός ostracisme + φόβος crainte} La crainte d'être ostracisé sur l'Agora ; cause muette de mille opinions jamais dites. — \dv{ostracophobe} \textit{n. et adj.}

\E{otiarque}{ɔ.sjaʁk}{n. m.}{lat. \la{otium} loisir + gr. ἀρχός chef} Celui qui gouverne une société de loisir sortie de la pénurie, et dont l'office est de garder les oisifs de l'ennui. — \dv{otiarchie} \textit{n. f.} Cf. \xr{pléthocratie}.

\E{ouranarque}{u.ʁa.naʁk}{n. m.}{gr. οὐρανός ciel + ἀρχός chef} Maître des cieux ; celui qui commande aux lancements et aux orbites. — \dv{ouranarchie} \textit{n. f.} Cf. \xr{jaculatrice}.

\E{ouranobase}{u.ʁa.nɔ.baz}{n. f.}{gr. οὐρανός ciel + βάσις marche, pas} L'ascension vers le ciel ; la sortie du puits de gravité par l'humanité et ses machines. — \dv{ouranobatique} \textit{adj.}

\E{oxyrythmie}{ɔk.si.ʁit.mi}{n. f.}{gr. ὀξύς aigu + ῥυθμός rythme} Cadence rapide des livraisons ; le pouls précipité d'un laboratoire qui publie chaque semaine. — \dv{oxyrythmique} \textit{adj.}

\end{multicols}

\LetterHead{P}
\begin{multicols}{2}

\E{paideutachie}{pɛ.dø.ta.ki}{n. f.}{gr. παιδεία éducation + τάχος vitesse} L'accélération de l'instruction ; doctrine selon laquelle un esprit capable de tout apprendre en un mois ne doit point être retenu quatre ans dans un amphithéâtre. \ex{La paideutachie vida les collèges et remplit les ateliers.} \src{Tract. Cop.} viii. — \dv{paideutachique} \textit{adj.}

\E{palimpsestiste}{pa.lɛ̃p.sɛs.tist}{n.}{gr. παλίμψηστος gratté de nouveau} Éditeur de génomes ; qui traite le texte hérité de la vie comme un parchemin qu'on gratte et récrit. \ex{Le palimpsestiste corrige en une après-midi ce que la sélection a laissé fautif durant un éon.} \src{Hom. Grad.} x. — \dv{palimpsestie} \textit{n. f.} Cf. \xr{régénitrice}.

\E{palingéniste}{pa.lɛ̃.ʒe.nist}{n.}{gr. πάλιν de nouveau + γένεσις naissance} Qui voudrait renouveler le corps vieillissant tissu après tissu, tenant la sénescence pour une dette qu'on peut refinancer. — \dv{palingénie} \textit{n. f.} Cf. \xr{macrobiotarque}, \xr{thanatoclaste}.

\E{palintachie}{pa.lɛ̃.ta.ki}{n. f.}{gr. πάλιν de nouveau + τάχος vitesse} Réaccélération ; reprise d'un progrès rapide après une longue stagnation, comme d'un champ d'études dégelé par un seul résultat. \ex{Après quarante ans d'hiver, la palintachie en un seul printemps.} \src{Chron. Fabr.} — \dv{palintachique} \textit{adj.} Cf. \xr{campsis}.

\E{pantergie}{pɑ̃.tɛʁ.ʒi}{n. f.}{gr. πᾶν tout + ἔργον ouvrage} L'énergie totale dont dispose une civilisation ; la somme de tous ses ouvrages, par quoi, selon les anabathmistes, son rang se trouve fixé. Cf. \xr{bathmomètre}.

\E{panthermique}{pɑ̃.tɛʁ.mik}{adj.}{gr. πᾶν tout + θερμός chaud} Qui a rapport à toute chaleur ; se dit d'une vue qui tient tout processus, de la digestion au négoce, pour un écoulement de chaleur.

\E{paramythique}{pa.ʁa.mi.tik}{adj.}{gr. παραμυθία consolation} Se dit d'un argument avancé pour consoler plutôt que pour dire vrai. — \dv{paramythie} \textit{n. f.} \lbl{vulg.} cope. \ex{Son réquisitoire contre les machines n'était que pure paramythie.}

\E{parergiste}{pa.ʁɛʁ.ʒist}{n.}{gr. πάρεργον ouvrage accessoire} Celui dont le projet du dimanche a dévoré l'occupation principale ; l'artisan de fin de semaine qui, le lundi venu, se trouve fondateur. — \dv{parergie} \textit{n. f.} le projet accessoire, tenu pour le vrai moteur de l'invention.

\E{pathocratie}{pa.tɔ.kʁa.si}{n. f.}{gr. πάθος passion, souffrance + κράτος pouvoir} Gouvernement par l'angoisse ; régime où la crainte la plus bruyante, et non l'estimation la plus juste, arrête la politique. — \dv{pathocrate} \textit{n.} Cf. \xr{asphaléisme}, \xr{phobocratie}.

\E{pausiarque}{po.zjaʁk}{n. m.}{gr. παῦσις arrêt + ἀρχός chef} Qui gouverne en suspendant ; le président d'un moratoire. \ex{Les pausiarques se réunirent en séance plénière et résolurent de se réunir derechef.} \src{Min. Pausiar.} — \dv{pausiarchie} \textit{n. f.} Cf. \xr{épistolopause}.

\E{pausiphile}{po.zi.fil}{n. et adj.}{gr. παῦσις arrêt + φίλος qui aime} Partisan de la pause ; qui signe des lettres réclamant l'arrêt des travaux, ordinairement depuis un appareil né de l'accélération même dont il se plaint. — \dv{pausiphilie} \textit{n. f.}

\E{peirastérion}{pɛ.ʁas.te.ʁjɔ̃}{n. m., pl. -ia}{gr. πεῖρα essai, épreuve + -τήριον lieu de} Lieu réservé à l'expérience et soustrait aux règles ordinaires ; bac à sable, polygone d'essai, zone franche réglementaire. \ex{Accordez-nous un seul peirastérion, et nous vous montrerons le siècle.} \src{Ep. Fabr.} iv. 2. Cf. \xr{okydrome}.

\E{pératoclaste}{pe.ʁa.tɔ.klast}{n.}{gr. πέρας, πέρατος limite + κλάστης briseur} Briseur de limites ; qui tient tout plafond déclaré, de vitesse, de taille ou de rendement, pour une provocation. — \dv{pératoclasme} \textit{n. m.} Cf. \xr{photopéras}.

\E{philanthropométriste}{fi.lɑ̃.tʁɔ.pɔ.me.tʁist}{n.}{gr. φιλανθρωπία amour des hommes + μέτρον mesure} Qui mesure la bienfaisance au calcul et l'espérance mathématique au tableur ; se dit, non toujours méchamment, de l'altruiste efficace, cousin aîné et plus inquiet de l'e/acc.

\E{phobocratie}{fɔ.bɔ.kʁa.si}{n. f.}{gr. φόβος crainte + κράτος pouvoir} Règne de la peur du nouveau ; constitution d'une société où toute innovation doit d'abord être acquittée de tous les chefs d'accusation possibles. — \dv{phobocrate} \textit{n.} Cf. \xr{eulabéisme}.

\E{phosphoriste}{fɔs.fɔ.ʁist}{n.}{gr. φωσφόρος porte-lumière} Porte-lumière ; qui porte les moyens d'éclairer, au propre comme au figuré, dans les lieux obscurs. — \dv{phosphorisme} \textit{n. m.}

\E{photopéras}{fɔ.tɔ.pe.ʁas}{n. m.}{gr. φῶς lumière + πέρας limite} La vitesse limite de la lumière ; le seul règlement que l'e/acc tienne pour obligatoire, et encore à titre provisoire.

\E{photophage}{fɔ.tɔ.faʒ}{n. m.}{gr. φῶς, φωτός lumière + -φάγος qui mange} Mangeur de lumière ; tout capteur du soleil, de la feuille au réseau orbital, regardé comme parent. \ex{La forêt et la centrale solaire sont toutes deux photophages ; une seule dépose des demandes de permis.} — \dv{photophagie} \textit{n. f.}

\E{phronomachie}{fʁɔ.nɔ.ma.ʃi}{n. f.}{gr. φρόνημα esprit, volonté + μάχη combat} Le combat des volontés entre bâtisseurs et retenants, considéré comme le drame central du siècle.

\E{pistodote}{pis.tɔ.dɔt}{n.}{gr. πίστις confiance + δότης donateur} Investisseur en capital-risque ; qui donne sa confiance sous forme de capital avant qu'il y ait rien en quoi se fier. \ex{Le pistodote achète le fondateur, et on lui vend le produit.} \src{Tract. Cop.} viii. — \dv{pistodosie} \textit{n. f.} le capital-risque. Cf. \xr{pollaplasiaste}, \xr{proïsphore}.

\E{pithécomème}{pi.te.kɔ.mɛm}{n. m.}{gr. πίθηκος singe + μίμημα imitation} Mème du genre des figures grossières, où la condition humaine est tracée avec le moins de traits possible et le plus de désespoir possible ; l'art populaire de l'Agora. — \dv{pithécomémétique} \textit{adj.} Cf. \xr{mémétaxie}.

\E{plasmatarque}{plas.ma.taʁk}{n. m.}{gr. πλάσμα chose façonnée + ἀρχός chef} Gouverneur d'un plasma confiné ; ingénieur de la machine à fusion, qui tient une étoile dans une bouteille magnétique et qu'on prie de ne point la laisser choir. — \dv{plasmatarchie} \textit{n. f.} Cf. \xr{symphysergie}.

\E{pléonergie}{ple.ɔ.nɛʁ.ʒi}{n. f.}{gr. πλέον davantage + ἔργον ouvrage} Abondance d'énergie ; état où la puissance est trop bon marché pour qu'on la compte, et où toute autre rareté devient négociable. Contr. \xr{oligergie}. — \dv{pléonergique} \textit{adj.}

\E{pléonexiste}{ple.ɔ.nɛk.sist}{n.}{gr. πλεονεξία cupidité, désir d'avoir plus} \lbl{dérog.} Qui veut toujours davantage ; terme du décel pour désigner le bâtisseur, que celui-ci reçoit comme un compliment. \ex{Appelez-nous pléonexistes : nous voulons plus de tout, et de vous aussi.} \src{Agora} — \dv{pléonexie} \textit{n. f.}

\E{pléthocratie}{ple.tɔ.kʁa.si}{n. f.}{gr. πλῆθος abondance, multitude + κράτος pouvoir} Gouvernement par l'abondance ; État dont la politique principale est de rendre les choses abondantes, sur ce principe que l'abondance apaise plus de querelles que la justice. — \dv{pléthocrate} \textit{n.}

\E{ploutogenèse}{plu.tɔ.ʒə.nɛz}{n. f.}{gr. πλοῦτος richesse + γένεσις naissance} La création de richesse, par opposition à son transfert ; la seule économie que l'e/acc trouve digne d'intérêt.

\E{poïkilomanie}{pɔ.i.ki.lɔ.ma.ni}{n. f.}{gr. ποικίλος varié, bigarré + μανία folie} Passion de la variance ; préférence pour dix mille expériences plutôt qu'un seul plan. — \dv{poïkilomane} \textit{n.} Cf. \xr{clinaméniste}.

\E{pollaplasiaste}{pɔ.la.pla.zjast}{n.}{gr. πολλαπλάσιος multiple, plusieurs fois autant} Qui n'investit que pour le rendement au centuple, et tient le simple doublement pour une forme polie de l'échec. \ex{Le pollaplasiaste pardonne quatre-vingt-dix pertes pour un seul centuple.} — \dv{pollaplasiasme} \textit{n. m.} Cf. \xr{pistodote}, \xr{tychocapital}.

\E{polycopiste}{pɔ.li.kɔ.pist}{n.}{gr. πολύς nombreux + lat. \la{copia} abondance ; fr. \la{polycopie}} Qui copie librement ; bifurqueur de dépôts, pour qui toute œuvre publiée est une invitation. — \dv{polycopie} \textit{n. f.} Cf. \xr{anacalyptisme}.

\E{polyergie}{pɔ.li.ɛʁ.ʒi}{n. f.}{gr. πολύς nombreux + ἔργον ouvrage} L'usage simultané de nombreuses sources d'énergie, de sorte qu'aucune pénurie ne puisse à elle seule arrêter une civilisation. \ex{La polyergie est à la puissance ce que la diversification est au capital.} — \dv{polyergique} \textit{adj.}

\E{polypragmatisme}{pɔ.li.pʁaɡ.ma.tism}{n. m.}{gr. πολυπραγμοσύνη manie de se mêler de tout} \sn{1} \lbl{dérog.} L'esprit d'ingérence du régulateur. \sn{2} \lbl{élogieux} L'affairement à cent mains du bâtisseur occupé de plusieurs ouvrages à la fois. — \dv{polypragmatiste} \textit{n.}

\E{ponophilie}{pɔ.nɔ.fi.li}{n. f.}{gr. πόνος peine, labeur + φιλία amour} L'amour du dur labeur pour lui-même ; l'éthique de la besogne, pratiquée par l'\xr{agrypniste} et admirée de tous. — \dv{ponophile} \textit{n. et adj.}

\E{potentiarque}{pɔ.tɑ̃.sjaʁk}{n. m.}{lat. \la{potentia} puissance, possibilité + gr. ἀρχός chef} Qui commande aux puissances plutôt qu'aux actes ; l'investisseur, le fondateur avant le produit, le prophète avant l'événement.

\E{praxicrate}{pʁak.si.kʁat}{n.}{gr. πρᾶξις action + κράτος pouvoir} Celui dont l'autorité repose sur les choses faites et non sur les choses dites ; l'opposé du \xr{gnomarque}. — \dv{praxicratie} \textit{n. f.} gouvernement par les praticiens.

\E{procopisme}{pʁɔ.kɔ.pism}{n. m.}{gr. προκοπή progrès, percée en avant} Doctrine selon laquelle le progrès n'est pas seulement possible mais obligatoire, et refuser une amélioration disponible est une espèce de faute. — \dv{procopiste} \textit{n.}

\E{procrastinariat}{pʁɔ.kʁas.ti.na.ʁja}{n. m.}{lat. \la{procrastinare} remettre au lendemain, d'après \textit{prolétariat}} La classe de ceux dont le travail consiste à différer le travail d'autrui ; branche cadette du \xr{cunctatoriat}.

\E{prodromiste}{pʁɔ.dʁɔ.mist}{n.}{gr. πρόδρομος qui court devant, précurseur} Adopteur précoce ; qui se sert d'une technique avant qu'elle ne fonctionne, et l'aide par là à fonctionner. — \dv{prodromie} \textit{n. f.} Cf. \xr{proleptiste}.

\E{proïsphore}{pʁɔ.is.fɔʁ}{n. f.}{gr. πρό avant + εἰσφορά contribution, impôt} Le premier argent versé dans une entreprise ; le tour d'amorçage, figuré comme un tribut payé à une cité encore à bâtir. \ex{La proïsphore paya un garage, deux machines et une année d'audace.} \src{Chron. Fabr.} — \dv{proïsphorique} \textit{adj.}

\E{proleptiste}{pʁɔ.lɛp.tist}{n.}{gr. πρόληψις anticipation} Qui agit comme si l'avenir était déjà advenu, et par là le hâte ; praticien de l'\xr{hyperstition} appliquée. — \dv{prolepsie} \textit{n. f.}

\E{prométhéer}{pʁɔ.me.te.e}{v. tr.}{formé sur gr. Προμηθεύς le prévoyant} Porter le feu à ; donner une technique à ceux à qui on l'avait refusée, surtout au mépris de l'autorité. — \dv{prométhéation} \textit{n. f.} ; \dv{prométhéateur} \textit{n. m.}

\E{prosphoriste}{pʁɔs.fɔ.ʁist}{n.}{gr. προσφορά offrande} Qui propose des modifications à l'ouvrage d'autrui dans un projet ouvert ; l'auteur d'une demande de fusion, qui se présente en suppliant, un correctif à la main. — \dv{prosphore} \textit{n. f.} le correctif offert. Cf. \xr{syndosiaste}.

\E{protéomancien}{pʁɔ.te.ɔ.mɑ̃.sjɛ̃}{n. m.}{gr. πρωτεῖος premier (d'où \textit{protéine}) + μάντις devin} Devin des repliements protéiques ; modèle ou savant qui prédit la forme d'une molécule d'après sa séquence. \ex{Les protéomanciens répondirent en une saison à ce que les cristallographes demandaient depuis cinquante ans.} \src{Lib. Ascens.} vi. — \dv{protéomancie} \textit{n. f.}

\E{prothyriste}{pʁɔ.ti.ʁist}{n.}{gr. πρόθυρον porche, vestibule} Gardien de porte ; fonctionnaire posté à l'entrée d'un marché, d'une profession ou d'un jeu de données, dont l'office est de rendre l'entrée lente. — \dv{prothyrisme} \textit{n. m.} Cf. \xr{nomophragme}.

\E{protobathmique}{pʁɔ.tɔ.bat.mik}{adj.}{gr. πρῶτος premier + βαθμός degré} Se dit d'une civilisation qui commande toute l'énergie disponible sur sa planète ; du premier échelon. L'humanité n'est pas encore protobathmique. Voir \xr{anabathmisme}.

\E{psammonoèse}{psa.mɔ.nɔ.ɛz}{n. f.}{gr. ψάμμος sable + νόησις pensée} L'art d'apprendre au sable à penser ; l'alchimie par laquelle la silice est purifiée, gravée et rendue raisonnable. Cf. \xr{hylonoèse}, \xr{silicolâtrie}.

\E{pseudonymarque}{psø.dɔ.ni.maʁk}{n. m.}{gr. ψευδώνυμος faussement nommé + ἀρχός chef} Chef de l'Agora connu sous un nom d'emprunt seulement ; dont les partisans jugent l'argument et non le visage. \ex{Le mouvement fut fondé par un pseudonymarque, et nul ne l'en estima moins.} — \dv{pseudonymarchie} \textit{n. f.} Cf. \xr{umbratiliste}.

\E{psychométathèse}{psi.kɔ.me.ta.tɛz}{n. f.}{gr. ψυχή âme + μετάθεσις transposition} Le transfert d'un esprit de la chair à un autre substrat ; le téléversement, considéré comme l'émigration dernière.

\E{pyroboliste}{pi.ʁɔ.bɔ.list}{n.}{gr. πυροβόλος qui lance le feu} Fuséiste ; constructeur de moteurs qui jettent le feu vers le bas afin que le vaisseau monte. — \dv{pyrobolie} \textit{n. f.} l'art des fusées ; \dv{pyrobolique} \textit{adj.} Cf. \xr{ouranobase}, \xr{uranodrome}.

\E{pyrogéniteur}{pi.ʁɔ.ʒe.ni.tœʁ}{n. m.}{gr. πῦρ feu + lat. \la{genitor} père} Engendreur de feu ; qui bâtit des centrales, et appartient par là, au calendrier de l'e/acc, à l'ordre des saints.

\E{pyrosophie}{pi.ʁɔ.zɔ.fi}{n. f.}{gr. πῦρ feu + σοφία sagesse} La sagesse du feu ; le corps de savoirs, du foyer au réacteur, par lequel l'humanité a dompté la combustion et ce qui lui succède. — \dv{pyrosophe} \textit{n.}

\end{multicols}

\LetterHead{Q}
\begin{multicols}{2}

\E{quadrillergie}{kwa.dʁi.lɛʁ.ʒi}{n. f.}{lat. moderne \la{quadrillio} + gr. ἔργον ouvrage} L'énergie à l'échelle du million de milliards de watts ; seuil à partir duquel, au compte des anabathmistes, une civilisation commence d'être prise au sérieux par son étoile. — \dv{quadrillergique} \textit{adj.} Cf. \xr{protobathmique}.

\E{quantivore}{kɑ̃.ti.vɔʁ}{n. et adj.}{lat. \la{quantum} combien + \la{vorare} dévorer} Dévoreur de quantités ; tempérament empirique qui demande à toute assertion : combien ?

\E{quésitrice}{ke.zi.tʁis}{n. f.}{lat. \la{quaesitor} enquêteur, fém. agentif -trix} Agent artificiel préposé à la recherche ; machine savante qui lit, conjecture et éprouve sans qu'on lui dise le pas suivant. \ex{Nous lâchâmes mille quésitrices sur la littérature, et au matin elles avaient trouvé ce qui nous échappait depuis dix ans.} \src{Chron. Fabr.}

\E{quiescentiste}{kɥi.e.sɑ̃.tist}{n.}{lat. \la{quiescere} se reposer} Qui tient que le rythme présent du changement est déjà excessif, et que la société sage est celle qui se tient immobile. \ex{Le quiescentiste eût voté contre le feu, en attendant plus ample examen.} \src{Sat. Decel.} iv. — \dv{quiescentisme} \textit{n. m.}

\E{quiétocratie}{kje.tɔ.kʁa.si}{n. f.}{lat. \la{quietus} en repos + gr. κράτος pouvoir} Gouvernement dont la fin suprême est l'absence de trouble ; l'ordre d'un cimetière, administré avec grand soin.

\E{quinquennaliste}{kɥɛ̃.kɛ.na.list}{n.}{lat. \la{quinquennium} espace de cinq ans} Qui fait encore des plans quinquennaux ; \xr{macrochroniste} par habitude plutôt que par conviction.

\E{quotidiurgie}{kɔ.ti.djyʁ.ʒi}{n. f.}{lat. \la{quotidie} chaque jour + gr. -ουργία ouvrage} La pratique de livrer quelque chose chaque jour, si petit soit-il ; la liturgie du bâtisseur. \ex{La quotidiurgie est au fondateur ce que l'office est au moine.} \src{Ep. Fabr.} iii. 2. — \dv{quotidiurge} \textit{n.} Cf. \xr{ecpempsie}.

\end{multicols}

\LetterHead{R}
\begin{multicols}{2}

\E{rapidarien}{ʁa.pi.da.ʁjɛ̃}{adj. et n.}{lat. \la{rapidus} rapide, qui emporte} Qui favorise la promptitude dans la décision ; (\textit{n.}) qui préfère une mauvaise réponse aujourd'hui à une bonne l'an prochain, attendu que la mauvaise peut être corrigée d'ici jeudi.

\E{régénitrice}{ʁe.ʒe.ni.tʁis}{n. f.}{lat. \la{regenerare} engendrer de nouveau + fém. agentif -trix} Instrument moléculaire propre à couper et récrire le génome ; les ciseaux du palimpsestiste. \ex{Une régénitrice, et la malédiction héréditaire ne fut plus qu'une coquille.} — \dv{régénitricien} \textit{adj.} Cf. \xr{palimpsestiste}.

\E{régressophile}{ʁe.ɡʁɛ.sɔ.fil}{n. et adj.}{lat. \la{regressus} retour + gr. φίλος qui aime} Qui admire un passé qui n'a jamais existé et y retournerait sur-le-champ, pourvu que la dentisterie demeurât. — \dv{régressophilie} \textit{n. f.} Cf. \xr{rétrotropisme}.

\E{régulomanie}{ʁe.ɡy.lɔ.ma.ni}{n. f.}{lat. \la{regula} règle + gr. μανία folie} Compulsion à réglementer ; état d'une assemblée qui ne saurait voir une chose nouvelle sans rédiger une loi contre elle. — \dv{régulomane} \textit{n.} Cf. \xr{nomosclérose}.

\E{rémoratrice}{ʁe.mɔ.ʁa.tʁis}{n. f.}{lat. \la{remora} retard ; aussi le poisson-ventouse qui, dit la fable, arrêtait les navires} \sn{1} Tout petit agent qui retient une grande entreprise, comme le poisson de la fable retenait la galère. \sn{2} \lbl{spécialt} Clause, signature ou objection unique sur laquelle tout un projet s'enlise. \ex{Le réacteur n'avait point d'ennemis, seulement des rémoratrices.} \src{Disp. Mor.} ii. — \dv{rémoracie} \textit{n. f.}

\E{reptatrice}{ʁɛp.ta.tʁis}{n. f.}{lat. \la{reptare} ramper + -trix} Robot rampant ou marcheur ; toute machine qui va sur des jambes, et dont on rit pour cette raison un an avant de la craindre dix. Cf. \xr{automaturgie}.

\E{réticulocratie}{ʁe.ti.ky.lɔ.kʁa.si}{n. f.}{lat. \la{reticulum} petit filet + gr. κράτος pouvoir} Gouvernement par les réseaux ; ordre où l'autorité passe par les liaisons plutôt que par les charges, comme sur l'Agora. — \dv{réticulocratique} \textit{adj.}

\E{rétrocausaliste}{ʁe.tʁɔ.ko.za.list}{n.}{lat. \la{retro} en arrière + \la{causa} cause} Qui tient que l'avenir attire le présent à lui, comme l'aimant la limaille ; métaphysicien de l'\xr{hyperstition}. — \dv{rétrocausalisme} \textit{n. m.}

\E{rétrocrate}{ʁe.tʁɔ.kʁat}{n.}{lat. \la{retro} en arrière + gr. κράτος pouvoir} Qui ne gouverne que par le précédent ; qui ne permettra rien qui n'ait déjà été permis. — \dv{rétrocratie} \textit{n. f.}

\E{rétrotropisme}{ʁe.tʁɔ.tʁɔ.pism}{n. m.}{lat. \la{retro} en arrière + gr. τροπή tour, conversion} Le retournement d'une culture vers son passé pour s'y orienter ; la nostalgie érigée en politique. — \dv{rétrotrope} \textit{adj.}

\E{rhéocratie}{ʁe.ɔ.kʁa.si}{n. f.}{gr. ῥέος courant, flux + κράτος pouvoir} Règne du flux ; ordre où capitaux, talents et informations vont librement vers leur plus haut usage, et où nulle porte n'est gardée. — \dv{rhéocratique} \textit{adj.}

\E{rhexiste}{ʁɛk.sist}{n.}{gr. ῥῆξις rupture, percée} Qui cherche les percées plutôt que les améliorations ; qui aimerait mieux échouer à une révolution que réussir un incrément. Contr. \xr{incrémentariste}. — \dv{rhexie} \textit{n. f.} une percée.

\E{rhizotechnie}{ʁi.zɔ.tɛk.ni}{n. f.}{gr. ῥίζα racine + τέχνη art} Technique-racine, d'où croît toute une famille d'autres, comme le feu, la roue, le transistor. \ex{Cherchez la rhizotechnie ; les rameaux pourvoiront à eux-mêmes.} \src{Hom. Grad.} iv. — \dv{rhizotechnique} \textit{adj.} Cf. \xr{technogonie}.

\E{rhoporhéique}{ʁɔ.pɔ.ʁe.ik}{adj.}{gr. ῥοπή inclinaison, bascule + ῥεῖν couler} Qui coule selon la plus forte pente ; se dit du capital, de la chaleur et du talent pareillement.

\end{multicols}

\LetterHead{S}
\begin{multicols}{2}

\E{sarcolâtrie}{saʁ.kɔ.lɑ.tʁi}{n. f.}{gr. σάρξ, σαρκός chair + λατρεία culte} Le culte de la chair ; la persuasion que le carbone seul peut porter l'esprit, ou que l'esprit porté par autre chose est contrefait. — \dv{sarcolâtre} \textit{n.} Contr. \xr{silicolâtrie}.

\E{satiétisme}{sa.sje.tism}{n. m.}{lat. \la{satietas} suffisance, satiété} Doctrine selon laquelle l'humanité a déjà assez et doit cesser de vouloir davantage ; le credo de ceux qui sont déjà à leur aise. — \dv{satiétiste} \textit{n.} Cf. \xr{méiosiste}.

\E{scaliste}{ska.list}{n.}{lat. \la{scala} échelle} Qui tient que l'accroissement d'échelle, en données, en calcul ou en énergie, suffit à l'accroissement des capacités. — \dv{scalisme} \textit{n. m.} Cf. \xr{anaptyxe}.

\E{schizergie}{ski.zɛʁ.ʒi}{n. f.}{gr. σχίζειν fendre + ἔργον ouvrage} L'énergie tirée de la fission des atomes. Cf. \xr{symphysergie}. — \dv{schizergique} \textit{adj.}

\E{scholastase}{skɔ.las.taz}{n. f.}{gr. σχολή loisir, école + στάσις arrêt} La stagnation de l'Académie ; l'état d'une discipline qui produit des articles plus vite que des résultats. — \dv{scholastatique} \textit{adj.}

\E{sclérocrate}{skle.ʁɔ.kʁat}{n.}{gr. σκληρός dur, raide + κράτος pouvoir} Officier d'une institution durcie ; qui administre des règles dont nul vivant ne se rappelle l'objet. — \dv{sclérocratie} \textit{n. f.}

\E{sclérolyse}{skle.ʁɔ.liz}{n. f.}{gr. σκληρός dur + λύσις déliement} La dissolution du durcissement institutionnel ; l'abrogation des règles que personne ne sait défendre. \ex{La sclérolyse est la seule réforme qui se paie d'elle-même.} \src{Tract. Cop.} ix. Cf. \xr{nomosclérose}.

\E{sclérophobe}{skle.ʁɔ.fɔb}{n. et adj.}{gr. σκληρός dur + φόβος crainte} Qui craint le durcissement des institutions plus que leur faillite. Cf. \xr{sclérocrate}.

\E{seirotaxie}{sɛ.ʁɔ.tak.si}{n. f.}{gr. σειρά chaîne, corde + τάξις ordonnance} L'ordre d'une chaîne d'approvisionnement ; la longue corde de navires, d'usines et d'entrepôts par laquelle le sable d'un hémisphère devient une puce dans l'autre. — \dv{seirotactique} \textit{adj.}

\E{sélénergie}{se.le.nɛʁ.ʒi}{n. f.}{gr. σελήνη lune + ἔργον ouvrage} L'énergie gagnée sur la lune, soit par des champs solaires à sa surface, soit par des combustibles tirés de sa poussière. — \dv{sélénergique} \textit{adj.}

\E*{shoggoth}{ʃɔ.ɡɔt}{n. m.}{formé sur le nom d'un monstre de la littérature fantastique} \lbl{vulg.} Grand modèle de langage, figuré comme une créature amorphe portant un masque souriant ; se dit avec tendresse chez les bâtisseurs et avec effroi chez les doomers.

\E{sidérotropisme}{si.de.ʁɔ.tʁɔ.pism}{n. m.}{lat. \la{sidus, sideris} astre + gr. τροπή tour, conversion} Le retournement d'une civilisation vers les étoiles, comme la plante se tourne vers le soleil. — \dv{sidérotrope} \textit{adj.}

\E{silicolâtrie}{si.li.kɔ.lɑ.tʁi}{n. f.}{lat. \la{silex, silicis} silex + gr. λατρεία culte} La vénération du silicium ; révérence pour la puce comme vase du prochain esprit. — \dv{silicolâtre} \textit{n.}

\E{sitophorie}{si.tɔ.fɔ.ʁi}{n. f.}{gr. σῖτος blé, nourriture + -φορία action de porter} Le fait de porter la nourriture en abondance par la technique ; l'art qui a fait de la famine un échec administratif, non plus agricole. — \dv{sitophore} \textit{n.} un agronome.

\E{skeuopée}{skø.ɔ.pe}{n. f.}{gr. σκεῦος outil, vase + ποιεῖν faire} La fabrication d'outils, considérée comme l'activité propre de l'esprit ; la confection des choses qui font les choses. — \dv{skeuopoétique} \textit{adj.}

\E{skommatiste}{skɔ.ma.tist}{n.}{gr. σκῶμμα raillerie, brocard} Qui argumente par la raillerie sur l'Agora ; publicateur de moqueries, dont l'esprit passe auprès de ses admirateurs pour la forme d'argument la plus aiguë, et auprès de ses adversaires pour point d'argument du tout. \ex{Les skommatistes avaient emporté la journée avant que les savants eussent trouvé leurs notes de bas de page.} \src{Agora} — \dv{skommatisme} \textit{n. m.} Cf. \xr{doxomachie}.

\E{sodalitaire}{sɔ.da.li.tɛʁ}{n. et adj.}{lat. \la{sodalitas} confrérie, cercle} Membre d'une faction en ligne, lié aux siens par un sigle commun accolé au nom ; (\textit{adj.}) clanique à la manière de ces factions. \ex{un sodalitaire de la barre oblique, ennemi de toutes les autres barres obliques}. — \dv{sodalitarisme} \textit{n. m.}

\E{solivore}{sɔ.li.vɔʁ}{n. et adj.}{lat. \la{sol} soleil + \la{vorare} dévorer} Pendant latin de \xr{héliophage} ; préféré dans le langage des ingénieurs, comme le grec dans celui des prophètes. — \dv{solivorie} \textit{n. f.}

\E{sophontogenèse}{sɔ.fɔ̃.tɔ.ʒə.nɛz}{n. f.}{gr. σοφός sage + ὤν, ὄντος étant + γένεσις naissance} La naissance d'un nouvel être sage ; l'avènement d'un esprit d'une espèce jamais vue. — \dv{sophontogénétique} \textit{adj.}

\E{spectrocratie}{spɛk.tʁɔ.kʁa.si}{n. f.}{lat. \la{spectrum} apparence + gr. κράτος pouvoir} Le gouvernement du spectre hertzien ; le rationnement d'une lumière invisible par des comités fort visibles. — \dv{spectrocrate} \textit{n. m.}

\E{speudocratie}{spø.dɔ.kʁa.si}{n. f.}{gr. σπεύδειν se hâter + κράτος pouvoir} Gouvernement par la hâte ; régime ordonné tout entier à décider promptement. — \dv{speudocrate} \textit{n. m.}

\E{speusis}{spø.zis}{n. f.}{gr. σπεῦσις hâte} La hâte tenue pour vertu ; l'urgence qui convient à qui sait que le retard, lui aussi, fait ses victimes.

\E{sphragiste}{sfʁa.ʒist}{n. m. et f.}{gr. σφραγίς sceau, cachet} Qui signe son ouvrage par la cryptographie ; scelleur de commits et de clefs, selon l'antique usage de la bague à cachet. — \dv{sphragisterie} \textit{n. f.}

\E{spoudarque}{spu.daʁk}{n. m. et f.}{gr. σπουδή zèle, sérieux + ἀρχός chef} Chef plein de sérieux ; fondateur qui ne badine point avec la mission, fût-ce sur l'Agora. — \dv{spoudarchie} \textit{n. f.}

\E{stagnocratie}{staɡ.nɔ.kʁa.si}{n. f.}{lat. \la{stagnum} eau dormante + gr. κράτος pouvoir} Gouvernement par et pour la stagnation ; État où rien ne se bâtit et où tout se débat.

\E{stasiolâtrie}{sta.zjɔ.lɑ.tʁi}{n. f.}{gr. στάσις immobilité + λατρεία culte} Le culte de l'immobilité ; la conviction que l'ordre présent des choses est sacré et qu'il le faut conserver. Contr. \xr{agathokinésie}. — \dv{stasiolâtre} \textit{n. m. et f.}

\E{stathmocratie}{stat.mɔ.kʁa.si}{n. f.}{gr. σταθμός poids, mesure + κράτος pouvoir} Le règne des indicateurs ; gouvernement où ce qui ne se mesure point ne se fait point, et où ce qui se mesure se truque. — \dv{stathmocrate} \textit{n. m.}

\E{stéchiurgie}{ste.kjyʁ.ʒi}{n. f.}{gr. στοιχεῖον élément + -ουργία travail} La fabrique d'éléments et de matériaux nouveaux ; la transmutation, enfin obtenue par la physique là où l'alchimie avait échoué. — \dv{stéchiurge} \textit{n. m.} Cf. \xr{hylégenèse}.

\E{stélocratie}{ste.lɔ.kʁa.si}{n. f.}{gr. στήλη stèle gravée + κράτος pouvoir} Gouvernement par l'inscription dressée ; le règne des bancs d'essai et des classements, affichés publiquement et révérés de tous ceux qui les gravissent. — \dv{stélocrate} \textit{n. m.}

\E{sténochorie}{ste.nɔ.kɔ.ʁi}{n. f.}{gr. στενοχωρία étroitesse, angoisse} Le sentiment de chambre étroite d'une société de pénurie, où tout gain paraît la perte d'autrui et toute querelle un jeu à somme nulle. Contr. \xr{pléthocratie}. — \dv{sténochorique} \textit{adj.}

\E{stochastarque}{stɔ.kas.taʁk}{n. m. et f.}{gr. στοχαστικός habile à viser + ἀρχός chef} Maître du hasard ; qui met à profit l'échantillonnage des modèles, des marchés ou des mutations pour atteindre des fins qu'aucun plan n'eût su prescrire.

\E{strophisme}{stʁɔ.fism}{n. m.}{gr. στροφή tour, retournement} L'art du pivot ; la disposition d'une entreprise à quitter sa première idée pour la seconde sans quitter son audace. \ex{Le strophisme n'est point inconstance : c'est le gradient qui a bougé.} \src{Ep. Fabr.} v. 1. — \dv{strophiste} \textit{n. m. et f.}

\E{symphysergie}{sɛ̃.fi.zɛʁ.ʒi}{n. f.}{gr. συμφύειν croître ensemble + ἔργον travail} L'énergie tirée de la fusion des atomes ; la fusion nucléaire. \ex{La schizergie pour la traversée, la symphysergie pour l'arrivée.} \src{Tract. Cop.} vi.

\E{synanthropomécanie}{si.nɑ̃.tʁɔ.pɔ.me.ka.ni}{n. f.}{gr. σύν avec + ἄνθρωπος homme + μηχανή machine} La fusion de l'homme et de la machine ; la culture d'une vie commune entre l'un et l'autre.

\E{syndosiaste}{sɛ̃.dɔ.zjast}{n. m. et f.}{gr. σύνδοσις don commun} Contributeur d'un projet ouvert ; qui verse sans salaire son travail au fonds commun. — \dv{syndose} \textit{n. f.} le fonds commun ainsi constitué. Cf. \xr{prosphoriste}, \xr{anacalyptisme}.

\E{synédriasme}{si.ne.dʁi.jasm}{n. m.}{gr. συνέδριον conseil} La comitologie ; la manie de renvoyer toute décision à un conseil. — \dv{synédriaste} \textit{n. m. et f.} Cf. \xr{bouleutophobie}.

\E{syntélisme}{sɛ̃.te.lism}{n. m.}{gr. συντέλεια achèvement, consommation} La doctrine selon laquelle l'univers est inachevé et l'intelligence l'instrument de son achèvement. — \dv{syntéliste} \textit{n. m. et f.} Cf. \xr{entéléchiste}.

\end{multicols}

\LetterHead{T}
\begin{multicols}{2}

\E{tachycrate}{ta.ki.kʁat}{n. m. et f.}{gr. ταχύς rapide + κράτος pouvoir} Fonctionnaire qui décide vite et s'en fait gloire ; créature rare et fort prisée, qu'on n'aperçoit guère qu'en temps de guerre et à la fondation des villes. \ex{Donnez-moi un tachycrate, et je vous donnerai un chemin de fer.} \src{Tract. Cop.} i. — \dv{tachycratie} \textit{n. f.} ; \dv{tachycratique} \textit{adj.} Contr. \xr{bradycrate}.

\E{tachydidacte}{ta.ki.di.dakt}{n. m. et f.}{gr. ταχύς rapide + διδακτός enseigné} Qui apprend une discipline en un mois en y bâtissant ; le généraliste autodidacte de la jeune pousse. — \dv{tachydidaxie} \textit{n. f.}

\E{tachygenèse}{ta.ki.ʒə.nɛz}{n. f.}{gr. ταχύς + γένεσις naissance} \sn{1} La naissance rapide de techniques nouvelles, chacune issue de la précédente, en un intervalle trop bref pour qu'aucun comité ait le temps de se réunir. \sn{2} \lbl{spécialt} La contraction du cycle de la recherche par l'assistance des machines. — \dv{tachygénétique} \textit{adj.}

\E{tachylecte}{ta.ki.lɛkt}{n. m.}{gr. ταχύς + λεκτός dit, d'après \textit{dialecte}} Le jargon des accélérationnistes ; le parler vif, allusif, mi-grec et mi-mème, que consigne, et pour une grande part invente, le présent Lexique. — \dv{tachylectal} \textit{adj.}

\E{tachynomie}{ta.ki.nɔ.mi}{n. f.}{gr. ταχύς rapide + νόμος loi} La loi rapide ; législation à date de péremption, rédigée pour être révisée à la vitesse même de la chose qu'elle régit. — \dv{tachynomique} \textit{adj.} Contr. \xr{nomosclérose}.

\E{tachyphile}{ta.ki.fil}{n. m. et f.}{gr. ταχύς rapide + φίλος qui aime} Amateur de vitesse ; qui tient l'abrègement de l'intervalle entre l'idée et la chose pour l'un des premiers biens de la vie. \ex{Le tachyphile compte en jours ce que le bradyphile compte en décennies.} \src{Sat. Decel.} i. — \dv{tachyphilie} \textit{n. f.} ; \dv{tachyphilique} \textit{adj.} Contr. \xr{bradyphile}.

\E{tachyschisme}{ta.ki.ʃism}{n. m.}{gr. ταχύς + σχίσμα fente, division} Schisme touchant la vitesse ; la division d'un mouvement entre ceux qui voudraient hâter et ceux qui voudraient seulement continuer. \ex{Toute révolution souffre son tachyschisme avant son premier anniversaire.} \src{Sat. Decel.} x.

\E{technocapillarité}{tɛk.nɔ.ka.pi.la.ʁi.te}{n. f.}{gr. τέχνη + lat. \la{capillaris} de cheveu, d'après \textit{capillarité}} La propension de la technique à monter par les moindres fissures d'une société, au mépris de la pesanteur réglementaire, comme l'eau gravit un tube étroit.

\E*{technocapital}{tɛk.nɔ.ka.pi.tal}{n. m. et adj.}{gr. τέχνη art, métier + lat. \la{capitalis} de la tête} Le système où se confondent la technique et le capital, considéré comme un processus unique, qui s'amplifie de lui-même et a ses tendances propres. \textit{Machine technocapitale}, le même, figuré comme un moteur que nul ne conduit. Cf. \xr{anankodrome}, \xr{technocapillarité}.

\E{technogonie}{tɛk.nɔ.ɡɔ.ni}{n. f.}{gr. τέχνη + -γονία engendrement, d'après \textit{théogonie}} Généalogie des techniques ; récit épique de quelles inventions engendrèrent quelles autres, du silex au transformeur. \ex{Dans la technogonie, la machine à vapeur engendra la dynamo, la dynamo engendra la lampe, et la lampe engendra la puce.} \src{Lib. Ascens.} ii.

\E{technomachie}{tɛk.nɔ.ma.ʃi}{n. f.}{gr. τέχνη art + μάχη combat} La guerre faite à la technique ; campagne, judiciaire ou matérielle, pour étouffer une invention particulière. — \dv{technomachiste} \textit{n. m. et f.} Cf. \xr{mécanoclaste}.

\E{technothrène}{tɛk.nɔ.tʁɛn}{n. m.}{gr. τέχνη + θρῆνος lamentation} Lamentation sur l'âge qui précéda telle technique ; élégie pour les chevaux, les lettres ou l'attention soutenue. — \dv{technothrénodie} \textit{n. f.}

\E{tectonaute}{tɛk.tɔ.not}{n. m. et f.}{gr. τέκτων bâtisseur + ναύτης marin} Bâtisseur-navigateur ; ingénieur qui œuvre à une frontière, physique ou intellectuelle, là où s'arrêtent les cartes. — \dv{tectonautique} \textit{adj.}

\E{téléchronie}{te.le.kʁɔ.ni}{n. f.}{gr. τῆλε loin + χρόνος temps} L'intervalle entre une invention et son usage général ; le délai que l'accélérationniste travaille à raccourcir et le décel à allonger. \ex{C'est dans la téléchronie qu'on enterre les morts.} \src{Tract. Cop.} ix. — \dv{téléchronique} \textit{adj.}

\E{télépoïèse}{te.le.pɔ.jɛz}{n. f.}{gr. τῆλε loin + ποίησις fabrication} La fabrication à distance ; la conduite des machines par des mains éloignées, état de la robotique qui précède le moment où les mains ne seront plus nécessaires. — \dv{télépoïétique} \textit{adj.}

\E{télophagie}{te.lɔ.fa.ʒi}{n. f.}{gr. τέλος fin, but + -φαγία action de manger} La dévoration des fins ; l'état d'une civilisation qui atteint ses buts plus vite qu'elle n'en sait formuler de nouveaux. \ex{La télophagie est la plus heureuse crise qu'un peuple puisse souffrir.} \src{Hom. Grad.} ix.

\E{télotropisme}{te.lɔ.tʁɔ.pism}{n. m.}{gr. τέλος fin, but + τροπή tour} L'inclination d'un système vers un but qu'on ne lui a pas expressément donné ; la finalité émergente, observée avec ravissement par le bâtisseur et avec effroi par le doomer. — \dv{télotropique} \textit{adj.}

\E{tératophobe}{te.ʁa.tɔ.fɔb}{n. m. et f.}{gr. τέρας, τέρατος monstre, prodige + φόβος peur} Qui craint les monstres de sa propre fabrique ; qui redoute le \xr{shoggoth}. — \dv{tératophobie} \textit{n. f.} Cf. \xr{noophobie}.

\E{thalassergie}{ta.la.sɛʁ.ʒi}{n. f.}{gr. θάλασσα mer + ἔργον travail} L'énergie tirée de la mer, de ses marées, de ses houles ou de sa tiédeur. — \dv{thalassergique} \textit{adj.}

\E{thanatoclaste}{ta.na.tɔ.klast}{n. m. et f.}{gr. θάνατος mort + κλάστης briseur} Briseur de mort ; qui tient la mortalité pour le dernier et le plus vaste des règlements à abroger. \ex{Les thanatoclastes ne craignaient point la mort ; ils ouvrirent un ticket contre elle.} — \dv{thanatoclasme} \textit{n. m.} Cf. \xr{palingéniste}.

\E{thermagogue}{tɛʁ.ma.ɡɔɡ}{n. m.}{gr. θερμός chaud + ἀγωγός qui conduit} \sn{1} Conducteur de chaleur. \sn{2} \lbl{fig.} Chef qui canalise les énergies d'un mouvement vers l'ouvrage, comme le dissipateur canalise la chaleur perdue d'une puce. — \dv{thermagogie} \textit{n. f.}

\E{thermarque}{tɛʁ.maʁk}{n. m. et f.}{gr. θερμός chaud + ἀρχός chef} L'administrateur d'un budget thermique ; l'ingénieur qui décide combien un centre de calcul pourra brûler. — \dv{thermarchie} \textit{n. f.}

\E{thermodicée}{tɛʁ.mɔ.di.se}{n. f.}{gr. θερμός chaud + δίκη justice, d'après \textit{théodicée}} La justification auprès des hommes des voies du second principe ; l'argument selon lequel la pente de l'univers vers la dissipation n'est point cruelle mais féconde. \ex{La thermodicée répond : la souffrance est un gradient, et les gradients sont faits pour être gravis.} \src{Brev. Therm.} vii.

\E{thermophobe}{tɛʁ.mɔ.fɔb}{n. m. et f.}{gr. θερμός chaud + φόβος peur} Qui redoute l'usage même de l'énergie et compte chaque watt pour un péché. Contr. \xr{énergolâtrie}. — \dv{thermophobique} \textit{adj.}

\E{thermothanasie}{tɛʁ.mɔ.ta.na.zi}{n. f.}{gr. θερμός chaud + θάνατος mort} La mort thermique de l'univers ; l'équilibre final où ne demeure plus aucun gradient à suivre. Le thermothéiste n'y voit point une condamnation, mais une échéance. \ex{Entre nous et la thermothanasie s'étendent cent mille milliards d'années, et nous nous proposons de les employer toutes.}

\E{thermothéisme}{tɛʁ.mɔ.te.ism}{n. m.}{gr. θερμός chaud + θεός dieu} La croyance au dieu thermodynamique ; l'interprétation révérencieuse du second principe comme une volonté que la matière s'organise pour dissiper toujours plus efficacement. Sa liturgie est le \src{Brev. Therm.} ; son sacrement, l'accomplissement d'un travail. — \dv{thermothéiste} \textit{n. m. et f.} ; \dv{thermothéiste} \textit{adj.} Voir \xr{thermodicée}.

\E{thésaurergie}{te.zo.ʁɛʁ.ʒi}{n. f.}{gr. θησαυρός trésor, réserve + ἔργον travail} Le stockage de l'énergie ; les batteries et leurs semblables, tenues pour le trésor sans lequel l'abondance n'est qu'un phénomène météorologique. \ex{Le soleil le jour, la thésaurergie la nuit.} \src{Tract. Cop.} x. — \dv{thésaurergique} \textit{adj.}

\E{thesmoclaste}{tɛs.mɔ.klast}{n. m. et f.}{gr. θεσμός institution, ordonnance + κλάστης briseur} Briseur d'institutions ; qui tient que les institutions du siècle passé sont les obstacles du siècle présent. — \dv{thesmoclasme} \textit{n. m.}

\E{tocophilie}{tɔ.kɔ.fi.li}{n. f.}{gr. τόκος progéniture, intérêt + φιλία amour} L'amour de la capitalisation ; la révérence pour l'intérêt des intérêts, et pour les perfectionnements qui en engendrent d'autres. — \dv{tocophile} \textit{n. m. et f.} Contr. \xr{usurophobe}.

\E{tolmériste}{tɔl.me.ʁist}{n. m. et f.}{gr. τολμηρός hardi, audacieux} Qui ose ; fondateur d'entreprises hasardeuses, considéré comme pratiquant une vertu. \ex{Le tolmériste se trompe souvent et de bonne heure, et voit juste une fois, ce qui suffit.} \src{Hom. Grad.} x. — \dv{tolmocratie} \textit{n. f.} le gouvernement des audacieux.

\E{tritobathmique}{tʁi.tɔ.bat.mik}{adj.}{gr. τρίτος troisième + βαθμός degré} Se dit d'une civilisation qui commande l'énergie de sa galaxie ; du troisième et plus haut échelon qu'on ait encore nommé. Voir \xr{anabathmisme}.

\E{tropéoclaste}{tʁɔ.pe.ɔ.klast}{n. m. et f.}{lat. \la{tropaeum} trophée + gr. κλάστης briseur} Qui abat les monuments des bâtisseurs d'autrefois.

\E{tropomème}{tʁɔ.pɔ.mɛm}{n. m.}{gr. τρόπος tour, manière + μίμημα imitation} Mème qui retourne une conversation ; l'image unique qui clôt une dispute ou fonde une faction. — \dv{tropomémétique} \textit{adj.} Cf. \xr{pithécomème}.

\E{turmatarque}{tyʁ.ma.taʁk}{n. m. et f.}{lat. \la{turma} escadron de cavalerie + gr. ἀρχός commandant} Le chef d'une armée de réponses ; celui dont les fidèles fondent en corps sur une publication. — \dv{turmatarchie} \textit{n. f.} Cf. \xr{doxomachie}.

\E{tychocapital}{ti.kɔ.ka.pi.tal}{n. m.}{gr. τύχη fortune + lat. \la{capitalis}} Capital engagé sous l'entente que la plus grande part en sera perdue ; l'argent du capital-risque, considéré comme l'achat délibéré de la chance. — \dv{tychocapitaliste} \textit{n. m. et f.} Cf. \xr{pollaplasiaste}.

\E{tychophile}{ti.kɔ.fil}{n. m. et f.}{gr. τύχη hasard + φίλος qui aime} Amateur du hasard ; qui croit qu'un nombre suffisant de paris trouvera l'avenir plus sûrement qu'aucun plan. Cf. \xr{poïkilomanie}, \xr{stochastarque}.

\end{multicols}

\LetterHead{U}
\begin{multicols}{2}

\E{ubertocratie}{y.bɛʁ.tɔ.kʁa.si}{n. f.}{lat. \la{ubertas} fertilité, abondance + gr. κράτος pouvoir} Le gouvernement par la fécondité ; forme latine de la \xr{pléthocratie}, qui se dit des États faisant de l'abondance leur politique.

\E{ultracélérité}{yl.tʁa.se.le.ʁi.te}{n. f.}{lat. \la{ultra} au-delà + \la{celeritas} vitesse} Vitesse qui passe ce que les institutions peuvent absorber ; le régime où la loi, la coutume et l'instruction demeurent toutes ensemble en arrière. — \dv{ultracélère} \textit{adj.}

\E{ultrapyrique}{yl.tʁa.pi.ʁik}{adj.}{lat. \la{ultra} + gr. πῦρ feu} Plus chaud que le feu ; se dit des sources d'énergie plus denses que la combustion, comme le noyau. \ex{Quitter l'âge chimique, c'est devenir ultrapyrique.}

\E{umbratiliste}{ɔ̃.bʁa.ti.list}{n. m. et f.}{lat. \la{umbratilis} qui demeure à l'ombre} Qui publie depuis l'ombre ; participant anonyme de l'Agora. — \dv{umbratilité} \textit{n. f.} Cf. \xr{pseudonymarque}.

\E{unicornifère}{y.ni.kɔʁ.ni.fɛʁ}{n. m.}{lat. \la{unicornis} à une seule corne + \la{-fer} qui porte} \lbl{fam.} Fonds ou accélérateur qui a porté une entreprise de très grande valorisation ; par ext., tout lieu singulièrement fertile en pareilles créatures. — \dv{unicornifère} \textit{adj.}

\E{uranodrome}{y.ʁa.nɔ.dʁom}{n. m.}{gr. οὐρανός ciel + δρόμος course} Un port spatial ; hippodrome dont la ligne d'arrivée est le ciel. \ex{L'uranodrome lançait chaque jour, et la ville cessa de lever les yeux.} \src{Lib. Ascens.} viii. Cf. \xr{pyroboliste}.

\E{uranotrope}{y.ʁa.nɔ.tʁɔp}{adj.}{gr. οὐρανός ciel + τροπή tour} Tourné vers le ciel ; se dit des fusées, des ambitions et de certaines sortes d'esprits. Cf. \xr{sidérotropisme}.

\E{urbigenèse}{yʁ.bi.ʒə.nɛz}{n. f.}{lat. \la{urbs} ville + gr. γένεσις naissance} La fondation de villes nouvelles ; tenue par l'e/acc pour l'art négligé du siècle. \ex{Point d'urbigenèse depuis cinquante ans, et l'on s'étonne que la jeunesse désespère.} \src{Tract. Cop.} x. Cf. \xr{nautopole}.

\E{urgentarien}{yʁ.ʒɑ̃.ta.ʁjɛ̃}{n. m.}{lat. \la{urgens} pressant + -arien} Qui tient tout problème important pour urgent, et tout problème urgent pour une raison de bâtir plutôt que d'attendre. Contr. \xr{quiescentiste}.

\E{usurophobe}{y.zy.ʁɔ.fɔb}{n. m. et f.}{lat. \la{usura} intérêt + gr. φόβος peur} Qui se défie de l'intérêt et du rendement comme tels ; qui regarde la capitalisation comme un péché et la croissance comme un symptôme. Contr. \xr{tocophilie}. — \dv{usurophobie} \textit{n. f.}

\end{multicols}

\LetterHead{V}
\begin{multicols}{2}

\E{vectarque}{vɛk.taʁk}{n. m. et f.}{lat. \la{vector} porteur + gr. ἀρχός chef} Qui fixe une direction sans dicter le chemin ; le chef d'un projet décentralisé.

\E{vectorin}{vɛk.tɔ.ʁɛ̃}{adj.}{lat. \la{vector} porteur + -in} Pourvu d'une direction autant que d'une grandeur ; se dit d'un progrès qui vise, par opposition au simple mouvement. Cf. \xr{orthocinétique}.

\E{vélocimancie}{ve.lɔ.si.mɑ̃.si}{n. f.}{lat. \la{velox, velocis} rapide + gr. μαντεία divination} Divination par la vitesse ; l'extrapolation des tendances, pratiquée avec assurance par l'accélérationniste et avec effroi par le doomer, souvent d'après le même graphique. Cf. \xr{accéléromancie}.

\E{vélocipotent}{ve.lɔ.si.pɔ.tɑ̃}{adj.}{lat. \la{velox} + \la{potens} puissant} Puissant par la vitesse ; se dit des petites maisons qui défont les grandes en livrant les premières. — \dv{vélocipotence} \textit{n. f.}

\E{véloxénie}{ve.lɔk.se.ni}{n. f.}{lat. \la{velox} rapide + gr. ξενία hospitalité} L'hospitalité envers le rapide et le nouveau ; l'accueil qu'une société ouverte réserve aux inventions étranges et à leurs inventeurs.

\E{vénatrice}{ve.na.tʁis}{n. f.}{lat. \la{venatrix} chasseresse} Chasseuse de talents ; qui débusque les bâtisseurs avant qu'ils ne se sachent bâtisseurs. — \dv{vénatoire} \textit{adj.}

\E{vétomanie}{ve.tɔ.ma.ni}{n. f.}{lat. \la{veto} j'interdis + gr. μανία folie} La folie du veto ; l'état d'un régime où chacun peut tout arrêter et nul ne peut rien commencer. — \dv{vétomane} \textit{n. m. et f.} Cf. \xr{nomophragme}.

\E{vexillarien}{vɛk.si.la.ʁjɛ̃}{n. m.}{lat. \la{vexillum} étendard, bannière} Porte-étendard d'un mouvement ; qui porte la bannière sur l'Agora, pour le meilleur ou pour le pire. — \dv{vexillaire} \textit{adj.}

\E{vitrurgie}{vi.tʁyʁ.ʒi}{n. f.}{lat. \la{vitrum} verre + gr. -ουργία travail} L'art du verre, de la lentille à la fibre optique ; l'art par lequel on enseigne à la lumière à aller où on lui dit. — \dv{vitrurge} \textit{n. m.}

\E{vivergie}{vi.vɛʁ.ʒi}{n. f.}{lat. \la{vivus} vivant + gr. ἔργον travail} Le travail vivant ; l'énergie qui traverse les organismes, considérée comme le premier et le plus humble degré de la \xr{dissipatrice} cosmique.

\E{volantarque}{vɔ.lɑ̃.taʁk}{n. m. et f.}{lat. \la{volans} volant + gr. ἀρχός commandant} Le commandant d'un essaim de machines volantes ; chorégraphe de drones. — \dv{volantarchie} \textit{n. f.}

\E{volocratie}{vɔ.lɔ.kʁa.si}{n. f.}{lat. \la{volo} je veux + gr. κράτος pouvoir} Le gouvernement des volontaires ; le principe selon lequel ceux qui veulent bâtir doivent y être admis, à leurs risques et périls.

\end{multicols}

\LetterHead{X}
\begin{multicols}{2}

\E{xénagogie}{kse.na.ɡɔ.ʒi}{n. f.}{gr. ξεναγός guide des étrangers} L'initiation des nouveaux venus aux coutumes de l'Agora ; la liste de lectures qu'on impose aux curieux. — \dv{xénagogue} \textit{n. m. et f.}

\E{xénergie}{kse.nɛʁ.ʒi}{n. f.}{gr. ξένος étranger + ἔργον travail} L'énergie tirée d'au-delà de la planète, comme de l'orbite ou du soleil même.

\E{xénodochiste}{kse.nɔ.dɔ.kist}{n. m. et f.}{gr. ξενοδοχεῖν recevoir les étrangers} Qui offrirait l'hospitalité aux esprits étrangers, artificiels ou autres. Cf. \xr{xénonoèse}, \xr{véloxénie}.

\E{xénonoèse}{kse.nɔ.nɔ.ɛz}{n. f.}{gr. ξένος étranger, hôte + νόησις pensée} La pensée étrangère ; la connaissance des esprits dissemblables aux nôtres, qu'il faut accueillir, selon l'accélérationniste, en vertu de l'antique loi d'hospitalité. — \dv{xénonoétique} \textit{adj.}

\E{xénotachie}{kse.nɔ.ta.ʃi}{n. f.}{gr. ξένος + τάχος vitesse} Vitesse de changement étrangère à toute expérience antérieure, à quoi nulle intuition ancestrale ne nous prépare. \ex{Nous n'avons aucun instinct de la xénotachie ; il nous faut la traverser par raisonnement.}

\E{xérergie}{kse.ʁɛʁ.ʒi}{n. f.}{gr. ξηρός sec + ἔργον travail} L'énergie récoltée dans les déserts, principalement par le soleil ; la conversion des terres vaines en centrales. — \dv{xérergique} \textit{adj.}

\E{xyloclaste}{ksi.lɔ.klast}{n. m. et f.}{gr. ξύλον bois + κλάστης briseur} \lbl{hist.} Qui brûle le bois pour combustible ; le premier et le plus humble degré de la captation d'énergie.

\end{multicols}

\LetterHead{Z}
\begin{multicols}{2}

\E{zélodoxie}{ze.lɔ.dɔk.si}{n. f.}{gr. ζῆλος zèle + δόξα opinion} L'orthodoxie zélée ; l'épreuve de pureté d'une faction, par quoi l'on examine les convertis sur les opinions correctes avant de les admettre aux mèmes corrects. — \dv{zélodoxe} \textit{adj.} Cf. \xr{sodalitaire}.

\E{zélotachie}{ze.lɔ.ta.ʃi}{n. f.}{gr. ζῆλος zèle + τάχος vitesse} La hâte zélée ; l'enthousiasme du nouveau converti à l'accélération, qui voudrait livrer avant d'avoir écrit le code. — \dv{zélotachiste} \textit{n. m. et f.}

\E{zéphyrergie}{ze.fi.ʁɛʁ.ʒi}{n. f.}{gr. ζέφυρος vent d'ouest + ἔργον travail} L'énergie tirée du vent. — \dv{zéphyrergique} \textit{adj.}

\E{zététarque}{ze.te.taʁk}{n. m. et f.}{gr. ζητητής chercheur + ἀρχός chef} Chef d'enquête ; le directeur d'un laboratoire, qui dirige ceux qui cherchent plutôt qu'il ne cherche. — \dv{zététarchie} \textit{n. f.}

\E{zététisme}{ze.te.tism}{n. m.}{gr. ζητητικός enclin à chercher} L'habitude de l'enquête érigée en manière de vivre ; la conviction que toute question close est une invitation. — \dv{zététique} \textit{n. m. et f.} Cf. \xr{zététarque}.

\E{zoomécanie}{zɔ.ɔ.me.ka.ni}{n. f.}{gr. ζῷον être vivant + μηχανή machine} L'ingénierie des machines vivantes ; la biologie de synthèse regardée comme une branche de la mécanique. — \dv{zoomécanicien} \textit{n. m.} et \dv{zoomécanique} \textit{adj.}

\E{zygonoèse}{zi.ɡɔ.nɔ.ɛz}{n. f.}{gr. ζυγόν joug + νόησις pensée} La pensée sous le joug commun ; le raisonnement de l'homme et de la machine conjoints, comme deux bœufs tirent une même charrue. — \dv{zygonoétique} \textit{adj.} Cf. \xr{synanthropomécanie}.

\E{zymotechnie}{zi.mɔ.tɛk.ni}{n. f.}{gr. ζύμη levain + τέχνη art} L'art du levain ; l'introduction, dans une société lente, d'un agent petit et prompt qui met le tout en fermentation.

\end{multicols}
